% Custom XeLaTeX/xepersian template for the Persian manuscript
% Karim Jabbary & Ali Jabbary — Quranic Emotion Spectrum
% Designed for compilation with XeLaTeX so xepersian can manage RTL.

\documentclass[12pt,a4paper]{article}

% Page geometry
\usepackage[a4paper,top=2.6cm,bottom=2.6cm,left=2.6cm,right=2.6cm]{geometry}

% Fonts (XeLaTeX) ------------------------------------------------------
\usepackage{fontspec}
\defaultfontfeatures{Ligatures=TeX}

% Latin (for inline English in citations, technical terms, references)
\setmainfont{EB Garamond}[
  Path        = ../fonts/,
  Extension   = .ttf,
  UprightFont = ebgaramond-regular,
  Numbers     = Lining,
]

% Arabic font for Quranic insertions
\newfontfamily\arabicfont{Amiri}[
  Path        = ../fonts/,
  Extension   = .ttf,
  UprightFont = amiri-regular,
  BoldFont    = amiri-bold,
  Script      = Arabic,
]
\newcommand{\textarabic}[1]{{\arabicfont #1}}
\DeclareRobustCommand{\arb}[1]{{\arabicfont #1}}

% xepersian — must come AFTER fontspec
\usepackage{xepersian}
\settextfont{Vazirmatn}[
  Path        = ../fonts/,
  Extension   = .ttf,
  UprightFont = vazirmatn-regular,
  BoldFont    = vazirmatn-bold,
  Script      = Arabic,
  Scale       = 1.0,
]
\setlatintextfont{EB Garamond}[
  Path        = ../fonts/,
  Extension   = .ttf,
  UprightFont = ebgaramond-regular,
]
% Digits: keep Persian digits in body, Latin digits in math
\setdigitfont{Vazirmatn}[
  Path        = ../fonts/,
  Extension   = .ttf,
  UprightFont = vazirmatn-regular,
  Script      = Arabic,
]

% Section numbering — Persian manuscript uses manual numbering ("۳. روش")
% so suppress LaTeX auto-numbers to prevent double prefixes.
\setcounter{secnumdepth}{-1}

% Section formatting
\usepackage[explicit]{titlesec}
\titleformat{\section}{\Large\bfseries}{\thesection}{0.6em}{#1}
\titleformat{\subsection}{\large\bfseries}{\thesubsection}{0.6em}{#1}
\titleformat{\subsubsection}{\normalsize\bfseries}{\thesubsubsection}{0.6em}{#1}
\titlespacing*{\section}{0pt}{1.6ex plus 0.3ex minus 0.2ex}{1ex plus 0.2ex}
\titlespacing*{\subsection}{0pt}{1.2ex plus 0.3ex minus 0.2ex}{0.8ex plus 0.2ex}

% Spacing
\usepackage{setspace}
\setstretch{1.35}        % Persian needs more line spacing for clarity
\setlength{\parskip}{0.3\baselineskip}
\setlength{\parindent}{1.4em}

% Pandoc compatibility macros
\providecommand{\tightlist}{\setlength{\itemsep}{0pt}\setlength{\parskip}{0pt}}
\usepackage{calc}
\providecommand{\real}[1]{#1}
\usepackage{xcolor}
\providecommand{\pandocbounded}[1]{#1}

% Tables, figures, captions
\usepackage{graphicx}
\graphicspath{{../../figures_fa/}{../../figures/}}
\usepackage[font=small,labelfont=bf,labelsep=period]{caption}
% Persianise the figure / table prefixes
\renewcommand{\figurename}{شکل}
\renewcommand{\tablename}{جدول}
\usepackage{booktabs}
\usepackage{array}
\usepackage{tabularx}
\usepackage{longtable}
\usepackage{makecell}
\renewcommand{\arraystretch}{1.25}

% Lists
\usepackage{enumitem}
\setlist{leftmargin=1.6em,topsep=0.3ex,itemsep=0.3ex}

% Hyperlinks (load near end)
\usepackage[hidelinks]{hyperref}
\urlstyle{same}
\hypersetup{
  pdftitle={پدیدارشناسی طیف هیجانی خشم در قرآن کریم},
  pdfauthor={Karim Jabbary; Ali Jabbary},
  pdfsubject={Quranic Studies; Cognitive Linguistics},
  pdfkeywords={Quran; semantic field; cognitive linguistics; emotion regulation},
  bookmarksdepth=2,
  bookmarksnumbered=true,
  colorlinks=false,
}

% Suppress orphan/widow lines
\widowpenalty=10000
\clubpenalty=10000

\begin{document}

\section{پدیدارشناسی طیفِ هیجانیِ خشم در قرآن کریم}\label{ux67eux62fux6ccux62fux627ux631ux634ux646ux627ux633ux6cc-ux637ux6ccux641-ux647ux6ccux62cux627ux646ux6cc-ux62eux634ux645-ux62fux631-ux642ux631ux622ux646-ux6a9ux631ux6ccux645}

\subsection{تحلیلِ شبکه معناییِ ناخشنودی، برافروختگی و ویرانگری بر اساس متغیرِ شدتِ کنش}\label{ux62aux62dux644ux6ccux644-ux634ux628ux6a9ux647-ux645ux639ux646ux627ux6ccux6cc-ux646ux627ux62eux634ux646ux648ux62fux6cc-ux628ux631ux627ux641ux631ux648ux62eux62aux6afux6cc-ux648-ux648ux6ccux631ux627ux646ux6afux631ux6cc-ux628ux631-ux627ux633ux627ux633-ux645ux62aux63aux6ccux631-ux634ux62fux62a-ux6a9ux646ux634}

\textbf{نویسنده مسئول:} کریم جباری
گروه زبان و ادبیات عرب، دانشکده ادبیات و علوم انسانی، دانشگاه ارومیه، ارومیه، ایران

\textbf{نویسنده دوم:} علی جباری
دانشگاه ارومیه، ارومیه، ایران

\begin{center}\rule{0.5\linewidth}{0.5pt}\end{center}

\subsection{چکیده ساختاریافته}\label{ux686ux6a9ux6ccux62fux647-ux633ux627ux62eux62aux627ux631ux6ccux627ux641ux62aux647}

\textbf{پیشینه و چشم‌انداز.} قرآن کریم در توصیف طیف هیجانی خشم در انسان، از یک دستگاه واژگانی‌-مفهومیِ بسیار دقیق بهره می‌گیرد؛ دستگاهی که در ادبیاتِ پژوهشیِ پیشین، عمدتاً به شیوهٔ تک‌واژه‌نگرانه یا اخلاق‌محور بررسی شده و سرشتِ شبکه‌ایِ آن---یعنی نسبت‌های همنشینی، جانشینی، و درجه‌داریِ شدتِ میان واژگان---تا کنون به‌صورتِ نظام‌مند تبیین نشده است. پژوهش حاضر، با تکیه بر ادغامِ سه چارچوبِ روش‌شناختی---معناشناسی میدانیِ ایزوتسو (۱۹۵۹، ۱۹۶۴، ۱۹۶۶)، معناشناسی طیفیِ کندی-ساسون \lr{(Kennedy 2007; Sassoon 2010)}، و نظریه استعاره مفهومیِ لیکاف-کوچش (Lakoff \& Kövecses 1987؛ Kövecses 2000)---درصدد است نشان دهد که چهارده ریشه کانونیِ عربی در قرآن---\textbf{أف، کره، ضيق، حزن، أسف، نقم، سخط، مقت، غضب، حرد، غيظ، بغي، طغي، عتو}---یک پیوستارِ معناییِ متصل و درجه‌بندی‌شده را شکل می‌دهند که بر روی محورِ «شدت کنش» سامان‌بندی می‌شود.

\textbf{روش‌شناسی.} داده‌های پژوهش از پیکره ریخت‌شناختیِ قرآنی (\lr{Quranic Arabic Corpus}، نسخه ۰٫۴، تنظیمِ کیس دوکس، ۲۰۱۱؛ منتشرشده تحت لیسانس \lr{GPL}) استخراج گردیده است. این پیکره ۱۲۸٬۲۱۹ پاره‌واژه قرآنی را با اطلاعاتِ کاملِ ریشه، شکل لمی، صورتِ سطحی، و برچسب جزء سخن \lr{(POS)} فراهم می‌آورد. خط لوله استخراج، در زبان پایتون و تحت لیسانس MIT، در پنج گام عمل می‌کند: تجزیه ریخت‌شناختی، فیلتر بر اساس ریشه‌های هدف، تجمیعِ پاره‌واژگان، پیوند با فراداده مکی/مدنی، و تولیدِ خروجی‌های جدولی، آماری و شبکه‌ای. تمامیِ یافته‌های گزارش‌شده با یک خط فرمان قابل بازتولید است.

\textbf{یافته‌ها.} تحلیلِ ۳۱۲ بسامد در چهارده ریشه کانونی شواهدی قاطع برای پیوستارِ پیشنهادی فراهم می‌آورد. هندسه طیفی شش‌مرحله‌ای به‌صورتِ زیر مشاهده می‌شود: مرحله ۱ (پیش‌از‌خشمِ ناخشنودی---أف، کره): ۴۴ بسامد؛ مرحله ۲ (فشار درونی و انقباض---ضیق، حزن، أسف): ۶۰ بسامد؛ مرحله ۳ (نفرتِ ارزیابانه---نقم، سخط، مقت): ۲۷ بسامد؛ مرحله ۴ (خشمِ فعال---غضب، حرد): ۲۵ بسامد؛ مرحله ۵ (خشمِ متراکم/انفجاری---غیظ، با تمیّز به‌مثابهِ تجلی): ۱۱ بسامد؛ مرحله ۶ (پیامدهای رفتاری همراه با ملاحظات---بغی، طغیان، عتوّ): ۱۴۵ بسامد. آزمونِ خی-دو برای فرضِ توزیعِ یکنواختِ واژگان در شش مرحله، رد می‌شود: \lr{χ²(5) = 227.15, p < 0.001.} آزمونِ تکمیلیِ مقایسهٔ هسته‌ی پدیدارشناختی (مرحله‌های ۱ تا ۵، ۱۶۷ بسامد) با قطبِ پیامدهای رفتاری (مرحله ۶، ۱۴۵ بسامد) \textbf{رد نمی‌شود} (\lr{χ²(1) = 1.55}). بنا بر اصلِ آماریِ متعارف، عدمِ ردِ فرضیه‌ی صفر، \emph{فقدانِ شواهدِ تفاوت} است و نه شواهدِ مساوات؛ نزدیکیِ مشاهده‌شده، با قرائتی که در §۴-۸ و §۵-۴ بسط می‌یابد \emph{سازگار} است (ریشه‌های مرحله ۶ تنها خاستگاهِ خشمی ندارند و در رابطه‌ای دوسویه با میدانِ خشم قرار دارند)، ولی به‌تنهایی این قرائت را اثبات نمی‌کند. دو الگوی توزیعیِ شایانِ توجهِ پیشین تأیید می‌شوند (سَخَط انحصاراً مدنی، ۰/۴، \lr{p = 0.005} در برابرِ خطِ پایه‌ی مستخرج از پیکره (۰٫۷۴ مکی)؛ أَسَف انحصاراً مکی، ۵/۰)، و الگوهای جدید گزارش می‌گردند: نقم با گرایشِ مکی (۱۲/۵)، مقت (۴/۲)، و حَرد به‌مثابهِ هاپَکس قرآنی (آیه‌ی ۲۵ سوره‌ی قلم، مکی). تحلیلِ شبکه هم‌رخدادگیریِ آیه‌ای نشان می‌دهد که «بَغی» با بالاترین مرکزیتِ نزدیکی \lr{(closeness = 0.55)} و در صدرِ مرکزیتِ بینابینی (\lr{betweenness = 0.15}، هم‌تراز با غضب و أسف) نقشِ یکی از گره‌های پلیِ اصلی میان طیفِ هیجانیِ خشم و میدانِ اخلاقی-اجتماعیِ ظلم/جرم/فسق/عداوت را ایفا می‌کند.

\textbf{نتیجه‌گیری و سهمِ نظری.} قرآن کریم، در بنیادی‌ترین لایه دستگاه واژگانیِ خود، یک «منطقِ تشدیدیِ خشم» را رمزگذاری کرده که از کمترین نشانه‌ی ناخشنودی («أف») آغاز می‌شود و، در صورت عدم مهار و در پیوند با عواملِ مستقلِ ساختاری (طمع، استکبار)، می‌تواند به پیامدهای رفتاری-ساختاری منجر گردد. این منطق با مدلِ فرآیندیِ تنظیم هیجانِ گراس \lr{(Gross 1998, 2014, 2015)}، چرخه شدتیِ پلاتچیک \lr{(Plutchik 1980)}، و سامانه سنجشِ خشمِ اسپیلبرگر \lr{(Spielberger 1999)} از حیث ساختاری هم‌سو است؛ ولی بر خلافِ این مدل‌ها که در محدوده روان‌شناسیِ فردی متوقف می‌شوند، گفتمانِ قرآنی، پیوستار را تا ساحتِ ساختارِ اجتماعی و عصیان وجودی گسترش می‌دهد. سهمِ این پژوهش سه‌وجهی است: (الف) \textbf{روش‌شناختی}، با ارائه‌ی ترکیبی نظام‌مند از سه چارچوبِ ایزوتسو + کندی + لیکاف بر روی متنِ قرآنی --- ترکیبی که، تا جایی که نگارندگان دیده‌اند، در ادبیاتِ مرور-شده‌ی پژوهش‌های قرآنی به این صورت گرد هم نیامده است؛ (ب) \textbf{توصیفی}، با عرضه‌ی یک کنکوردانسِ بازتولیدپذیر برای ۳۱۲ بسامد در چهارده ریشه؛ و (ج) \textbf{نظری}، با صورت‌بندیِ «فرضیه‌ی پیوستارِ شدتیِ قرآنی» \lr{(Qurʾānic Action-Intensity Continuum Hypothesis)} و به‌رسمیت‌شناختنِ علیتِ دوسویه در مرحله‌ی ۶، به‌مثابه چارچوبی برای پژوهش‌های آتی در حوزه‌های قرآن‌کاوی، روان‌شناسیِ اسلامی، ترجمه، و تنظیم هیجان.

\textbf{واژگان کلیدی:} قرآن کریم، طیف هیجانی خشم، شبکه معنایی، زبان‌شناسی شناختی، معناشناسی طیفی، خشم، طغیان، تنظیم هیجان، ایزوتسو، استعاره مفهومی، کنکوردانس قرآن.

\begin{center}\rule{0.5\linewidth}{0.5pt}\end{center}

\subsection{۱. مقدمه}\label{ux645ux642ux62fux645ux647}

\subsubsection{۱-۱. بیان مسئله و شکافِ پژوهشی}\label{ux628ux6ccux627ux646-ux645ux633ux626ux644ux647-ux648-ux634ux6a9ux627ux641-ux67eux698ux648ux647ux634ux6cc}

تأمل در زبانِ قرآن نشان می‌دهد که این متن، در صورت‌بندیِ طیفِ هیجانیِ خشم، از یک دستگاه واژگانیِ غنی و مدرَج بهره می‌گیرد که در آن، هر واژه، نه یک هم‌معنای منفعلِ واژه‌های هم‌میدان، بلکه نقطه مشخصی در یک پیوستارِ شناختی-رفتاری را اشغال می‌کند. واژگانی نظیر «أُفّ» (کمترین نشانه‌ی کلامیِ ناخشنودی، اسراء: ۲۳)، «کره» (کراهتِ درونی، بقره: ۲۱۶)، «ضِيق» (تنگیِ سینه، حجر: ۹۷)، «حُزن» (اندوه، بقره: ۳۸)، «أَسَف» (اندوهِ آمیخته با خشم، زخرف: ۵۵)، «نَقم» (نکوهشِ همراه با اراده‌ی انتقام، اعراف: ۱۲۶)، «سَخَط» (نارضایتیِ شدید، محمد: ۲۸)، «مَقت» (نفرتِ اخلاقیِ شدید، غافر: ۱۰)، «غَضَب» (خشمِ هدف‌مند، فتح: ۶)، «حَرد» (خشمِ همراه با قصدِ منع، قلم: ۲۵)، «غَیظ» (خشمِ متراکم، آل عمران: ۱۳۴)، «تَمَیُّز مِن الغَیظ» (فروپاشیِ هیجانی، ملک: ۸)، «بَغی» (تجاوز، شوری: ۴۲)، «طُغیان» (سرکشی، علق: ۶) و «عُتُوّ» (سرکشیِ لجوجانه، فرقان: ۲۱)، در پراکندگیِ آیاتیِ خود تصویرگر یک هندسه معناییِ منسجم‌اند که در آن، نقطه آغاز، نشانه‌ای کلامی یا درونی از ناخشنودی است و نقطه پایان، یک رفتارِ اجتماعی-ساختاریِ تثبیت‌شده در ضدیت با نظمِ اخلاقیِ متن---با این هشدار که، چنان‌که در §۴-۶ و §۵-۴ بسط می‌یابد، رابطه‌ی نقطه‌ی پایان با خشم، \textbf{دوسویه} است و نه صرفاً پیامدِ یک‌جانبه.

با این‌همه، پژوهش‌های پیشین، چه در سنت تفسیریِ کلاسیک و چه در دو دهه اخیر در پژوهش‌های زبان‌شناختیِ قرآن، این دستگاه را عمدتاً به‌مثابه مجموعه‌ای از واحدهای منفصل تحلیل کرده‌اند. این رویکردِ ذره‌گرایانه، اگرچه در دقتِ لغوی و توضیحِ خاستگاهِ هر واژه ارزشمند بوده، یک مسئله بنیادی را پیشِ روی پژوهشگر قرار می‌دهد: \textbf{اگر این واژگان، اعضای یک شبکه واحدند، نسبتِ ساختاریِ آنها چیست؟} آیا «غضب» شدیدتر از «سخط» است یا متفاوت‌بافت با آن؟ آیا «بَغی» در امتداد «غضب» قرار می‌گیرد یا در میدانی متفاوت؟ آیا توالیِ «ضیق → غضب → غیظ → بَغی» از حیثِ پیکره‌ای قابل پشتیبانی است یا صرفاً یک شِمای حدسی؟ پاسخ به این پرسش‌ها مستلزم آن است که از چارچوبِ تک‌واژه‌نگری به چارچوبِ شبکه‌نگری و طیف‌نگری گذار شود؛ گذاری که هم نیازمندِ مبانیِ نظریِ نوین (در تلاقیِ زبان‌شناسیِ شناختی، معناشناسیِ صوری، و معناشناسیِ متنیِ ایزوتسو) است و هم نیازمندِ پشتوانه‌ی پیکره‌ایِ قابلِ راستی‌آزمایی.

پژوهش حاضر، در سه مسیرِ همسو، به این شکاف پاسخ می‌دهد. \textbf{نخست}، با تلفیقِ سه نظریه‌ی به‌ظاهر مستقل---معناشناسیِ میدانیِ ایزوتسو که شبکه‌ها را در یک پیکره‌ی متنیِ کلان نظام می‌بخشد، معناشناسیِ طیفیِ کندی-ساسون که درجه‌بندیِ شدت را در پیشایی‌های گرادپذیر صورت‌بندی می‌کند، و نظریه استعاره مفهومیِ لیکاف-کوچش که سازوکارهای شناختیِ هیجان را آشکار می‌سازد---یک چارچوبِ نظریِ ترکیبی پیشنهاد می‌دهد که، تا جایی که نگارندگان در ادبیاتِ مرور-شده دیده‌اند، در پژوهش‌های قرآنیِ پیشین به این صورتِ تلفیقی به کار گرفته نشده است. \textbf{دوم}، این چارچوب را در یک خط لوله‌ی محاسباتیِ کاملاً بازتولیدپذیر، بر روی پیکره‌ی ریخت‌شناختیِ قرآن، عملیاتی می‌سازد و کنکوردانسی برای ۳۱۲ بسامد در چهارده ریشه‌ی کانونیِ طیف هیجانی خشم (به همراهِ پنج ریشه‌ی مفهومیِ مجاور) ارائه می‌دهد. \textbf{سوم}، یافته‌های توصیفیِ این تحلیل را با مدل‌های مدرنِ روان‌شناسیِ هیجان \lr{(Gross, Plutchik, Spielberger, Lazarus, Scherer)} در یک گفت‌وگوی نظریِ نظام‌مند قرار می‌دهد و پیامدهای آن را برای ترجمه قرآن و روان‌شناسیِ اسلامیِ کاربردی صورت‌بندی می‌کند.

\subsubsection{۱-۲. پیشینه پژوهش}\label{ux67eux6ccux634ux6ccux646ux647-ux67eux698ux648ux647ux634}

پژوهش‌های پیشین در حوزه‌ی طیف هیجانی خشم در قرآن را می‌توان در پنج خوشه نسبتاً مستقل دسته‌بندی کرد:

\textbf{الف) پژوهش‌های لغویِ تک‌واژه‌محور.} این خوشه، که گسترده‌ترین و قدیمی‌ترین خوشه‌ی پژوهشی است، به ریشه‌شناسی و کاربردِ قرآنیِ یک واژه می‌پردازد. در زبان فارسی، نَریمانی، اقبالی و چاری (۱۴۰۰) در «بازشناسی معنایی واژگان «غیظ»، «غضب» و «سَخَط» در قرآن کریم با رویکرد تفسیری»، نزدیک‌ترین کارِ پیشین به پژوهش حاضر را عرضه کرده‌اند. ایشان با رویکردِ لغوی-تفسیری نشان داده‌اند که «غضب» کاربردِ کانونی‌اش در توصیفِ خشمِ خداوند است و گاهی به انبیا و مؤمنان نیز نسبت داده می‌شود؛ «سخط» به اوجِ نارضایتیِ الهی، به‌ویژه نسبت به منافقان، اطلاق می‌گردد؛ و «غیظ» در موارد قرآنی، ویژگیِ منکران و منافقان است در حالی که ساحتِ الهی از آن منزه شناخته شده. نتایجِ ایشان از حیثِ تفسیری استوار است، ولی پژوهشِ مذکور، سه واژه را در یک بافتِ مقایسه‌ایِ تک‌مرحله‌ای تحلیل کرده و مرحله‌های پیشینِ هیجان (تنگی-اندوه) و پسینِ آن (طغیان-عتوّ) را پوشش نمی‌دهد. به‌علاوه، چارچوبِ نظریِ ایشان، فاقدِ ابزار صوریِ معناشناسیِ طیفی است؛ بنابراین نمی‌تواند مدلی صریح برای \textbf{درجه‌بندیِ شدت} میان این سه واژه ارائه دهد.

به‌همین قیاس، قاری‌زاده، سلمانی مارَوست، میمندی و فارسی (۱۴۰۲) در «چندمعنایی واژۀ «طغیان» در قرآن کریم با رویکردی بر زبان‌شناختی»، تنها واژه «طغیان» را با ابزار چندمعنایی و روابطِ هم‌نشینی واکاوی کرده‌اند و نشان داده‌اند که این واژه در ۲۷ سوره و ۱۳ صورتِ صرفی به‌کار رفته و معنای کانونیِ آن «تجاوز از حدّ» است. مقالاتی مشابه نیز برای حُزن و بَغی منتشر شده‌اند (مثلاً پایان‌نامه‌های دانشگاهی در دو دهه اخیر در دانشگاه‌های قم و فردوسی مشهد). همه این آثار، به‌رغمِ ارزشِ موضعی، نقصِ بنیادینِ تک‌واژه‌نگری را به‌اشتراک دارند: هیچ‌یک نمی‌تواند نسبتِ ساختاریِ یک واژه را با کلِ شبکه روشن سازد.

\textbf{ب) پژوهش‌های اخلاق‌محور (موعظه‌گرا).} خوشه دوم به موضوع «کظم غیظ» (آل عمران: ۱۳۴) و راهکارهای روایی-قرآنی برای مهار خشم می‌پردازد. این پژوهش‌ها از منظر تربیت اخلاقی ارزشمندند ولی متمرکز بر یک نقطه از پیوستار (ساحتِ غیظ-کظم) باقی می‌مانند و سامانه واژگانیِ گسترده‌تر را مغفول می‌گذارند.

\textbf{ج) پژوهش‌های معناشناسیِ شناختیِ ایزوتسویی.} خوشه سوم، که پس از ترجمه آثار ایزوتسو به فارسی توسط احمد آرام (۱۳۶۱) و فریدون بدره‌ای (۱۳۶۰) رو به گسترش بود، تلاش دارد روشِ میدانی-شبکه‌ایِ ایزوتسو را به مفاهیمِ مختلف قرآنی بسط دهد. اثر بنیادینِ قائمی‌نیا، \textbf{معناشناسی شناختی قرآن} (۱۳۹۰) و کتابِ مکمل \textbf{استعاره‌های مفهومی و فضاهای قرآنی} (۱۳۹۳)، ادغامِ ایزوتسو و استعاره مفهومیِ لیکاف-جانسون را در پژوهش‌های قرآنیِ ایرانی نهادینه کرد. آثار پاکتچی نیز روش‌شناسیِ دقیق‌تری برای تحلیل میدانیِ واژگانِ قرآن فراهم آورد. این آثار، چارچوبِ نظریِ موجود را متحول کرد ولی هیچ‌یک، به‌طورِ خاص، میدانِ طیفِ هیجانیِ خشم را به‌مثابهِ یک پیوستارِ متصل تحلیل نکرده است.

\textbf{د) پژوهش‌های ترجمه‌ای.} خوشه چهارم به مسئله ترجمه واژگانِ هم‌میدان توجه می‌کند و گرایش به یکسان‌سازیِ ترجمه (مثلاً برگردانیِ همه‌ی غضب، غیظ، سخط به «خشم») را به‌مثابه یک خطای ساختاری شناسایی می‌کند. پژوهش‌های اخیرِ محاسباتی این نگرانی را تجربی‌تر می‌سازند: \lr{Almuhanna et al.\ (2025)} نشان می‌دهند که ترجمه‌های انگلیسیِ معتبرِ قرآن در آیاتِ هیجانی نشتیِ معناییِ نظام‌مند دارند، و مجموعه‌داده‌ی چندزبانه‌ی \lr{ELQV} از آیاتِ هیجان‌برچسب‌خورده (۲٬۱۰۰ آیه؛ \lr{Authors 2025} در \emph{\lr{Journal of King Saud University--CIS}}) دقیقاً برای ممیّزی‌پذیرکردنِ این یکسان‌سازی طراحی شده. این پژوهش‌ها، اگرچه به نقصِ سیستمیِ ترجمه‌ها اشاره می‌کنند، فاقدِ یک مدلِ مرجعِ شدتیِ پشتوانه‌ای‌اند که در برابرِ آن بتوان ترجمه‌های موجود را راستی‌آزمایی کرد.

\textbf{ه) پژوهش‌های روان‌شناسیِ اسلامی.} خوشه پنجم، که در دو دهه اخیر گسترش یافته، به استخراجِ مدل‌های تنظیم هیجان از منابع اسلامی و تطبیقِ آن‌ها با مدل‌های روان‌شناسیِ غربی می‌پردازد. از جمله آثار اخیرِ مرورشده‌ی این خوشه می‌توان به تحلیل‌های موضوعیِ ترس، خشم، اندوه و شرم در پیکره‌ی قرآن و حدیث اشاره کرد (نظیر مقاله‌ی منتشرشده در \emph{\lr{International Journal of Counseling and Psychotherapy}}، ۲۰۲۳، شماره‌ی ۳۱۱)، و نیز پروژه‌ی گسترده‌ترِ تطبیقِ نظریه‌ی تنظیم هیجانِ گراس با انسان‌شناسیِ اخلاقیِ کلاسیکِ اسلامی. این آثار، با وجود کاربردگرایی، عمدتاً بر داده‌های روایی تکیه دارند و تحلیلِ زبان‌شناختیِ نظام‌مند بر روی متنِ قرآنی را در درجه‌ی دوم اهمیت قرار می‌دهند --- همان شکافی که پژوهش حاضر، از سمتِ پیکره‌ی واژگانی، در پی پُرکردنِ آن است.

\textbf{شکافِ پژوهشی.} پنج خوشه فوق، در مجموع، همه‌ی نقاطی را که برای صورت‌بندیِ یک تحلیلِ پیوستارِ شدتیِ قرآنی لازم است، پوشش می‌دهند، ولی هیچ‌یک این نقاط را در یک چارچوبِ نظریِ منسجم گرد هم نیاورده است. مشخصاً، تا کنون پژوهشی منتشر نشده که به‌صورتِ هم‌زمان: (۱) همه‌ی واژگانِ کانونیِ طیفِ هیجانیِ خشم را در یک چارچوب واحد بررسی کند؛ (۲) نسبتِ ساختاریِ آن‌ها را با ابزارِ معناشناسیِ طیفی صورت‌بندی نماید؛ (۳) یافته‌ها را با تحلیلِ پیکره‌ایِ بازتولیدپذیر راستی‌آزمایی کند؛ (۴) آن را با مدل‌های مدرنِ روان‌شناسیِ هیجان در گفت‌وگوی نظری قرار دهد؛ و (۵) پیامدهای آن را برای ترجمه و عمل بالینی صورت‌بندی نماید. پژوهش حاضر، هر پنج بُعد را در یک ساختارِ یکپارچه پیش می‌برد.

\subsubsection{۱-۳. سؤال‌ها، اهداف و فرضیه‌های پژوهش}\label{ux633ux624ux627ux644ux647ux627-ux627ux647ux62fux627ux641-ux648-ux641ux631ux636ux6ccux647ux647ux627ux6cc-ux67eux698ux648ux647ux634}

\textbf{سؤال‌های پژوهشی} به‌صورتِ زیر تدوین شده‌اند:

\begin{itemize}
\tightlist
\item
  \textbf{س۱.} میدانِ معناییِ طیفِ هیجانیِ خشم در قرآن از چه واژگانی تشکیل شده و این واژگان در چه ساختارِ شبکه‌ای با یکدیگر مرتبط‌اند؟
\item
  \textbf{س۲.} چگونه می‌توان این واژگان را بر اساسِ متغیرِ «شدت کنش» در یک پیوستارِ منسجم رده‌بندی کرد؟
\item
  \textbf{س۳.} آیا الگوی توزیعِ پیکره‌ایِ این واژگان (فراوانی، توزیع مکی-مدنی، هم‌رخدادگیریِ آیه‌ای، مرکزیتِ شبکه‌ای) رده‌بندیِ نظریِ پیشنهادی را تأیید می‌کند؟
\item
  \textbf{س۴.} پیوستارِ پیشنهادی با کدام مدل‌های مدرنِ روان‌شناسیِ هیجان قابلِ تطبیق است و چه تمایزهای ساختاری‌ای با آن‌ها دارد؟
\item
  \textbf{س۵.} پیامدهای این یافته‌ها برای ترجمه قرآن و روان‌شناسیِ اسلامیِ کاربردی چیست؟
\end{itemize}

\textbf{فرضیه اصلی} که در ادامه «فرضیه پیوستارِ شدتیِ قرآنی» \lr{(Qurʾānic Action-Intensity Continuum Hypothesis)} نامیده می‌شود، چنین است:

\begin{quote}
چهارده ریشه‌ی کانونیِ أف، کره، ضيق، حزن، أسف، نقم، سخط، مقت، غضب، حرد، غیظ، بغي، طغي و عتو در قرآن کریم، نه به‌مثابه واژگانِ مترادف یا نقاطِ منفصل، بلکه به‌مثابه چهارده گرهِ متمایز در یک شبکه معناییِ پیوستار-شکل عمل می‌کنند؛ شبکه‌ای که توسط متغیرِ شدتِ کنش، در شش مرحله ساختار می‌یابد و توالیِ آن، در هر دو ساحتِ کیفی (ریشه‌شناسی، تفسیر، استعاره مفهومی) و کمی (فراوانی پیکره‌ای، توزیع زمانی، تحلیل شبکه)، با شواهدِ قرآنی همخوان است.
\end{quote}

\textbf{فرضیه‌های فرعی} که در بخش یافته‌ها آزموده می‌شوند:

\begin{itemize}
\tightlist
\item
  \textbf{ف.۱.} توزیعِ بسامدیِ چهارده ریشه در شش مرحله از فرضِ یکنواختی به‌طورِ معناداری منحرف است (آزمونِ خی-دو، \lr{χ²(5) = 227.15}).
\item
  \textbf{ف.۲.} هسته‌ی پدیدارشناختی (مرحله‌های ۱ تا ۵) و قطبِ پیامدهای رفتاری (مرحله ۶) از حیثِ شمارشِ خام تفاوتِ معناداری ندارند (\lr{χ²(1) = 1.55}، رد نمی‌شود)؛ این یافته از قرائتِ علیتِ دوسویه‌ی مرحله ۶ (که در §۴-۸ و §۵-۴ بسط می‌یابد) پشتیبانی می‌کند.
\item
  \textbf{ف.۳.} واژه «سَخَط» الگوی توزیعی متمایز و تخصصی‌شده در سور مدنی دارد.
\item
  \textbf{ف.۴.} در شبکه هم‌رخدادگیریِ آیه‌ای، واژه «بَغی» نقشِ گرهِ پلی میان میدانِ هیجانی و میدانِ اخلاقیِ ظلم/جرم/فسق ایفا می‌کند.
\end{itemize}

\subsubsection{۱-۴. ضرورت و کاربردهای پژوهش}\label{ux636ux631ux648ux631ux62a-ux648-ux6a9ux627ux631ux628ux631ux62fux647ux627ux6cc-ux67eux698ux648ux647ux634}

ضرورتِ این پژوهش از سه افق قابلِ تبیین است:

از منظرِ \textbf{علوم قرآنی}، آشکارسازیِ هندسه‌ی دقیقِ شبکه‌ی معناییِ طیفِ هیجانیِ خشم، نشان می‌دهد که زبانِ قرآنی، بدون استفاده از اصطلاحاتِ انتزاعیِ روان‌شناسیِ نوین، یک سامانه‌ی چهارده‌نقطه‌ای از مراحلِ شدتِ هیجان را به‌گونه‌ای هم‌بسته و قابلِ راستی‌آزمایی صورت‌بندی کرده است.

از منظرِ \textbf{روان‌شناسی و علومِ تربیتی}، پیوستارِ پیشنهادی می‌تواند پایه نظریِ مدلِ بومیِ تنظیم هیجان قرار گیرد؛ مدلی که در آن، مفاهیمِ قرآنی نظیر «کظم غیظ» نه به‌مثابه توصیه‌های اخلاقیِ مجزا، بلکه به‌مثابه نقاطِ مداخله در یک سامانه‌ی پیشرفته‌ی تنظیم هیجان عمل می‌کنند.

از منظرِ \textbf{مطالعاتِ ترجمه قرآن}، پژوهش حاضر معیاری دقیق برای ارزیابیِ کیفیت ترجمه‌ها در میدانِ طیف هیجانی خشم فراهم می‌آورد. ترجمه‌ای که تمایز شدتیِ غضب، غیظ، سخط، مقت و تمیّز را از میان ببرد، آسیبی جدی به سامانه مفهومیِ متن وارد می‌کند؛ پیشنهادهای روشن این پژوهش (در بخش ۵-۵) می‌تواند راهنمای مترجمانِ آینده باشد.

\begin{center}\rule{0.5\linewidth}{0.5pt}\end{center}

\subsection{۲. مبانی نظری}\label{ux645ux628ux627ux646ux6cc-ux646ux638ux631ux6cc}

\subsubsection{۲-۱. معناشناسی میدانیِ ایزوتسو: شبکه به‌مثابه واحدِ تحلیل}\label{ux645ux639ux646ux627ux634ux646ux627ux633ux6cc-ux645ux6ccux62fux627ux646ux6cc-ux627ux6ccux632ux648ux62aux633ux648-ux634ux628ux6a9ux647-ux628ux647ux645ux62bux627ux628ux647-ux648ux627ux62dux62f-ux62aux62dux644ux6ccux644}

روشِ معناشناسی میدانی \lr{(Semantic Field Theory)} که توسط توشیهیکو ایزوتسو در سه اثرِ کلیدیِ او---\textbf{ساختار اصطلاحاتِ اخلاقیِ قرآن} (۱۹۵۹)، \textbf{خدا و انسان در قرآن} (۱۹۶۴)، و \textbf{مفاهیم اخلاقی-دینی در قرآن مجید} (۱۹۶۶)---پایه‌ریزی شد، انقلابی روش‌شناختی در پژوهش‌های قرآنیِ معاصر به ارمغان آورد. ایزوتسو، با الهام از زبان‌شناسیِ ساختاریِ تریر-وایسگربر و در پیوند با فلسفه زبانِ مشرق‌زمین، استدلال کرد که واژگانِ کلیدیِ قرآن---نظیر اللّه، الإیمان، الکُفر، الظُّلم، التَّقوی---نه گزاره‌هایی منفصل بلکه گره‌هایی در شبکه‌ای از روابطِ همنشینی، جانشینی و تقابل‌اند، شبکه‌ای که در مجموع، یک «جهان‌بینیِ زبانیِ قرآنی» \lr{(Qurʾānic Weltanschauung)} را شکل می‌دهد.

سه مفهومِ کانونیِ روش ایزوتسو عبارت‌اند از: \textbf{واژه‌کانون} \lr{(focus word)} که در یک میدانِ معنایی جایگاهِ محوری دارد و تعریف‌گرِ بنیادِ معناییِ آن میدان است؛ \textbf{میدان معنایی} \lr{(semantic field)} که شبکه‌ای از واژگانِ مرتبط است و حول واژه‌کانون شکل می‌گیرد؛ و \textbf{واژه‌کلید} \lr{(key word)} که در گفتمانِ کلیِ متن، اهمیتی فراگیر دارد و چندین میدان را به یکدیگر می‌پیوندد. این سه‌گانه، چنان‌که قائمی‌نیا (۱۳۹۰: ۸۲-۹۸) به‌خوبی نشان داده، ابزارِ نظریِ مناسبی برای تحلیلِ مفاهیمِ پیچیده در متونِ مقدس فراهم می‌آورد.

با این وجود، روشِ ایزوتسو، در صورتِ اولیه آن، یک محدودیتِ قابلِ توجه دارد: این روش به ما می‌گوید که واژگانِ یک میدان «در ارتباط‌اند»، ولی صورتِ ساختاریِ این ارتباط را---به‌ویژه آن جا که ارتباطِ مذکور درجه‌داری \lr{(gradability)} یا شدت داشته باشد---به‌نحوی صریح صورت‌بندی نمی‌کند. ایزوتسو از هندسه فرکانسی، آماری، یا طیفیِ این روابط سخن نمی‌گوید. در واقع، ایزوتسو محدوده ارتباطاتِ مفهومی را تعیین می‌کند ولی \textbf{ترتیب} و \textbf{شدتِ نسبی} را به‌مثابه پدیده‌هایی مستقل بررسی نمی‌کند. این نقطه دقیقاً همان جایی است که نیاز به یک ابزارِ نظریِ مکمل احساس می‌شود.

\subsubsection{۲-۲. معناشناسی طیفی \lr{(Scalar Semantics)}: از شبکه به پیوستار}\label{ux645ux639ux646ux627ux634ux646ux627ux633ux6cc-ux637ux6ccux641ux6cc-scalar-semantics-ux627ux632-ux634ux628ux6a9ux647-ux628ux647-ux67eux6ccux648ux633ux62aux627ux631}

پاسخ به این محدودیت در زبان‌شناسیِ صوری، در نظریه «معناشناسی طیفی» یافت می‌شود. این نظریه، با ریشه‌هایی در اثر کِنِدی \lr{(Kennedy 1999, 2007)} و توسعه‌های بعدیِ ساسون \lr{(Sassoon 2010)} و موری \lr{(Morey 2014)}، چارچوبی صوری برای تحلیلِ پیشایی‌های گرادپذیر \lr{(gradable predicates)} و قیدهای درجه‌گذار \lr{(degree modifiers)} فراهم می‌کند. در این چارچوب، یک پیشای گرادپذیر، بر روی یک «طیف» \lr{(scale)} قرار می‌گیرد که سه مؤلفه دارد: (الف) یک \textbf{بُعدِ} \lr{(dimension)} که به یک ویژگیِ قابلِ سنجش اشاره می‌کند (مثلاً دما، طول، شدت)، (ب) یک \textbf{ترتیبِ} \lr{(ordering)} که نقاطِ این بُعد را به‌گونه‌ای نظام‌مند مرتب می‌سازد، و (ج) یک \textbf{واحدِ سنجشِ} \lr{(unit of measure)} صریح یا ضمنی.

پژوهش حاضر استدلال می‌کند که میدانِ معناییِ طیفِ هیجانیِ خشم در قرآن، یک طیفِ مدرَج به این معنا است که:

\begin{itemize}
\tightlist
\item
  \textbf{بُعد}: شدتِ کنش \lr{(intensity of action)}، که خود ترکیبی از سه زیربعد است: \lr{(i)} محل تجسدِ هیجان، از درون‌بدنی به بیرون‌بدنی؛ \lr{(ii)} شدتِ تجربیِ ذهنی، از پایین به بالا؛ \lr{(iii)} دامنه پیامدِ اجتماعی، از فرد به ساختار.
\item
  \textbf{ترتیب}: توالیِ شش مرحله‌ای از کمترین نشانه‌ی ناخشنودی (أف، کره) به فشار درونی، نفرتِ ارزیابانه، خشمِ فعال، خشمِ متراکم، و سرانجام پیامدهای رفتاری-ساختاری.
\item
  \textbf{واحد ضمنی}: درجه گذار از تجربه ذاتی به کنشِ بیرونی.
\end{itemize}

این انتقال از شبکه ایزوتسویی به پیوستارِ طیفی، یک پیشرفتِ نظریِ کلیدی است که در پژوهش حاضر به‌صورت روش‌شناختی صریح می‌شود. به دیگر بیان، ما ادعا نمی‌کنیم که قرآن یک «نظریه» معناشناسی طیفی دارد؛ بلکه ادعا می‌کنیم که اگر این چارچوب نظریِ مدرن را بر داده‌های قرآنی اعمال کنیم، یک ساختارِ پیوستار-شکلِ پنهان آشکار می‌شود که با چشم‌اندازِ تک‌واژه‌نگرِ سنتی از دیدِ پژوهشگر مخفی می‌ماند.

\subsubsection{۲-۳. نظریه استعاره مفهومی \lr{(Conceptual Metaphor Theory)}}\label{ux646ux638ux631ux6ccux647-ux627ux633ux62aux639ux627ux631ux647-ux645ux641ux647ux648ux645ux6cc-conceptual-metaphor-theory}

سومین لایه نظریِ این پژوهش، نظریه استعاره مفهومیِ لیکاف-جانسون (Lakoff \& Johnson 1980) و توسعه‌های بعدیِ آن توسط کوچش \lr{(Kövecses 1986, 2000, 2010)} و مَعلِج \lr{(Maalej 2004)} در زمینه عربی است. ایده مرکزیِ این نظریه آن است که استعاره‌ها صرفاً آرایه‌های ادبی نیستند، بلکه ابزارهای شناختیِ بنیادینی‌اند که به ما اجازه می‌دهند مفاهیمِ انتزاعی (نظیر هیجان) را از طریق نگاشت بر روی مفاهیمِ ملموس‌تر (نظیر بدن، حرکت، ظرف، حرارت) درک کنیم.

لیکاف و کوچش (۱۹۸۷) در پژوهشِ کلاسیکِ خود نشان دادند که زبان انگلیسی برای بازنماییِ خشم از یک مجموعه‌ی منسجم از استعاره‌های مفهومی استفاده می‌کند که محورِ آن «خشم گرمای مایعی در یک ظرف فشرده است» \lr{(ANGER IS THE HEAT OF A FLUID IN A PRESSURIZED CONTAINER)} می‌باشد. عباراتی نظیر «\lr{to boil with rage}», «\lr{to be steaming}», «\lr{to blow one's top}», «\lr{to keep a lid on it}» همگی این نگاشتِ بنیادی را عملیاتی می‌کنند.

نکته نظریِ بسیار شایان توجهی که در پژوهش حاضر برجسته خواهد شد آن است که قرآن کریم، قرن‌ها پیش از صورت‌بندیِ مدرنِ این استعاره، \textbf{دقیقاً همین ساختارِ شناختی را در سه‌گانه غیظ-کظم-تمیّز عملیاتی کرده است}. واژه «غَیظ» در لغت به «خشمِ متراکم در درون» اشاره می‌کند (یعنی محتوای فشرده درونِ ظرف)؛ فعلِ «کَظَمَ» در ریشه‌شناسیِ اولیه به «بستنِ دهانه مَشک برای جلوگیری از سرریزِ آب» اشاره دارد (یعنی عملیاتِ مهر کردنِ ظرف)؛ و عبارتِ «تَكَادُ تَمَيَّزُ مِنَ الْغَيْظِ» در آیه ۸ سوره ملک، چنان شدتی از خشم را تصویر می‌کند که گویی خودِ ظرف در حالِ از هم گسستن است (یعنی فروپاشیِ ساختارِ ظرف). شکل ۶ این نگاشتِ مفهومی را به‌نحوی تصویری بازنمایی می‌کند.

این هم‌سوییِ شایانِ توجه، در ادبیاتِ بین‌المللیِ زبان‌شناسیِ شناختی به‌قدرِ کافی برجسته نشده است. ادعای محتاطانه‌ی پژوهش حاضر آن است که آیه‌ی ۶۷:۸، نمونه‌ای \textbf{بسیار شفاف} از پدیده‌ی «فروپاشیِ ظرف» را در یک متنِ پیشامدرن عرضه می‌کند --- نمونه‌ای که در آن، فعلِ «تَمَیَّز» نه خروجِ محتوا بلکه گسستگیِ خودِ ظرف را تصویر می‌کند. ادعای قوی‌تر مبنی بر «قدیمی‌ترین» یا «کامل‌ترین» بودنِ این بازنمایی، نیازمندِ پیمایشِ تطبیقیِ نظام‌مند در ادبیاتِ پیشامدرنِ سامی، یونانی، چینی، و سنسکریت است که در محدوده‌ی پژوهش حاضر نمی‌گنجد و در §۶ به‌عنوان مسیرِ پژوهشِ آتی پیشنهاد می‌شود.

\subsubsection{۲-۴. هم‌گرایی نظری: چارچوبِ سه‌گانه}\label{ux647ux645ux6afux631ux627ux6ccux6cc-ux646ux638ux631ux6cc-ux686ux627ux631ux686ux648ux628-ux633ux647ux6afux627ux646ux647}

سه چارچوبِ نظریِ بالا، در پژوهش حاضر، در یک ساختارِ ترکیبی همگرا می‌شوند:

\begin{longtable}[]{@{}
  >{\raggedright\arraybackslash}p{(\linewidth - 4\tabcolsep) * \real{0.3333}}
  >{\raggedright\arraybackslash}p{(\linewidth - 4\tabcolsep) * \real{0.3333}}
  >{\raggedright\arraybackslash}p{(\linewidth - 4\tabcolsep) * \real{0.3333}}@{}}
\toprule\noalign{}
\begin{minipage}[b]{\linewidth}\raggedright
چارچوب
\end{minipage} & \begin{minipage}[b]{\linewidth}\raggedright
کارکرد در این پژوهش
\end{minipage} & \begin{minipage}[b]{\linewidth}\raggedright
محدودیتی که توسط چارچوب بعدی برطرف می‌شود
\end{minipage} \\
\midrule\noalign{}
\endhead
\bottomrule\noalign{}
\endlastfoot
ایزوتسو (میدانِ معنایی) & تعریفِ مرزهای میدان، شناسایی واژگانِ کانونی، توصیفِ روابطِ همنشینی-جانشینی & فقدانِ صورت‌بندیِ صریحِ شدت/درجه \\
کندی-ساسون (طیف) & درجه‌بندیِ شدت در درونِ میدان، صورت‌بندیِ بُعد و ترتیب & فقدانِ تبیینِ سازوکارِ شناختیِ صورت‌بندیِ شدت \\
لیکاف-کوچش (استعاره مفهومی) & تبیینِ سازوکارِ شناختیِ ساختاربخشِ پیوستار، به‌ویژه در مرحله ۵ (غیظ → تمیّز) & -- \\
\end{longtable}

هیچ‌یک از این سه چارچوب به‌تنهایی برای پاسخ‌گویی به سؤال‌های پنج‌گانه‌ی پژوهش کفایت نمی‌کند؛ ولی ادغامِ آن‌ها در یک ساختارِ سه‌لایه‌ای (شبکه → طیف → استعاره) ابزاری قدرتمند برای واکاویِ هندسه‌ی طیفِ هیجانیِ خشم در قرآن فراهم می‌آورد.

\begin{center}\rule{0.5\linewidth}{0.5pt}\end{center}

\begin{figure}
\centering
\includegraphics[width=1\linewidth,height=\textheight,keepaspectratio]{fig1_continuum.pdf}
\caption{پیوستارِ شش-مرحله‌ایِ طیفِ هیجانیِ خشم در قرآن کریم، تصویرشده روی یک محورِ افقیِ واحد از \emph{شدتِ کنش}. هر گره با یکی از چهارده ریشه‌ی عربیِ کانونی متناظر است؛ سطحِ گره با بسامدِ آن ریشه در پیکره‌ی ریخت‌شناختیِ قرآنی (دوکس، ۲۰۱۱) متناسب است. مرحله ۱ (پیش‌از‌خشمِ ناخشنودی---أف، کره)؛ مرحله ۲ (فشارِ درونی و انقباض---ضیق، حزن، أسف)؛ مرحله ۳ (نفرتِ ارزیابانه---نقم، سخط، مقت)؛ مرحله ۴ (خشمِ فعال---غضب، حرد)؛ مرحله ۵ (خشمِ متراکم/انفجاری---غیظ، با تمیّز به‌مثابه‌ی تجلی)؛ مرحله ۶ (پیامدهای رفتاری همراه با ملاحظات---بغی، طغیان، عتوّ). جرمِ بصریِ مرحله ۶ \lr{(n = 145)} با مجموعِ مراحلِ ۱ تا ۵ \lr{(n = 167)} قابلِ مقایسه است؛ نزدیکیِ این دو مقدار (\lr{χ²(1) = 1.55}، رد نمی‌شود)، با تفسیرِ علیتِ دوسویه‌ی مرحله ۶ که در §۴-۸ و §۵-۴ بسط می‌یابد، سازگار است.}
\end{figure}

\subsection{۳. روش‌شناسی}\label{ux631ux648ux634ux634ux646ux627ux633ux6cc}

\subsubsection{۳-۱. پیکره و منابع داده}\label{ux67eux6ccux6a9ux631ux647-ux648-ux645ux646ux627ux628ux639-ux62fux627ux62fux647}

پیکره داده‌های اصلیِ پژوهش، \textbf{پیکره ریخت‌شناختیِ قرآنی} (\lr{Quranic Arabic Corpus}، نسخه ۰٫۴) است که توسط کیس دوکس در دانشگاه لیدز در سال ۲۰۱۱ منتشر شده و تحت لیسانس GPL در دسترس است \lr{(Dukes 2011)}. این پیکره، برای هریک از ۱۲۸٬۲۱۹ پاره‌واژه قرآن، یک رکوردِ ساختاری شامل: (الف) موقعیت در قالب سوره:آیه:واژه:پاره (مثلاً (۱۵:۹۷:۲:۱))، (ب) صورتِ سطحی به ترانویسی بَکوالتر، (ج) برچسب جزء سخن \lr{(POS)}، و (د) ویژگی‌های ریخت‌شناختی شامل ریشه، شکل لمی، جنسیت، شمار، حالت، و گاهی مفاهیمِ معنایی، فراهم می‌آورد.

برای فراداده مکی/مدنیِ سور و متنِ قرآن به ویرایشِ عثمانیِ رسم‌الخط، از منابع پروژه تنزیل \lr{(Tanzil 2008)} استفاده شده است. این منبع، تحت لیسانس CC-BY-ND 3.0 منتشر شده و در میان پژوهشگرانِ قرآنیِ معاصر به‌مثابه استانداردِ متنیِ معیار شناخته می‌شود.

\subsubsection{۳-۲. واژگان مورد بررسی و معیارهای انتخاب}\label{ux648ux627ux698ux6afux627ux646-ux645ux648ux631ux62f-ux628ux631ux631ux633ux6cc-ux648-ux645ux639ux6ccux627ux631ux647ux627ux6cc-ux627ux646ux62aux62eux627ux628}

چهارده ریشه‌ی کانونی در پژوهش حاضر، بر اساسِ پنج معیارِ \emph{شمول و خروج} انتخاب شده‌اند:

\textbf{معیارهای شمول:}
\begin{enumerate}
\tightlist
\item \textbf{قرابتِ معناییِ مستقیم با میدانِ خشم}: حضورِ ریشه به‌مثابهِ بازنمایی‌کننده‌ی ناخشنودی، نفرت، خشم، یا پیامدهای رفتاریِ آن.
\item \textbf{وقوعِ قرآنی}: دستِ‌کم یک بسامد در پیکره‌ی ریخت‌شناختیِ قرآنی، با پذیرشِ هاپَکس‌ها (مانند «حَرد» در قلم: ۲۵) به‌شرطِ غنیِ بودنِ پروفایلِ تفسیری.
\item \textbf{پشتیبانیِ لغوی}: حضورِ ریشه به‌مثابهِ مدخلِ مستقل در دستِ‌کم سه منبع از چهار لغت‌نامه‌ی معیار (مقاييس، المفردات، لسان، التحقيق).
\end{enumerate}

\textbf{معیارهای خروج:}
\begin{enumerate}
\setcounter{enumi}{3}
\tightlist
\item \textbf{حذفِ ریشه‌های اخلاقیِ بسیار گسترده مگر در نقشِ چترِ مفهومی}: ریشه‌هایی نظیر «ظلم» (با ۳۱۵ بسامد) و «جرم» (۶۶) که گستره‌ی مفهومیِ بسیار وسیع‌تر از طیفِ خشم دارند، از پیوستارِ اصلی کنار گذاشته و تنها در آزمونِ پایداری (§۴-۹) به‌مثابهِ \emph{چترِ مفهومی} استفاده می‌شوند.
\item \textbf{راستی‌آزماییِ بافتاریِ دستی}: برای هر ریشه، حداقل پنج آیه‌ی نمونه به‌صورتِ دستی بازخوانی شده تا اطمینان حاصل شود که کاربردِ غالبِ آن در میدانِ خشم است (نه در میدانی هم‌نام مانند کراهت/إکراه که در ادامه تفکیک می‌شود).
\end{enumerate}

\textbf{ریشه‌های مرزی و توضیحِ گزینش:}

\begin{itemize}
\tightlist
\item \textbf{«کره»} در دو سامانه به‌کار می‌رود: کراهتِ درونی (\lr{Stage 1}، در میدانِ هیجانی) و إکراهِ ساختاری (سازوکارِ تحمیل، نه هیجان). تحلیلِ هر دو در §۴-۳ گزارش می‌شود.
\item \textbf{«شنأ»} با ۳ بسامد، به‌رغمِ معنایِ معناییِ مرتبط با نفرت، به دلیلِ بسامدِ بسیار اندک و نامتعارف بودنِ کاربردِ آیاتیِ آن، از پیوستارِ اصلی کنار گذاشته شد و در گروهِ چترِ مفهومی بررسی می‌شود.
\end{itemize}

افزون بر این چهارده ریشه، پنج ریشه‌ی گسترده‌تر---\textbf{شنأ، جرم، فسق، ظلم، عدو}---به‌عنوان «واژگانِ چترِ مفهومی» \lr{(umbrella terms)} نیز استخراج شده‌اند تا در آزمونِ پایداری (§۴-۹) مورد استفاده قرار گیرند.

\paragraph{خطِ پایه‌ی مکی/مدنی --- توضیحِ صریح.} برای آزمون‌های دوجمله‌ای که گرایشِ یک ریشه به یک دوره‌ی نزول را می‌سنجند، خطِ پایه‌ی ۰٫۷۴ (نسبتِ مکی-بودنِ پایه) استخراج می‌شود به‌صورتِ زیر:

\[
\text{baseline} = \frac{\text{تعداد آیاتِ مکی}}{\text{تعداد کلِ آیات}} = \frac{4{,}613}{6{,}236} \approx 0{.}74
\]

این یک خطِ پایه‌ی \emph{آیه-محور} (\lr{verse-level}) است و نه سوره-محور یا توکن-محور. توضیحِ گزینش: چون مقایسه‌ها در سطحِ \emph{بسامدِ ریشه در آیات} (نه در سوره‌ها یا توکن‌های سطحی) انجام می‌شود، باید خطِ پایه نیز در همین سطح باشد. تحلیلِ حساسیتِ خطِ پایه در دامنه‌ی ۰٫۷۰ تا ۰٫۷۸ (به‌منظورِ پوششِ عدمِ قطعیتِ سنتیِ رده‌بندیِ سور) نشان داد که نتیجه‌ی معناداریِ \lr{p < 0.01} برای سَخَط در سرتاسرِ این دامنه باقی می‌ماند.

\paragraph{راستی‌آزماییِ تخصیصِ مرحله‌ای.} هر یک از چهارده ریشه به یکی از شش مرحله تخصیص یافته است. این تخصیص بر مبنایِ معیارهای زیر، با همگرایی شواهد، صورت گرفته:

\begin{longtable}{p{0.10\linewidth}p{0.13\linewidth}p{0.32\linewidth}p{0.30\linewidth}}
\toprule
\textbf{ریشه} & \textbf{مرحله‌ی تخصیص‌یافته} & \textbf{شواهدِ لغوی-تفسیری} & \textbf{دلیلِ عدم تخصیص به مرحله‌ی مجاور} \\
\midrule
\endhead
\bottomrule
\endlastfoot
أف & ۱ & ابراز کلامیِ ناخشنودیِ خفیف؛ اسراء: ۲۳ & فاقدِ مؤلفه‌ی فشارِ درونی (مرحله ۲) \\
کره & ۱ & کراهتِ درونیِ پیش‌از‌خشم؛ بقره: ۲۱۶ & فاقدِ مؤلفه‌ی نفرتِ ارزیابانه‌ی شدید (مرحله ۳) \\
ضیق & ۲ & تنگیِ بدنی-روانی (ضیق الصدر)؛ حجر: ۹۷ & فاقدِ مؤلفه‌ی تقابلی-اخلاقی (مرحله ۳) \\
حزن & ۲ & اندوهِ آمیخته با فشارِ درونی؛ بقره: ۳۸ & تجلیِ غالبِ آن انفعالی است، نه فعال (مرحله ۴) \\
أسف & ۲ & اندوهِ همراه با خشم؛ زخرف: ۵۵ & محورِ بدنی-تنگی غالب نیست (همراه ضیق نیست) \\
نقم & ۳ & نکوهشِ همراه با اراده‌ی انتقام؛ اعراف: ۱۲۶ & فاقدِ بُعدِ خشمِ کنشیِ آنی (مرحله ۴) \\
سخط & ۳ & نارضایتیِ ارزیابانه‌ی شدید؛ محمد: ۲۸ & بُعدِ نهفته‌گی غالب است، نه ابرازِ کنشی \\
مقت & ۳ & نفرتِ اخلاقیِ تثبیت‌شده؛ غافر: ۱۰ & ارزیابانه و نه فوران‌کننده (مرحله ۵) \\
غضب & ۴ & خشمِ هدف‌مندِ ابراز-شده؛ فتح: ۶ & آنی و نه متراکم (مرحله ۵) \\
حرد & ۴ & خشمِ همراه با قصدِ مهار؛ قلم: ۲۵ & مؤلفه‌ی فعالِ کنش‌گرایانه دارد (نه صرفاً ارزیابی) \\
غیظ & ۵ & خشمِ متراکم/مهار-شده؛ آل عمران: ۱۳۴ & فشارِ درونی فراتر از مرحله ۴؛ کظم تجلیِ آن است \\
بغی & ۶ & تجاوزِ ساختاری؛ شوری: ۴۲ & علیتِ دوسویه با خشم (نه پیامدِ یک‌جانبه) \\
طغی & ۶ & سرکشیِ ساختاری؛ علق: ۶ & ریشه در استکبار، نه فقط در مهارِ ناموفقِ هیجان \\
عتو & ۶ & سرپیچیِ تثبیت‌شده؛ فرقان: ۲۱ & ویژگیِ شخصیتیِ پایدار، نه حالتِ گذرا \\
\end{longtable}

این جدول، عمومی‌بودنِ تصمیم‌های تخصیصی را قابلِ راستی‌آزماییِ مستقل می‌سازد و خواننده می‌تواند با مراجعه به منابع لغوی و تفسیری، مرتبه‌بندیِ بدیل پیشنهاد دهد.

\subsubsection{۳-۳. خط لوله محاسباتی}\label{ux62eux637-ux644ux648ux644ux647-ux645ux62dux627ux633ux628ux627ux62aux6cc}

سامانه استخراج به زبان پایتون نوشته شده و از پنج گام تشکیل می‌شود:

\begin{verbatim}
گام ۱: تجزیه پرونده QAC \lr{(analysis/qac_parser.py)}
       ↓
گام ۲: فیلتر بر اساسِ ریشه‌های هدف \lr{(analysis/spectrum_roots.py)}
       ↓
گام ۳: تجمیعِ پاره‌واژه‌ها به سطحِ واژه \lr{(collapse_to_words)}
       ↓
گام ۴: پیوند با فراداده مکی/مدنیِ سور
       ↓
گام ۵: تولید خروجیِ جدولی، آماری و شبکه‌ای
       \lr{(analysis/extract_concordance.py
        analysis/advanced_metrics.py)}
\end{verbatim}

تمام کد و داده‌ها در انباره‌ی عمومیِ پژوهش (به آدرسِ \lr{\url{https://github.com/ali-kin4/quranic-emotion-spectrum}}) در دسترس‌اند و نتایج با اجرای دو فرمان (\texttt{python\ analysis/extract\_concordance.py\ \&\&\ python\ analysis/visualize\_spectrum.py}) به‌طور کامل قابلِ بازتولید است. این رویکردِ بازتولیدپذیر، در مطالعاتِ قرآنیِ سنتی غایب بوده و یکی از سهم‌های روش‌شناختیِ پژوهش حاضر را تشکیل می‌دهد.

\subsubsection{۳-۴. ابزارهای تحلیلی}\label{ux627ux628ux632ux627ux631ux647ux627ux6cc-ux62aux62dux644ux6ccux644ux6cc}

برای تحلیلِ یافته‌ها، از مجموعه‌ای از ابزارهای کمّی و کیفی به‌صورت همگرا بهره گرفته می‌شود:

\textbf{کمّی:}
- شمارش بسامدِ خام و توزیعِ مکی/مدنی هر ریشه
- آزمون خی-دو (χ²) برای راستی‌آزماییِ فرضیه‌های توزیعی (ف.۱، ف.۲)
- آزمون دوجمله‌ای دوسویه برای راستی‌آزماییِ ادعاهای اختصاصیتِ توزیعی (ف.۳)
- \textbf{راستی‌آزماییِ استواری در شمارش‌های اندک:} برای ریشه‌هایی با بسامدِ بسیار کم (نظیر سَخَط با ۴، عتوّ با ۱۰، أسف با ۵، حَرد با ۱)، از مجموعه‌ای از آزمون‌های جایگزین برای کنترلِ خطای نوع اول استفاده می‌شود: \lr{(i)} آزمون دقیقِ دوجمله‌ای (به‌جای تقریبِ نرمال)، \lr{(ii)} تحلیلِ حساسیتِ خطِ پایه (در دامنه ۰٫۷۰ تا ۰٫۷۸ به جای نقطه‌ی واحدِ ۰٫۷۴)، \lr{(iii)} گزارشِ توانِ آماری برای \lr{n < 10}، و \lr{(iv)} هشدارِ صریح در گزارشِ نتایج برای ریشه‌هایی که \lr{n < 5}؛ تفسیرِ هر یافته‌ی آماری در این رژیم، صرفاً به‌عنوانِ شواهدِ اولیه و نه شواهدِ قطعی، صورت می‌پذیرد.
- تحلیلِ ریخت‌شناختی: بخش‌بندیِ فعل/اسم/صفت/اسم فاعل-مفعول
- تحلیل شبکه: ساختِ گراف هم‌رخدادگیریِ آیه‌ای، محاسبه مرکزیتِ درجه‌ای \lr{(degree)}، بینابینی \lr{(betweenness)}، نزدیکی \lr{(closeness)}، و بردارِ ویژه \lr{(eigenvector)}.

\paragraph{مشخصاتِ کاملِ تحلیلِ شبکه (\lr{network specification}).} برای رفعِ ابهام و تضمینِ بازتولیدپذیری، مؤلفه‌های گرافِ هم‌رخدادگیری به‌صورت زیر تعریف می‌شوند:

\begin{longtable}{p{0.30\linewidth}p{0.62\linewidth}}
\toprule
\textbf{مؤلفه} & \textbf{تعریفِ عملیاتی} \\
\midrule
\endhead
\bottomrule
\endlastfoot
گره \lr{(node)} & ریشه‌ی عربی (نه واژه‌ی صرف‌شده). در سامانه‌ی پایه: ۱۴ ریشه‌ی کانونی؛ در آزمونِ پایداری: ۱۴ + ۵ ریشه‌ی چترِ مفهومی. \\
یال \lr{(edge)} & حضورِ دو ریشه در یک آیه‌ی واحد (نه در سوره‌ی واحد). \\
وزنِ یال \lr{(edge weight)} & تعداد آیاتی که هر دو ریشه در آن‌ها هم‌زمان حاضرند (شمارشِ خام، بدونِ نرمال‌سازی). \\
نوعِ گراف & جهت‌دار-نشده (\lr{undirected}); ساده (بدون یالِ موازی)؛ بدونِ خود-حلقه (\lr{self-loop} نشده). \\
نرم‌افزار & \lr{NetworkX 3.x} در پایتون \lr{3.10+}. \\
درجه‌ی مرکزیتِ خام & شمارشِ یال‌های متصل به هر گره. \\
مرکزیتِ بینابینی & فرمولِ \lr{NetworkX.betweenness\_centrality} با \lr{normalized=True} (نرمال‌سازی به ضریبِ \lr{$2/((n-1)(n-2))$}). \\
مرکزیتِ نزدیکی & \lr{NetworkX.closeness\_centrality} با نرمال‌سازیِ \lr{Wasserman--Faust}. \\
مرکزیتِ بردارِ ویژه & توانِ ۱۰۰ تکرار با \lr{tol=1e-6}; در صورتِ عدمِ همگرایی، مقدار \lr{NaN} ثبت می‌شود (پایپ‌لاین این حالت را مدیریت می‌کند). \\
هم‌رخدادگیریِ صفر-بسامد & ریشه‌هایی که هیچ هم‌رخدادگیری ندارند (نظیر حَرد به‌عنوان هاپَکس) به‌مثابه گره‌های جدا \lr{(isolated)} باقی می‌مانند. \\
\end{longtable}

این مشخصات مستقیماً در فایلِ \lr{analysis/advanced\_metrics.py} پیاده‌سازی شده‌اند.

\textbf{کیفی:}
- ریشه‌شناسیِ تطبیقی در چهار لغت‌نامه‌ی معیار (مقاییس، مفردات، لسان، تحقیق)
- تحلیلِ تفسیریِ تطبیقی برای آیات کانونی، با تکیه بر چهار تفسیرِ معتبر (المیزان طباطبائی، الکشاف زمخشری، مجمع البیان طبرسی، التحریر و التنویر ابن عاشور)
- تحلیلِ پدیدارشناختیِ نشانه‌های متنی (نشانگرها: ضیق الصدر، حزن القلب، أسف النفس، ضاق ذرعاً، و غیره)
- تحلیلِ ترجمه‌ایِ تطبیقیِ آیات کانونی در پنج ترجمه فارسی معیار

\subsubsection{۳-۵. روایی پیکره و راستی‌آزمایی}\label{ux631ux648ux627ux6ccux6cc-ux67eux6ccux6a9ux631ux647-ux648-ux631ux627ux633ux62aux6ccux622ux632ux645ux627ux6ccux6cc}

برای اطمینان از صحتِ خط لوله، چهارده آیه‌ی شاخص (یکی برای هر ریشه‌ی کانونی) به‌صورتِ دستی در برابر خروجیِ پیکره راستی‌آزمایی شدند:

\begin{itemize}
\tightlist
\item
  اسراء:۲۳ (أف)؛ بقره:۲۱۶ (کره)؛ حجر:۹۷ (ضیق)؛ بقره:۳۸ (حزن)؛ زخرف:۵۵ (أسف)؛ اعراف:۱۲۶ (نقم)؛ محمد:۲۸ (سخط)؛ غافر:۱۰ (مقت)؛ فتح:۶ (غضب)؛ قلم:۲۵ (حرد)؛ آل عمران:۱۳۴ (غیظ، با تجلیِ آن در ملک:۸ به‌صورتِ تمیّز)؛ شوری:۴۲ (بغی)؛ علق:۶ (طغیان)؛ فرقان:۲۱ (عتوّ).
\end{itemize}

هر چهارده آیه با موفقیت در پیکره خروجی، با صورتِ سطحیِ صحیح و متنِ مطابق با مصحفِ عثمانی، یافت گردید. این راستی‌آزمایی، اطمینان می‌دهد که الگوریتمِ ترانویسی بکوالتر و فرآیندِ تجمیعِ پاره‌واژه‌ها بدونِ خطای سامانمند عمل می‌کند.

\begin{center}\rule{0.5\linewidth}{0.5pt}\end{center}

\subsection{۴. یافته‌ها}\label{ux6ccux627ux641ux62aux647ux647ux627}

تحلیلِ چهارده ریشه‌ی کانونیِ میدانِ طیف هیجانی خشم در پیکره‌ی ریخت‌شناختیِ قرآنی، شواهدی منسجم و چندلایه برای پیوستارِ پیشنهادی فراهم می‌آورد. در ادامه، یافته‌های هر مرحله به‌تفصیل و در پیوستِ منابعِ لغویِ تطبیقی، تحلیلِ تفسیری، و داده‌های توزیعی-آماری-شبکه‌ای، ارائه می‌شود.

\subsubsection{۴-۱. ریشه‌شناسیِ تطبیقی: جدول معیار چهار-لغت‌نامه‌ای}\label{ux631ux6ccux634ux647ux634ux646ux627ux633ux6cc-ux62aux637ux628ux6ccux642ux6cc-ux62cux62fux648ux644-ux645ux639ux6ccux627ux631-ux686ux647ux627ux631-ux644ux63aux62aux646ux627ux645ux647ux627ux6cc}

پیش از ورود به تحلیلِ مرحله‌ای، جدول ۱ معنای کانونیِ هر ریشه را در چهار لغت‌نامه‌ی معیار عربی---\textbf{معجم مقاييس اللغة} ابن فارس (د. ۳۹۵ ق)، \textbf{المفردات في غريب القرآن} راغب اصفهانی (د. ۵۰۲ ق)، \textbf{لسان العرب} ابن منظور (د. ۷۱۱ ق)، و \textbf{التحقيق في كلمات القرآن} مصطفوی (معاصر)---فهرست می‌کند.

\textbf{جدول ۱.} ریشه‌شناسیِ تطبیقی چهارده ریشه‌ی کانونی در چهار لغت‌نامه‌ی معیار

\begin{longtable}[]{@{}
  >{\centering\arraybackslash}p{(\linewidth - 10\tabcolsep) * \real{0.1667}}
  >{\centering\arraybackslash}p{(\linewidth - 10\tabcolsep) * \real{0.1667}}
  >{\raggedright\arraybackslash}p{(\linewidth - 10\tabcolsep) * \real{0.1667}}
  >{\raggedright\arraybackslash}p{(\linewidth - 10\tabcolsep) * \real{0.1667}}
  >{\raggedright\arraybackslash}p{(\linewidth - 10\tabcolsep) * \real{0.1667}}
  >{\raggedright\arraybackslash}p{(\linewidth - 10\tabcolsep) * \real{0.1667}}@{}}
\toprule\noalign{}
\begin{minipage}[b]{\linewidth}\centering
مرحله
\end{minipage} & \begin{minipage}[b]{\linewidth}\centering
ریشه
\end{minipage} & \begin{minipage}[b]{\linewidth}\raggedright
مقاییس (ابن فارس)
\end{minipage} & \begin{minipage}[b]{\linewidth}\raggedright
المفردات (راغب)
\end{minipage} & \begin{minipage}[b]{\linewidth}\raggedright
لسان العرب (ابن منظور)
\end{minipage} & \begin{minipage}[b]{\linewidth}\raggedright
التحقیق (مصطفوی)
\end{minipage} \\
\midrule\noalign{}
\endhead
\bottomrule\noalign{}
\endlastfoot
۱ & أ-ف-ف & تَضَجُّر، نشانه‌ی کلامیِ ناپسندی {[}تکمیل از منابع لغوی{]} & کم‌ترین ابرازِ کلامیِ ناخشنودی & کلمه‌ی تَضَجُّر در برابرِ شخص یا شیء & کمترین نشانه‌ی کلامی از نفرتِ درونی \\
۱ & ک-ر-ه & کراهت، یافتنِ سنگین؛ اِکراه (تحمیل) & کراهتِ نفس؛ امتداد ساختاریِ آن به اِکراه & آنچه طبع یا عقل از آن گریزان است & کراهتِ درونی در برابرِ اِکراهِ ساختاری (دومی هیجان نیست) \\
۲ & ض-ي-ق & تنگی، عدمِ گشایش & حالتِ فشاری بر فرد، در برابرِ سعه & کم بودنِ مساحت، از سعه به ضیق & فقدانِ گشایش (بَسط) درونی یا بیرونی \\
۲ & ح-ز-ن & ناهمواری زمین → سنگینیِ روان & اندوهِ مقابلِ فرح & زمینِ خشن، خشونتِ نفس & تأثرِ ادراکی-عاطفی از فقدان \\
۲ & أ-س-ف & اندوهِ شدید آمیخته با خشم & شدت غضب در همراهیِ حزن & به نهایتِ اندوه و خشم رسیدن & ترکیبِ حزن و غضبِ نسبی \\
۳ & ن-ق-م & نکوهش همراه با اراده‌ی انتقام {[}تکمیل از منابع لغوی{]} & اعتراض و طلبِ کیفر در برابرِ کنش & استنکارِ شدید؛ مقایسه با المنتقم به‌مثابهِ صفتِ الهی & خشمِ ارزیاب با جهت‌گیریِ انتقام \\
۳ & س-خ-ط & کراهتِ شدید و ردِ اخلاقی & ضدِّ رضا با شدت & عدمِ خشنودیِ شدید با حالتِ سرزنشی & ابرازِ نفرتِ همراه با داوریِ اخلاقی \\
۳ & م-ق-ت & شدیدترین گونه‌ی إبغاض (زجّاج، اللیث در تهذیب اللغة) & بُغضِ ناشی از مشاهده‌ی فعلِ قبیح (راغب) & أَشَدُّ الإِبغاض؛ روی‌گردانی-فاصله‌گیری به‌جای تهاجم & نفرتِ صرفاً ارزیابانه‌ی اخلاقی \\
۴ & غ-ض-ب & شدتِ هیجان و حرکت در درون & ثوران دم القلب طلباً للانتقام & تغییرِ حالتِ روانیِ جوشنده & هیجانِ جهت‌مند با امکانِ انتقام \\
۴ & ح-ر-د & خشمِ همراه با قصد؛ نیتِ منع & بنا بر طبرسی در مجمع البیان: غضب و قصد & بنا بر طباطبائی در المیزان: خشم به‌علاوه‌ی منعِ آگاهانه {[}تکمیل از منابع لغوی{]} & خشمِ مسلح‌شده با امساک (آیه‌ی ۲۵ سوره‌ی قلم) \\
۵ & غ-ي-ظ & غلیان و فشار درونی & شدید‌ترین حالاتِ غضب پنهان & امتلاءِ نفس از خشمِ مهارشده & تراکمِ هیجانی در آستانه‌ی انفجار \\
۶ & ب-غ-ي & طلبِ فراتر رفتن از حق & تجاوز و اخراج از طریقِ حق & استیلاءِ به ستم & کنشِ تجاوزآمیز با هدفِ ستم \\
۶ & ط-غ-ي & سرریز شدنِ آب از حدّ & عبور از مقدارِ مأذون & تعدّی شدید از حد & شکستنِ مرزِ هنجارها به‌صورتِ سرکش \\
۶ & ع-ت-و & سرکشیِ سختِ همراه با تکبر & شدتِ کبر و عناد در عصیان & استکبار به‌نهایت، در پیوند با اِستِکبار & تثبیتِ موضعِ سرکشانه در برابرِ حق (ویژگیِ پایدارِ شخصیت) \\
\end{longtable}

این جدول، پشتوانه‌ی لغوی-تاریخیِ پیوستارِ پیشنهادی را فراهم می‌آورد و نشان می‌دهد که گذار از أُف به عتوّ، نه ساخته‌ی حدسیِ پژوهش حاضر، بلکه بازتابی از ساختارِ مفهومیِ تثبیت‌شده در سنتِ چهار لغت‌نامه‌ی معیار است.

\subsubsection{۴-۲. یافته‌های توزیعیِ کلی}\label{ux6ccux627ux641ux62aux647ux647ux627ux6cc-ux62aux648ux632ux6ccux639ux6cc-ux6a9ux644ux6cc}

جدول ۲، نمای کلیِ توزیعِ چهارده ریشه‌ی کانونی را ارائه می‌دهد. این جدول، شالوده‌ی عددیِ تحلیل‌های مرحله‌ای و آماریِ بعدی را تشکیل می‌دهد.

\textbf{جدول ۲.} نمای توزیعی چهارده ریشه‌ی کانونی در پیکره ریخت‌شناختیِ قرآن

\begin{longtable}[]{@{}
  >{\centering\arraybackslash}p{(\linewidth - 10\tabcolsep) * \real{0.1667}}
  >{\centering\arraybackslash}p{(\linewidth - 10\tabcolsep) * \real{0.1667}}
  >{\centering\arraybackslash}p{(\linewidth - 10\tabcolsep) * \real{0.1667}}
  >{\centering\arraybackslash}p{(\linewidth - 10\tabcolsep) * \real{0.1667}}
  >{\centering\arraybackslash}p{(\linewidth - 10\tabcolsep) * \real{0.1667}}
  >{\raggedright\arraybackslash}p{(\linewidth - 10\tabcolsep) * \real{0.1667}}@{}}
\toprule\noalign{}
\begin{minipage}[b]{\linewidth}\centering
مرحله
\end{minipage} & \begin{minipage}[b]{\linewidth}\centering
ریشه
\end{minipage} & \begin{minipage}[b]{\linewidth}\centering
بسامد کل
\end{minipage} & \begin{minipage}[b]{\linewidth}\centering
مکی
\end{minipage} & \begin{minipage}[b]{\linewidth}\centering
مدنی
\end{minipage} & \begin{minipage}[b]{\linewidth}\raggedright
یادداشت
\end{minipage} \\
\midrule\noalign{}
\endhead
\bottomrule\noalign{}
\endlastfoot
۱ & أف & ۳ & ۳ & ۰ & همگی مکی؛ کم‌ترین تَرازِ کلامیِ ناخشنودی \\
۱ & کره & ۴۱ & ۱۴ & ۲۷ & گرایشِ مدنی \lr{p = 0.005}؛ کراهت/إکراه دو‌گانه \\
۲ & ضيق & ۱۳ & ۹ & ۴ & تنگی درونی (ضِيقُ الصدر؛ ضاق ذرعاً) \\
۲ & حزن & ۴۲ & ۲۶ & ۱۶ & پربسامدترینِ مرحله ۲؛ ولی ارزشِ معنایی فراگذرنده (نک. §۴-۴-۲-۱) \\
۲ & أسف & ۵ & ۵ & ۰ & انحصاراً مکی \\
۳ & نقم & ۱۷ & ۱۲ & ۵ & گرایش مکی؛ نکوهشِ همراه با اراده‌ی انتقام؛ المنتقم \\
۳ & سخط & ۴ & ۰ & ۴ & انحصاراً مدنی \lr{p = 0.005} \\
۳ & مقت & ۶ & ۴ & ۲ & شدیدترین نفرتِ اخلاقی (زجّاج) \\
۴ & غضب & ۲۴ & ۱۳ & ۱۱ & واژه‌ی کانونیِ خشم؛ خداوند، انبیا، مؤمنان \\
۴ & حرد & ۱ & ۱ & ۰ & هاپَکس قرآنی: قلم: ۲۵ \\
۵ & غيظ & ۱۱ & ۳ & ۸ & خشمِ متراکم؛ تجلیِ تمیّز در ملک: ۸ \\
۶ & بغي & ۹۶ & ۵۵ & ۴۱ & گرهِ پلی؛ علیتِ دوسویه (نک. §۴-۸-۵) \\
۶ & طغي & ۳۹ & ۲۹ & ۱۰ & پارادایمِ فرعونی؛ استکبارِ ساختاری \\
۶ & عتو & ۱۰ & ۹ & ۱ & عنادِ تثبیت‌شده؛ پیوسته با استکبار \\
\textbf{مجموع} & -- & \textbf{۳۱۲} & \textbf{۱۸۳} & \textbf{۱۲۹} & -- \\
\end{longtable}

تحلیلِ توزیعِ مرحله‌ای، بسامدها را به‌صورت زیر تجمیع می‌کند: مرحله ۱ = ۴۴؛ مرحله ۲ = ۶۰؛ مرحله ۳ = ۲۷؛ مرحله ۴ = ۲۵؛ مرحله ۵ = ۱۱؛ مرحله ۶ = ۱۴۵.

\textbf{آزمونِ خی-دوی فرضِ توزیعِ شش‌مرحله‌ایِ یکنواخت:} \lr{χ²(df=5) = 227.15, p < 0.001} --- مقادیرِ مشاهده‌شده به‌طور قاطع از فرضِ یکنواختی منحرف‌اند. تأیید فرضیه ف.۱.

\textbf{آزمونِ هسته‌ی پدیدارشناختی (مراحل ۱ تا ۵ = ۱۶۷) در برابرِ قطبِ پیامدهای رفتاری (مرحله ۶ = ۱۴۵):} \lr{χ²(df=1) = 1.55}، \textbf{رد نمی‌شود}. بنا بر اصلِ آماریِ متعارف، عدمِ ردِ فرضیه‌ی صفر، \emph{فقدانِ شواهدِ تفاوت} است و نه شواهدِ مساوات؛ این آزمون، به‌تنهایی، هیچ معماریِ علّیِ خاصی را اثبات نمی‌کند. نزدیکیِ مشاهده‌شده، با ملاحظه‌ی مفسرانه‌ی §۴-۸ \emph{سازگار} است: ریشه‌های مرحله ۶ (بغی، طغیان، عتوّ) تنها خاستگاهی خشمی ندارند، بلکه می‌توانند به‌صورتِ مستقل از مسیرهای ساختاریِ غیرهیجانی (طمع، علوّ فی الأرض، استکبار) برخیزند. این قرائت، نیازمندِ پشتیبانیِ تفسیری و بافتاری است که در §۴-۸ ارائه می‌شود.

\begin{figure}
\centering
\includegraphics[width=0.95\linewidth,height=\textheight,keepaspectratio]{fig2_frequency_by_stage.pdf}
\caption{بسامدِ ریشه‌ای و مرحله‌ایِ ۳۱۲ بسامدِ کانونیِ طیف. \emph{پنلِ راست} (مقیاسِ خطی، به‌تفکیکِ ریشه): رنگِ هر میله کدِ مرحله‌ی مربوطه را نشان می‌دهد. مرحله ۶ نیز عدمِ تقارنِ نمایانی دارد (بغی ۹۶ در برابر عتو ۱۰) که بازتابِ تقسیمِ کارِ واژگانی میانِ پرخاشگریِ توزیع‌شده (بغی) و سرکشیِ لجوجانه (عتو) است که در §۴-۸ تحلیل می‌شود. \emph{پنلِ چپ} (مجموعِ هر مرحله): مرحله‌ها به‌ترتیب: ۴۴، ۶۰، ۲۷، ۲۵، ۱۱، ۱۴۵. آزمونِ شش‌مرحله‌ایِ خی-دو فرضِ یکنواختی را رد می‌کند (\lr{χ²(5) = 227.15, p < 0.001})؛ ولی مقایسه‌ی هسته‌ی پدیدارشناختی (مراحلِ ۱ تا ۵ = ۱۶۷) با قطبِ پیامدهای رفتاری (مرحله ۶ = ۱۴۵) رد نمی‌شود (\lr{χ²(1) = 1.55})، یافته‌ای که هم‌سو با ملاحظه‌ی مفسرانه‌ی §۴-۸ و §۵-۴ است: ریشه‌های مرحله ۶ صرفاً پیامدِ یک‌جانبه‌ی خشم نیستند.}
\end{figure}

\begin{figure}
\centering
\includegraphics[width=0.95\linewidth,height=\textheight,keepaspectratio]{fig3_meccan_medinan.pdf}
\caption{توزیعِ تجمیعیِ مکی/مدنی برای هر یک از چهارده ریشه‌ی کانونی، به‌صورتِ نمودارِ میله‌ایِ اشباع‌شده. اعدادِ بالای هر میله، نسبتِ خامِ مکی/مدنی را نشان می‌دهند. الگوهای توزیعیِ شایانِ تأمل: \emph{أَسَف} (مرحله ۲) \textbf{به‌صورتِ انحصاری} در سور مکی ظاهر می‌شود (۵/۰)؛ \emph{سَخَط} (مرحله ۳) \textbf{به‌صورتِ انحصاری} در سور مدنی ظاهر می‌شود (۰/۴؛ \lr{p = 0.005} در برابرِ خطِ پایه‌ی پیکره ۰٫۷۴)؛ \emph{نقم} (مرحله ۳) گرایشِ مکی دارد (۱۲/۵)؛ \emph{مقت} (مرحله ۳) متوازن (۴/۲)؛ \emph{حَرد} هاپَکس قرآنی در آیه‌ی ۲۵ سوره‌ی قلم (مکی)؛ \emph{کره} گرایشِ مدنی (۱۴/۲۷، \lr{p = 0.005}). تصویرسازی، یافته‌ای غیرپیش‌پاافتاده را رمزگذاری می‌کند: همه‌ی ریشه‌های طیف، منبعِ واژگانیِ بافت‌خنثی نیستند --- چندین ریشه از آن‌ها \emph{تخصص‌یافتگیِ گفتمانی-بافتی} دارند، الگویی که در §۵-۵ به‌مثابهِ ویژگیِ عمومیِ واژگانِ قرآنی نظریه‌پردازی می‌شود.}
\end{figure}

\begin{figure}
\centering
\includegraphics[width=0.95\linewidth,height=\textheight,keepaspectratio]{fig5_morphology_stack.pdf}
\caption{پروفایلِ ریخت‌شناختیِ چهارده ریشه‌ی کانونی به‌تفکیکِ جزء سخن (فعل/اسم/صفت/اسم فاعل). شیبِ چپ‌به‌راست در نمودار، یک \emph{جابه‌جاییِ گونه‌شناختی} در سراسرِ پیوستار را نشان می‌دهد که شواهدِ پیکره‌ایِ مستقلی برای فرضیه‌ی شدتِ کنش فراهم می‌آورد. مراحلِ ۱ و ۲ به‌صورتِ غالب فعلی‌اند --- حزن ۸۸٪ فعلی (۳۷ فعل / ۵ اسم)، ضيق ۶۲٪ --- هم‌خوان با هستی‌شناسیِ نشانه‌های پیشاخشمی به‌مثابهِ \emph{حالاتِ گذرا و موقعیت‌انگیخته}. مرحله ۴ نومینالیزه می‌شود: غضب با ۱۶ اسم در برابرِ ۶ فعل، نشان‌دهنده‌ی آن است که این واژه به یک ارجاعِ پایدار و قابلِ انتساب («غضبِ خداوند») تبدیل شده. مرحله ۶ بیشترین درجه‌ی نومینالیزاسیون را دارد، با ظهورِ صورت‌های اسمِ فاعل (طاغوت، باغی)؛ واژگانِ پیامدِ رفتاری به‌مثابهِ ویژگیِ پایدارِ کنشگر تثبیت می‌گردند، نه فرآیندِ گذرا. این شیبِ ریخت‌شناختی --- از فعل‌محوریِ مراحلِ ۱ و ۲ به اسم‌محوریِ مرحله ۶ --- شاهدی پیکره-درونی برای پیوستارِ شدتِ کنش است.}
\end{figure}

\subsubsection{۴-۳. مرحله ۱ --- پیش‌از‌خشمِ ناخشنودی}\label{ux645ux631ux62dux644ux647-ux6f1-ux67eux6ccux634ux627ux632ux62eux634ux645-ux646ux627ux62eux634ux646ux648ux62fux6cc}

مرحله نخست از پیوستار، با دو ریشه (أف، کره) آغاز می‌شود؛ ریشه‌هایی که به ترازِ پیش‌از‌خشمِ هیجان تعلق دارند: ابرازِ کلامیِ تَضَجُّر در أف، و کراهتِ درونی در کره.

\textbf{۴-۳-۰. أُف (۳ بسامد؛ همگی مکی).} این ریشه، که در منابعِ لغوی به «تَضَجُّر» (آزردگیِ کلامی) ترجمه می‌شود، کم‌ترین تَرازِ ابرازِ کلامیِ ناخشنودی در دستگاهِ قرآنی است. هر سه بسامدِ آن جایگاه‌های پارادایمیِ اخلاقی-بلاغی دارند:

\begin{itemize}
\tightlist
\item
  اسراء: ۲۳، که گفتنِ أف به والدین را نهی می‌کند («فَلَا تَقُل لَّهُمَا أُفٍّ»): \textbf{کم‌ترین} نشانه‌ی کلامیِ ناخشنودی، در بالاترین رابطه‌ی اخلاقی، بارِ هنجاری می‌گیرد.
\item
  انبیاء: ۶۷، که حضرت ابراهیم (ع) در رویاروییِ با قومِ بت‌پرستِ خود می‌گوید: «أُفٍّ لَّكُمْ وَلِمَا تَعْبُدُونَ مِن دُونِ اللَّهِ» --- اعتراضِ پیامبرانه.
\item
  احقاف: ۱۷، که أُف از زبانِ فرزندِ گریزان از والدینِ مؤمن گفته می‌شود --- نقطه‌ی عکسِ آیه‌ی اسراء.
\end{itemize}

به‌مثابهِ کم‌ترین تَرازِ پیوستار، أف یک نقطه‌ی بنیادی از گفتمانِ قرآنی را تثبیت می‌کند: حتی کوچک‌ترین نشانه‌ی کلامیِ ناخشنودی، اخلاقی‌اند و وزنِ اخلاقی دارند. بنابراین این ریشه، مرزِ پایینیِ پیوستار را در ترازی \textbf{زیرِ خشم} ولی \textbf{بالای بی‌اعتنایی} پایه‌ریزی می‌کند.

\textbf{۴-۳-۱. کره (۴۱ بسامد، ۱۴ مکی / ۲۷ مدنی؛ آزمونِ دوجمله‌ای \lr{p = 0.005} به‌سویِ مدنی).} کره ساختاری دو‌گانه دارد که باید تحلیلاً تفکیک شود:

(الف) \textbf{کراهتِ درونی (کراهت).} کراهت به‌مثابهِ یک حالتِ عاطفی، نفرتِ ارزیابانه‌ی درونیِ کنشگر نسبت به یک شیء، شخص، یا مسیرِ کنش است. آیه‌ی ۲۱۶ سوره‌ی بقره، «عَسَى أَن تَكْرَهُوا شَيْئًا وَهُوَ خَيْرٌ لَّكُمْ»، کراهتِ درونی را به‌مثابهِ یک حالتِ عاطفیِ \textbf{اصلاح‌پذیر} صورت‌بندی می‌کند: ارزشِ معرفتیِ آن نسبت به موضوع می‌تواند با دانش یا اصلاحِ الهی بازنگری گردد. این، به‌درستی، معنایی است که در میدانِ هیجانیِ مرحله ۱ قرار می‌گیرد.

(ب) \textbf{اِکراهِ ساختاری (اکراه).} اکراه \textbf{هیجان نیست} بلکه سازوکارِ تحمیل است: امتدادِ ساختاریِ نفرتِ یک طرف به اجبارِ طرفِ دیگر. آیه‌ی ۲۵۶ سوره‌ی بقره، «لَا إِكْرَاهَ فِي الدِّينِ»، حدِّ اخلاقیِ این امتدادِ ساختاری را تعیین می‌کند؛ آیه‌ی ۱۰۶ سوره‌ی نحل («إِلَّا مَنْ أُكْرِهَ») استثنای حقوقیِ شهادت تحتِ اِکراه را فراهم می‌آورد؛ و آیه‌ی ۳۳ سوره‌ی نور («وَلَا تُكْرِهُوا فَتَيَاتِكُمْ عَلَى الْبِغَاءِ») اِکراهِ کنیزان را نهی می‌کند. ما کره را در پیوستار، بر اساسِ معنای کراهت می‌گنجانیم؛ معنای اکراه از دایره‌ی میدانِ عاطفی بیرون است ولی به‌مثابهِ یک هم‌ریشه‌ی ساختاری، گزارشِ آن لازم است زیرا همان امتدادِ رفتاریِ نفرت را عملیاتی می‌کند --- پدیده‌ای که در مرحله ۶ نیز به‌صورتی متفاوت ظاهر می‌گردد.

\subsubsection{۴-۴. مرحله ۲ --- فشارِ درونی و انقباض}\label{ux645ux631ux62dux644ux647-ux6f2-ux641ux634ux627ux631-ux62fux631ux648ux646ux6cc-ux648-ux627ux646ux642ux628ux627ux636}

مرحله‌ی دوم، با سه ریشه (ضيق، حزن، أسف) و در مجموع ۶۰ بسامد، شکل می‌گیرد. هر یک، یک نمایندگیِ پدیدارشناختیِ متمایز در خانواده‌ی فشارِ درونی است.

\textbf{۴-۴-۱. ضیق (۱۳ بسامد، ۹ مکی / ۴ مدنی).} این ریشه، بنیادی‌ترین مفهومِ فشارِ درونی را ارائه می‌کند. ابن فارس در \textbf{مقاییس اللغة} (ج۳، ص۳۸۲) معنای کانونیِ آن را «تنگی، عدمِ گشایش» می‌داند و راغب در \textbf{المفردات} آن را در تقابلِ مستقیم با «سعه» (گشایش) قرار می‌دهد. در قرآن، این ریشه عمدتاً در ترکیبِ «ضِيق الصدر» (تنگیِ سینه) ظاهر می‌شود؛ ترکیبی که خود، یک نشانه پدیدارشناختیِ صریح را در بر دارد: مکانِ تجلیِ این هیجان، در قلمروی بدنِ زنده، \textbf{سینه} است.

آیه ۹۷ سوره حِجر، شاخص‌ترین کاربرد این ریشه را عرضه می‌کند: «وَلَقَدْ نَعْلَمُ أَنَّكَ يَضِيقُ صَدْرُكَ بِمَا يَقُولُونَ» (حقا که ما می‌دانیم که سینه‌ات از آنچه می‌گویند، تنگ می‌شود). این آیه، که در فضای تکذیب‌گریِ مکیِ اولیه نازل شده، نه به یک خشمِ فعال، بلکه به یک فشارِ روانیِ درونی اشاره دارد که در پاسخ به انکارِ اجتماعی شکل می‌گیرد. علامه طباطبایی در \textbf{المیزان} (ج۱۲، ص۱۹۵) به‌درستی توجه می‌دهد که در این آیه، خداوند به پیامبر (ص) \textbf{علم} خود را نسبت به این تنگی ابراز می‌کند، نه \textbf{رضایت} از آن یا \textbf{توصیه} بر مهارش؛ یعنی این تنگی، یک پدیده طبیعی و قابلِ پذیرشِ روانی است، نه یک خطای رفتاری.

نکته‌ی معناشناختیِ ظریفی که در ۸ مورد از ۱۳ بسامد قابلِ مشاهده است، آن است که فعلِ ضیق (با ۸ بسامدِ فعلی در مقابل ۳ اسمی، نسبت ۲٫۷:۱) در قرآن غالباً به‌صورتِ مضارع («یضیق»، «ضاق»، «یضیقون») به‌کار می‌رود. این ریخت‌شناسیِ غالباً فعلی، یک ویژگیِ پدیدارشناسانه را آشکار می‌سازد: ضیق، نه یک حالتِ پایدارِ شخصیتی (که با اسم بیان می‌شد)، بلکه یک واکنشِ موقتی و موقعیتی (که با فعل بیان می‌شود) است.

\textbf{۴-۴-۱-۱. ضاق ذرعاً --- لوط، آیه‌های ۷۷ و ۸۰ سوره‌ی هود --- پرخاشگریِ بازداشته.} یک کاربردِ ترکیبیِ بسیار مهمِ ضیق در روایتِ لوط (ع) رخ می‌دهد و یک مدلِ چهارگامیِ خشمِ درونی را عرضه می‌کند:

(۱) \textbf{آغازگر}: خودِ تجاوز (آیه ۷۷: «سِيءَ بِهِم»).
(۲) \textbf{انقباض (ضیق + ذرع)}: آیه ۷۷: «وَضَاقَ بِهِمْ ذَرْعًا» --- تحت‌اللفظ: «دامنه‌ی بازوی او تنگ شد». «ذَرع» در اصل، مقدارِ گستره‌ی بازو است؛ استعاره از کنشگری است که \textbf{توانِ عملیاتیِ} خود را برای پاسخ‌دهی، تنگ یافته. تصویر، دقیقاً تصویرِ کنشگری است که نمی‌تواند خشمِ خود را \textbf{عمل} کند.
(۳) \textbf{بیانِ زبانیِ آرزومندانه}: آیه‌ی ۸۰: «لَوْ أَنَّ لِي بِكُمْ قُوَّةً أَوْ آوِي إِلَىٰ رُكْنٍ شَدِيدٍ» --- هرگاه عملِ مستقیم بسته شده باشد، خشم به‌صورتِ آرزویی برای قدرت یا یارگیریِ نهادی بیان می‌شود؛ «قُوَّة» نیروی شخصی است، و «رُکنٍ شَدید» یاوریِ نهادی.
(۴) \textbf{انتقالِ کنشگری}: ضیقِ بازداشته‌ی لوط، توسطِ کنشگریِ الهی برداشته می‌شود (فرشتگان می‌آیند، شهر ویران می‌شود). روایتِ قرآنی بدین‌سان \textbf{پرخاشگریِ بازداشته} را الهیاتی می‌کند: آن‌گاه که کنشگرِ مظلوم قدرتی انسانی ندارد، کنشگریِ الهی مسیرِ خشم را به انجام می‌رساند.

این پاره از قرآن، خوانشِ ضیق را به‌مثابهِ صرفاً انقباضِ پیشاخشمی منع می‌کند: ضاق ذرعاً، \textbf{پرخاشگریِ پسا-خشمیِ بازداشته} را عملیاتی می‌کند --- خشمی که به آستانه‌ی کنش رسیده ولی به‌صورتِ ساختاری از انجام بازداشته شده است. واژه‌ی ضیق در قرآن، بنابراین، از یک برچسبِ مرحله‌ایِ ساده، غنی‌تر است.

\textbf{۴-۴-۲. حُزن (۴۲ بسامد، ۲۶ مکی / ۱۶ مدنی).} پربسامدترین واژه‌ی مرحله ۲، حزن است که ریشه‌شناسیِ آن نکته‌ای ظریف را آشکار می‌سازد. ابن منظور در \textbf{لسان العرب} یادآور می‌شود که این ریشه در اصل به «ناهمواری و خشونتِ زمین» اشاره دارد و به‌صورتِ استعاری به «حالتِ سنگینی و درشتیِ روانی» منتقل شده است. این انتقالِ استعاری، خود یک شاهدِ موضعی از نظریه‌ی استعاره مفهومی است: در ادراکِ زبانیِ عربی، اندوه، نوعی «ناهمواریِ روانی» است که مسیرِ نرم و آرامِ زندگی را با موانعی دشوار مواجه می‌کند. این استعاره با مدلِ غربیِ EMOTION IS A LANDSCAPE (که توسط کوچش در ۲۰۰۰ صورت‌بندی شد) به‌نحوِ شایانِ توجهی همخوان است.

آیه ۳۸ سوره بقره، ساختاری فشرده ولی مهم را عرضه می‌کند: «فَمَنْ تَبِعَ هُدَايَ فَلَا خَوْفٌ عَلَيْهِمْ وَلَا هُمْ يَحْزَنُونَ» (پس هر که از هدایتم پیروی کند، نه بیمی بر آنان است و نه اندوهگین می‌شوند). در این آیه، حزن در کنار «خوف» قرار می‌گیرد و در یک تقابلِ صریح با «هدایت» می‌نشیند. این هم‌نشینی، حزن را به‌مثابه یک حالتِ «دور افتادن از هدایت» معرفی می‌کند، نه یک واکنشِ صرفاً موقعیتی. علامه طباطبایی در \textbf{المیزان} (ج۱، ص۱۴۶) این آیه را «اساسی برای علم النفس قرآنی» می‌نامد و می‌نویسد که نفیِ خوف و حزن در آن، «نه نفیِ تجربه عاطفی، بلکه نفیِ ریشه دواندنِ این عواطف در عمق وجود است».

ریخت‌شناسیِ حزن نیز نکته‌ای جالب را آشکار می‌سازد: نسبتِ فعل به اسم در آن ۳۷ به ۵ است (۸۸٪ فعلی)، که چنان‌که در شکل ۵ مشاهده می‌شود، شدیدترین گرایش به فعلی‌سازی در میان چهارده ریشه است. این یافته را می‌توان به‌مثابه نشانه‌ای از این تفسیر کرد که حزن، در زبان قرآنی، نه یک ویژگیِ شخصیتی پایدار (اسم: حُزن)، بلکه یک عمل/حالتِ مستمر (فعل: یَحزَنُ، حَزِنَ، تَحزَنُ) است که می‌تواند تنظیم شود.

\textbf{۴-۴-۲-۱. ملاحظه: حزن، تنها حالتِ پسا-خشمی نیست.} یک ملاحظه‌ی مفسرانه باید در این‌جا گزارش شود. حزن \textbf{انحصاراً} از خشم برنمی‌خیزد؛ این واژه در دستگاهِ قرآنی، دست‌کم چهار ارزشِ معنایی دارد:

(الف) \textbf{اندوهِ مستقل بر فقدان}: آیه‌ی ۸۴ سوره‌ی یوسف («وَتَوَلَّىٰ عَنْهُمْ وَقَالَ يَا أَسَفَىٰ عَلَىٰ يُوسُفَ وَابْيَضَّتْ عَيْنَاهُ مِنَ الْحُزْنِ»)، اندوهِ ممتدِ یعقوب را بر فقدانِ یوسف به تصویر می‌کشد --- اندوهِ بر فقدانِ محبوب، بدون آمیزشی با خشم.
(ب) \textbf{اندوهِ ناشی از ستمِ بیرونی} (اندوهِ پیامبران بر تکذیبِ قوم).
(ج) \textbf{پشیمانیِ اخلاقیِ مثبت} --- اندوه بر گناه، که ستوده‌اشان است (تَوَجُّهِ توبه‌ای).
(د) \textbf{خشمِ به‌درون‌برگشته} آن‌گاه که بیرونی‌سازیِ آن از طریقِ ساختاری بسته شده --- نزدیک‌تر به جایگاهی که در مرحله ۲ بدان داده‌ایم.

این واژه را بر اساسِ تراکمِ پیکره‌ای و نزدیکیِ معنایی (ب) و (د) به میدانِ خشم در پیوستار می‌گنجانیم؛ ولی \textbf{ارزشِ معنایی} آن، بیش از آن‌چه یک برچسبِ مرحله‌ایِ تک‌بعدی اجازه می‌دهد، فراگذرنده است.

\textbf{۴-۴-۳. أَسَف (۵ بسامد، ۵ مکی / ۰ مدنی).} کم‌بسامدترین ولی نظریه‌نمای‌ترین واژه‌ی مرحله ۲، أَسَف است. ابن فارس در \textbf{مقاییس} و راغب در \textbf{المفردات} هردو، این ریشه را به «اندوهِ شدید آمیخته با خشمِ درونی» (الجمع بین الحزن و الغضب) تعبیر می‌کنند. این تعریف، أَسَف را در نقطه‌ای منحصر به‌فرد در پیوستار قرار می‌دهد: نقطه‌ی گذار میان فشارِ درونی و خشمِ آشکار.

آیه ۵۵ سوره زخرف، این ساختارِ گذاری را به‌نحوی روشن نشان می‌دهد: «فَلَمَّا آسَفُونَا انْتَقَمْنَا مِنْهُمْ فَأَغْرَقْنَاهُمْ أَجْمَعِينَ» (و چون ما را به أسَف آوردند، از ایشان انتقام گرفتیم و همگی را غرق کردیم). در این آیه، اسف به‌مثابه شرطِ زمینه‌ای انتقام تصویر می‌شود؛ یعنی اسف، نه یک هیجانِ منفعلِ پایان‌یافته، بلکه یک حالتِ گذاری است که مستقیماً به کنش (انتقام) منتهی می‌شود. شیخ طبرسی در \textbf{مجمع البیان} (ج۹، ص۸۲) این آیه را با ظرافت تحلیل کرده و نوشته است: «الأسف ههنا الغضب الشدید الذی ینشأ من الحزن السابق» (اسف در این‌جا، خشمِ شدیدی است که از اندوهِ پیشین برمی‌خیزد).

\textbf{یافته‌ی توزیعی استثنایی} درباره أسف آن است که هر ۵ بسامدِ آن، \textbf{به‌طور انحصاری در سور مکی} ظاهر می‌شود (آزمون دوجمله‌ای: \lr{p = 0.16}، که با n کم، ضعیف ولی همخوان است). تحلیلِ بافتیِ این پنج آیه نشان می‌دهد که اسف در همه موارد در روایت‌های مرتبط با مواجهه‌ی پیامبرانِ پیشین یا حضرت موسی (ع) با تکذیب‌گرانِ خود به‌کار رفته است.

\textbf{۴-۴-۳-۱. غَضبانَ أَسِفاً --- موسی (ع)، آیه‌های ۱۵۰ سوره‌ی اعراف و ۸۶ سوره‌ی طه --- خشمِ مقدس.} ترکیبِ توصیفیِ بسیار مهمِ غَضبانَ أَسِفاً، در آیه‌ی ۱۵۰ اعراف و آیه‌ی ۸۶ طه، در توصیفِ بازگشتِ موسی (ع) به نزدِ قومِ گوساله‌پرست به‌کار رفته. این ترکیب، یک دیالکتیکِ سه‌لایه را تولید می‌کند که خوانشِ خشم را به‌مثابهِ صرفاً امرِ منفی، دشوار می‌سازد:

(الف) \textbf{لایه‌بندیِ ساختاریِ درون/بیرون}: غضب، انرژیِ \textbf{بیرونی-جنبشیِ} پاسخِ موسی است --- مولفه‌ی حرکتیِ مرئی (انداختنِ الواح، گرفتنِ ریشِ هارون). أسف، انرژیِ \textbf{درونی-بالقوه} است --- اندوهی شفقت‌آمیز بر قومِ گم‌گشته، بر چهل روزِ هدررفته، بر بازگشتِ به جاهلیت.
(ب) \textbf{زمینه‌سازیِ علّی}: أسف \textbf{زمینه‌ی} غضب در این روایت است. اگر أسف نبود، خشمِ موسی صرفاً خشونت بود؛ ولی با حضورِ أسف، خشمِ او «خشمِ مقدس» می‌شود --- خشمی در خدمتِ هدایت.
(ج) \textbf{جهت‌گیریِ نیتی}: غضب \textbf{به‌بیرون} جهت‌گیری شده (به‌سویِ قوم، گوساله، برادر). أسف \textbf{به‌درون} جهت‌گیری شده (سوگواری بر فرصتِ هدررفته، بر بازگشتِ به نادانی).

این ترکیب، خوانشِ خشم را به‌مثابهِ صرفاً پدیده‌ای منفی، مشکل‌دار می‌کند. در موردِ پیامبران، خشم آن‌گاه که در شفقتِ حزنی بنیاد یافته باشد، \textbf{تولیدِ اخلاقی} دارد. دستگاهِ واژگانیِ قرآن، با جفت کردنِ غضب و أسف در یک پیشای واحد، این ساختارِ اخلاقی را با دقتی شایانِ توجه رمزگذاری می‌کند.

\textbf{۴-۴-۴. شواهدِ شبکه‌ای برای انسجامِ مرحله ۲.} تحلیلِ شبکه‌ی هم‌رخدادگیریِ آیه‌ای (شکل ۴) نشان می‌دهد که قوی‌ترین یالِ شبکه، میان «ضیق» و «حزن» است (وزن = ۳)، یعنی این دو واژه در ۳ آیه به‌طور هم‌زمان به‌کار رفته‌اند. این یافته، پشتوانه‌ی پیکره‌ایِ مستقلی برای گروه‌بندیِ این دو واژه در یک مرحله مشترک فراهم می‌آورد. تحلیلِ مرکزیتِ شبکه نیز نشان می‌دهد که حُزن با مرکزیتِ درجه‌ایِ وزن‌دارِ ۵ (هم‌بالاترین در شبکه با غضب)، گرهِ کانونیِ این مرحله را تشکیل می‌دهد.

\subsubsection{۴-۵. مرحله ۳ --- نفرتِ ارزیابانه}\label{ux645ux631ux62dux644ux647-ux6f3-ux646ux641ux631ux62a-ux627ux631ux632ux6ccux627ux628ux627ux646ux647}

مرحله‌ی سوم، با سه ریشه (نقم، سخط، مقت) و در مجموع ۲۷ بسامد، نفرتِ ارزیابانه را به‌مثابهِ پاسخی اخلاقی-داورانه به یک موضوع، صورت‌بندی می‌کند.

\textbf{۴-۵-۱. نَقم (۱۷ بسامد، ۱۲ مکی / ۵ مدنی).} نَقم، واژه‌ی «نکوهشِ همراه با اراده‌ی انتقام» است: خشمی که در ارزیابیِ اخلاقی بنیاد یافته و به‌سویِ کیفر جهت‌گیری شده. آیه ۱۲۶ سوره‌ی اعراف («وَمَا نَقَمُوا مِنَّا إِلَّا أَنْ آمَنَّا بِآيَاتِ رَبِّنَا») نمونه‌ی کانونی است: خشم \textbf{بر ایمانِ دیگران} --- ساحرانِ فرعون، پس از پذیرشِ پیامِ موسی (ع)، توصیف می‌شوند که موردِ نَقمِ فرعونی قرار گرفته‌اند \textbf{زیرا} ایمان آورده‌اند. ساختِ دستوریِ «نَقَمَ مِن + (سبب)» یک قالبِ کانونیِ قرآنی برای خشمِ ارزیابی‌بنیاد است.

این واژه، صفتِ الهیِ «المنتقم» (آیه ۲۲ سوره سجده و موازی‌های آن) را نیز تولید می‌کند، که پلِ مستقیمِ الهیاتی-واژگانی از ساحتِ هیجانیِ انسان به ساحتِ کیفریِ الهی را برقرار می‌سازد. جایگاه در پیوستار: نَقم میانِ سخط و غضب است؛ کنش‌مندتر از سخط (که داوریِ پایدار است) ولی کم‌جنبش‌تر از غضب (که کنشگرِ حرکتی را بسیج می‌کند). گرایشِ مکی (۱۲/۵) با حضورِ مکرر در روایت‌های پیامبرانه (ساحرانِ فرعون، قومِ لوط، اهلِ مَدیَن) سازگار است.

\textbf{۴-۵-۲. سَخَط (۴ بسامد، ۰ مکی / ۴ مدنی).} اگرچه سخط یکی از کم‌بسامدترین واژه‌های پیوستار است، الگوی توزیعیِ آن، یکی از مهم‌ترین یافته‌های پژوهش حاضر را در بر دارد. هر چهار بسامدِ این ریشه، \textbf{به‌طور انحصاری در سور مدنی} ظاهر می‌شود (آزمون دوجمله‌ایِ دوسویه با خطِ پایه‌ی مستخرج از پیکره (۰٫۷۴ مکی): \lr{p = 0.005}، معنادار در سطح \lr{α = 0.05}).

تحلیلِ بافتیِ این چهار آیه نشان می‌دهد که در همه موارد، سخط در پیوند با مفهومِ منافقان یا شبه‌مؤمنانِ نیمه‌پیمان به‌کار رفته است. آیه شاخصِ این کاربرد در سوره محمد (آیه ۲۸) ظاهر می‌شود: «ذَٰلِكَ بِأَنَّهُمُ اتَّبَعُوا مَا أَسْخَطَ اللَّهَ وَكَرِهُوا رِضْوَانَهُ فَأَحْبَطَ أَعْمَالَهُمْ» (این به سببِ آن است که از آنچه خدا را به سخط آورده، پیروی کرده و رضایتش را خوش نداشته‌اند؛ پس {[}خداوند{]} اعمالشان را تباه ساخت). در این آیه، چنان‌که نَریمانی، اقبالی و چاری (۱۴۰۰) نیز در تحلیلِ تفسیریِ خود به‌درستی نشان می‌دهند، سخط نه به‌مثابه یک هیجانِ آنی، بلکه به‌مثابه یک داوریِ ارزشیِ مستمر تصویر می‌شود.

پژوهش حاضر، با تکیه بر یافته‌ی توزیعیِ پیکره‌ای، گامی فراتر برمی‌دارد و فرضیه زیر را پیشنهاد می‌کند:

\begin{quote}
\textbf{فرضیه ۳ (تخصص‌یافتگیِ گفتمانیِ سَخَط):} سخط در گفتمانِ قرآنی، یک واژه‌ی گفتمانی-تخصصی است که برای توصیفِ نارضایتیِ الهی نسبت به منافقان و شبه‌مؤمنان (که در فضای مدنی و در مواجهه با ساختارِ نوبنیادِ امّت ظاهر می‌شوند) رزرو شده است. این تخصص‌یافتگی، نمونه‌ای از پدیده‌ی «تخصیصِ گفتمانیِ-بافتیِ یک لغتِ هم‌میدان» در زبانِ قرآن است.
\end{quote}

\textbf{۴-۵-۳. مَقت (۶ بسامد، ۴ مکی / ۲ مدنی).} مَقت، واژه‌ی «نفرتِ اخلاقیِ شدید» است. در چهار لغت‌نامه‌ی معیار، تعریفِ مشترکی غالب است: «أَشَدُّ الإِبغاضِ» (شدیدترین گونه‌ی إبغاض)، چنان‌که زجّاج و اللیث (در تهذیب اللغة) ضبط کرده‌اند. راغب توضیحِ تکمیلی می‌دهد: مَقت \textbf{از مشاهده‌ی فعلِ قبیح برمی‌خیزد} («عن أمرٍ قبیحٍ رَکِبَه») --- یعنی مَقت یک پاسخِ \textbf{صرفاً ارزیابانه} است.

این بنیاد، مقت را از غضب در دو جنبه‌ی کلیدی متمایز می‌کند:

(الف) \textbf{علیت}: مقت صرفاً ارزیابانه است (با مشاهده‌ی فعلِ قبیح برمی‌خیزد)، در حالی که غضب می‌تواند با آسیبِ شخصی یا ناکامی، علاوه بر ارزیابیِ اخلاقی، آغاز گردد.
(ب) \textbf{پروفایلِ کنش}: پاسخِ مقت \textbf{روی‌گردانی-فاصله‌گیری} است (پشت‌کردن، گسستِ تأیید) و نه \textbf{رویاروییِ-هجومی} (بسیجِ کنش برای کیفر). مقت می‌گریزد؛ غضب پیش می‌رود.

آیات کانونی، مقت را به‌مثابهِ یک نقطه‌ی پایه‌ی نفرتِ ارزیابانه تثبیت می‌کنند:

\begin{itemize}
\tightlist
\item
  آیه ۱۰ سوره‌ی غافر («لَمَقْتُ اللَّهِ أَكْبَرُ مِن مَّقْتِكُمْ أَنفُسَكُمْ»): یک \textbf{فراخوانِ دو‌برابره} که «مَقت النفس» (نفرتِ از خود در روزِ قیامت) را به‌مثابهِ یک طبقه‌ی روان‌اخلاقیِ قرآنی تثبیت می‌کند.
\item
  آیه ۲۲ سوره‌ی نساء (در بافتِ ازدواج‌های ممنوع) چنین ازدواج‌هایی را «فَاحِشَةً وَمَقْتًا وَسَاءَ سَبِيلًا» توصیف می‌کند --- نفرتِ اخلاقیِ \textbf{نهادینه} به‌مثابهِ ویژگیِ یک طبقه از کنش‌ها.
\item
  آیه ۳۹ سوره‌ی فاطر («وَلَا يَزِيدُ الْكَافِرِينَ كُفْرُهُمْ عِندَ رَبِّهِمْ إِلَّا مَقْتًا») مقتِ الهی را تابعِ افزایشِ کفرِ کافران معرفی می‌کند.
\item
  آیه ۳ سوره‌ی صف («كَبُرَ مَقْتًا عِندَ اللَّهِ أَن تَقُولُوا مَا لَا تَفْعَلُونَ») گفتارِ بدونِ کردار را آماجِ مقتِ شدیدِ الهی می‌داند.
\end{itemize}

در همه‌ی این آیات، این واژه به‌مثابهِ اصطلاحی برای \textbf{نفرتِ اخلاقی به‌مثابهِ پاسخِ ارزیابانه} تثبیت می‌گردد، که از غضبِ جنبشی-حرکتیِ مرحله ۴ متمایز است.

\subsubsection{۴-۶. مرحله ۴ --- خشمِ فعال}\label{ux645ux631ux62dux644ux647-ux6f4-ux62eux634ux645-ux641ux639ux627ux644}

مرحله‌ی چهارم، با دو ریشه (غضب، حرد) و در مجموع ۲۵ بسامد، خشمِ فعالِ بسیج‌کننده‌ی کنشگرِ حرکتی را در بر می‌گیرد.

\textbf{۴-۶-۱. غَضَب (۲۴ بسامد، ۱۳ مکی / ۱۱ مدنی).} غضب، مرکزی‌ترین و پربسامدترین واژه‌ی این مرحله است و توزیعِ نسبتاً متوازنِ مکی-مدنیِ آن (۵۴٪ مکی) آن را به‌مثابه واژه‌ی «بدون نشانگر» \lr{(unmarked)} در میدانِ خشمِ فعال معرفی می‌کند. در چهار لغت‌نامه‌ی معیار، غضب به «هیجانی شدید ناشی از احساسِ ظلم یا تعدی، که فرد را به واکنش وامی‌دارد» (راغب) یا به «ثَوَران دَم القَلب طلباً للانتقام» (مصطفوی) تعبیر شده است. این تعریف، غضب را به‌مثابه یک هیجانِ \textbf{هدف‌مند} و \textbf{معطوف به دیگری} صورت‌بندی می‌کند.

تحلیلِ ریخت‌شناختیِ غضب، نکته بسیار شایان توجهی را آشکار می‌سازد: نسبتِ اسم به فعل در آن ۱۶ به ۶ است (۶۷٪ اسمی)، که برخلافِ حُزن (۸۸٪ فعلی) است. این تفاوتِ ریخت‌شناختی، یک یافته‌ی نظریِ مهم را پشتیبانی می‌کند: حُزن یک «حالتِ گذرنده» \lr{(transient state)} و غضب یک «حالتِ-مفهومِ ساختاری» \lr{(structural concept-state)} است.

آیه‌ی شاخصِ کاربرد غضب آیه ۶ سوره فتح است: «غَضِبَ اللَّهُ عَلَيْهِمْ وَلَعَنَهُمْ» (خداوند بر آنان غضب کرد و آنان را لعنت نمود). فخر رازی در \textbf{مفاتيح الغيب} (ج۲۸، ص۸۴) این ویژگیِ پیوندیِ غضب را در یک عبارتِ نظریِ عمیق چنین خلاصه می‌کند: «الغضبُ من الله إرادَةُ إنزال العقاب» (غضبِ الهی، اراده‌ی نزول عقوبت است) --- یعنی غضب در ساحتِ الهی، نه یک تجربه‌ی عاطفی، بلکه نوعی \textbf{داوریِ نظام‌مند} است که به اقدامِ اخلاقی منتهی می‌شود.

\textbf{۴-۶-۲. تحلیلِ شبکه‌ای: غضب به‌مثابه‌ی گرهِ پیوندی.} تحلیلِ شبکه (شکل ۴ و ۷) نشان می‌دهد که غضب با مرکزیتِ بینابینیِ ۰٫۱۷، یکی از سه گرهِ پلیِ شبکه است. به‌طور خاص، یال‌های قابلِ توجهِ آن عبارت‌اند از: غضب-غیظ (که مسیرِ گذار به مرحله ۵ را پیکره‌سازی می‌کند)، غضب-حُزن (که با مرحله ۲ پیوند دارد)، و غضب-بَغی (که با مرحله ۶ پیوند می‌خورد).

\textbf{۴-۶-۳. حَرد (۱ بسامد: آیه ۲۵ سوره‌ی قلم --- هاپَکس قرآنی).} حَرد، یک هاپَکس قرآنی است که تنها در آیه‌ی ۲۵ سوره‌ی قلم («وَغَدَوْا عَلَىٰ حَرْدٍ قَادِرِينَ») در داستانِ اصحاب الجنّه (صاحبانِ باغ) ظاهر می‌شود. صاحبانِ باغ، که شب پیش سوگند خورده بودند که پیش از طلوع دروی کنند تا حقِ مساکین را \textbf{منع} کنند، «در پگاه، بر حَرد، در حالی که خود را قادر می‌پنداشتند، به سوی باغ شتافتند». تفسیرِ سنتی، یک تفسیرِ دو‌گانه‌ی شایانِ توجه ارائه می‌دهد:

\begin{itemize}
\tightlist
\item
  مجمع البیانِ طبرسی، حَرد را به هر دو معنای \textbf{غضب} و \textbf{قصد} (نیتِ هدف‌مند) معنا می‌کند. این واژه، خشمِ همراه با قصد است.
\item
  المیزانِ علامه طباطبائی، تفسیر را دقیق‌تر می‌کند: حَرد یعنی \textbf{خشم به‌علاوه‌ی منعِ آگاهانه}؛ یعنی خشمی هدف‌مند که از طریقِ خِسّت، مسلح شده.
\end{itemize}

پدیدارشناسیِ حَرد، بنابراین، در درونِ مرحله ۴ متمایز است. غضب، کنشگرِ حرکتی را به‌سویِ هدفی بیرونی بسیج می‌کند؛ حَرد، کنشگرِ حرکتی را به‌سویِ \textbf{منع} آگاهانه و خِسّت‌آمیز هدایت می‌کند --- خشمی که جهان را از آن‌چه به آن طلب می‌شود، عاری می‌کند. امضای جامعه‌شناختیِ حَرد، خشمی است که از طریقِ خِسّتِ نهادی یا قراردادی، مسلح می‌شود. پیامدِ روایتیِ آن در سوره‌ی قلم، نابودیِ باغ توسطِ توفانِ شب است: کیفرِ الهی بر طرحِ مبتنی بر حَرد.

جایگاه حَرد در پیوستار: \textbf{خشمِ فعال} (مرحله ۴)، ولی به‌نحو متمایزی \textbf{هدف‌مند-و-مانع‌انگیز} و نه صرفاً رویاروی-هجومی. وضعیتِ هاپَکس بودنِ این واژه --- جفت با سنتِ غنیِ تفسیری --- نمونه‌ای روشن از اقتصادِ واژگانیِ قرآن است: یک بسامدِ یگانه به‌مثابهِ لنگرگاهِ زبانیِ یک طبقه‌ی روان‌اخلاقیِ کامل، عمل می‌کند.

\subsubsection{۴-۷. مرحله ۵ --- خشمِ متراکم/انفجاری}\label{ux645ux631ux62dux644ux647-ux6f5-ux62eux634ux645-ux645ux62aux631ux627ux6a9ux645ux627ux646ux641ux62cux627ux631ux6cc}

\textbf{۴-۷-۱. غَیظ (۱۱ بسامد، ۳ مکی / ۸ مدنی).} غیظ، در چهار لغت‌نامه‌ی معیار، به «خشمِ متراکم، انباشته و آماده انفجار» تعبیر می‌شود؛ یعنی خشمی که هنوز کاملاً به بیرون نریخته است ولی فشارِ درونیِ آن چنان است که اگر مهار نشود، به فورانِ رفتاری بدل خواهد شد. این تعریف، تمایزِ ساختاریِ غیظ از غضب را روشن می‌سازد: غضب، خشمِ آشکار و معطوف به هدف است؛ غیظ، خشمِ فروبرده در درون و در آستانه‌ی انفجار است.

این تمایز در آیه ۱۳۴ سوره آل عمران به‌نحوی برجسته نمایان می‌شود: «الَّذِينَ يُنْفِقُونَ فِي السَّرَّاءِ وَالضَّرَّاءِ وَالْكَاظِمِينَ الْغَيْظَ وَالْعَافِينَ عَنِ النَّاسِ ۗ وَاللَّهُ يُحِبُّ الْمُحْسِنِينَ» ({[}آنان که{]} در فراخی و تنگدستی انفاق می‌کنند، خشم خود را فرو می‌خورند و از مردم درمی‌گذرند، و خدا نیکوکاران را دوست دارد).

از منظرِ نظریه استعاره مفهومی، فعلِ «کاظِم» از ریشه «ک-ظ-م» در لغت به «بستنِ دهانه ظرف (مَشک)» اشاره دارد. این فعل، در ترکیب با غیظ، استعاره‌ی مفهومیِ «خشم به‌مثابه‌ی محتوای فشرده درون ظرف» را به‌نحو صریح عملیاتی می‌کند. مؤمنان، در این تصویر، کسانی هستند که با کنترلِ فعالِ این فشار درونی، از فروپاشیِ ظرف و سرریز هیجان به ساحتِ رفتاری جلوگیری می‌کنند. این تصویر دقیقاً همان «خشم گرمای مایعی در یک ظرف فشرده» \lr{(ANGER IS THE HEAT OF A FLUID IN A PRESSURIZED CONTAINER)} است که لیکاف و کوچش (۱۹۸۷) چهارده قرن بعد در زبان انگلیسی شناسایی کردند.

\textbf{۴-۷-۲. تَمَیُّز و فروپاشیِ ظرف (آیه ۸ سوره ملک --- تجلی، نه ریشه‌ی مستقل).} شدیدترین تصویر قرآنی از خشم در آیه ۸ سوره ملک یافت می‌شود: «تَكَادُ تَمَيَّزُ مِنَ الْغَيْظِ ۖ كُلَّمَا أُلْقِيَ فِيهَا فَوْجٌ سَأَلَهُمْ خَزَنَتُهَا أَلَمْ يَأْتِكُمْ نَذِيرٌ» ({[}آتش جهنم{]} چنان است که گویی می‌خواهد از شدتِ غیظ، از هم بگسلد\ldots). فعلِ «تَمَيَّز» از ریشه «م-ي-ز» در اصل به «جدا شدن، از هم گسستن» اشاره دارد؛ در این بافت، شدتِ غیظ به حدی است که گویی ساختارِ وجودیِ خودِ ظرفِ هیجان (در این‌جا، آتش جهنم) رو به فروپاشی می‌گذارد. در تحلیلِ حاضر، تمیّز را به‌مثابهِ \textbf{تجلیِ} مرحله ۵ (نشانه‌ی رفتاریِ شکستِ مهار) می‌دانیم، نه به‌مثابهِ ریشه‌ای مستقل؛ بدین‌سان شمارِ ریشه‌های کانونی، چهارده می‌ماند.

این تصویرِ استعاری، در ادبیاتِ بین‌المللیِ زبان‌شناسیِ شناختی، تا کنون به‌قدرِ کافی برجسته نشده است. ادعای محتاطانه‌ی پژوهش حاضر آن است که آیه‌ی ۶۷:۸، یک نمونه‌ی \textbf{بسیار شفاف} از پدیده‌ی «فروپاشیِ ظرف» \lr{(container collapse)} را در یک متنِ پیشامدرن فراهم می‌آورد. ادعاهای قوی‌تر درباره‌ی جایگاهِ تطبیقیِ این آیه (نظیر «قدیمی‌ترین» یا «کامل‌ترین») نیازمندِ پیمایشِ نظام‌مندِ تصاویرِ ظرفی-انفجاریِ خشم در ادبیاتِ پیشامدرن است که هنوز انجام نشده. لیکاف و کوچش (۱۹۸۷) و کوچش (۲۰۰۰) این پدیده را به‌مثابه «گذار از فشار به انفجار» در زبان انگلیسی شناسایی کردند، ولی هیچ‌یک از مثال‌هایِ ایشان به وضوحِ تصویرِ قرآنی نزدیک نیست؛ زیرا در زبان انگلیسیِ روزمره، عبارتی نظیر «\lr{to explode with rage}» انفجار را گزارش می‌کند، ولی فروپاشیِ ساختاریِ خودِ ظرف را به‌نحو صریح بازنمایی نمی‌کند. در حالی که فعلِ «تَمَيَّز» در آیه ۶۷:۸، \textbf{ساختارِ خود ظرف را در معرض گسستگی} قرار می‌دهد.

شکل ۶ این نگاشتِ سه‌مرحله‌ای را به‌نحوِ تصویری بازنمایی می‌کند: ظرفِ باز (مرحله ۲: غضب) → ظرفِ مهر شده با فشارِ بالا (مرحله ۳الف: غیظ + کظم) → ظرفِ در آستانه گسستن (مرحله ۳ب: تمیّز).

\textbf{۴-۷-۳. توزیعِ مدنی-غالبِ غیظ.} نسبتِ ۳ مکی به ۸ مدنی برای غیظ، در نگاهِ نخست با الگوی مکی-غالبِ بسیاری از واژگانِ هیجانی متفاوت است. تبیین این الگو در ساحتِ بافتارِ پیامبریِ مدنی است: مهار خشم (کظم غیظ)، یک مهارتِ تنظیم هیجان است که در فضای پرتنشِ امّتِ نوبنیادِ مدینه، اهمیتِ ویژه‌ای می‌یابد.

\subsubsection{۴-۸. مرحله ۶ --- پیامدهای رفتاری (همراه با ملاحظات)}\label{ux645ux631ux62dux644ux647-ux6f6-ux67eux6ccux627ux645ux62fux647ux627ux6cc-ux631ux641ux62aux627ux631ux6cc-ux647ux645ux631ux627ux647-ux628ux627-ux645ux644ux627ux62dux638ux627ux62a}

مرحله‌ی ششم، که سه واژه «بَغی» (۹۶ بسامد)، «طُغیان» (۳۹ بسامد) و «عُتُوّ» (۱۰ بسامد) را در بر می‌گیرد، در بسامدِ خود (۱۴۵) با مجموعِ مراحلِ ۱ تا ۵ (۱۶۷) قابلِ مقایسه است (\lr{χ²(1) = 1.55}، رد نمی‌شود). تفسیرِ این نزدیکی در §۴-۸-۵ به‌عنوان \textbf{علیتِ دوسویه} بسط می‌یابد.

\textbf{۴-۸-۱. بَغی (۹۶ بسامد، ۵۵ مکی / ۴۱ مدنی).} پربسامدترین واژه‌ی پیوستار و در عین حال، گرهِ کانونیِ شبکه است. در منابع لغوی، بَغی به «طلبِ تجاوزکارانه و فراتر رفتن از حقِ دیگران» (راغب) تعبیر می‌شود. آیه ۴۲ سوره شوری مصداقِ شاخصِ این کاربرد است: «إِنَّمَا السَّبِيلُ عَلَى الَّذِينَ يَظْلِمُونَ النَّاسَ وَيَبْغُونَ فِي الْأَرْضِ بِغَيْرِ الْحَقِّ ۚ أُولَٰئِكَ لَهُمْ عَذَابٌ أَلِيمٌ» (راهِ {[}اعتراض{]} تنها بر کسانی است که بر مردم ستم می‌کنند و در زمین به ناحق سرکشی می‌نمایند؛ آنانند که عذابی دردناک خواهند داشت).

ابن عاشور در \textbf{التحرير والتنوير} (ج۲۵، ص۱۲۶) به‌درستی اشاره می‌کند که بَغی در این آیه، در یک هم‌نشینیِ معنادار با ظلم می‌نشیند: «البغیُ والظلمُ مترادفانِ تَقریباً، إلّا أن البغیَ یَخصُّ التعدی غیرَ المستحَقّ علی الآخر، والظلمُ یعمّ کلّ الجور» (بغی و ظلم تقریباً مترادف‌اند، با این تفاوت که بغی به‌طورِ خاص به تعدیِ ناروا بر دیگری اشاره دارد، ولی ظلم همه‌ی انواعِ ستم را در بر می‌گیرد).

\textbf{شاهدِ شبکه‌ای کلیدی:} تحلیلِ مرکزیتِ شبکه (جدول ۳ و شکل ۷) نشان می‌دهد که بَغی با مرکزیتِ نزدیکیِ ۰٫۵۵ بالاترین مقدارِ این معیار را در شبکه دارد و در مرکزیتِ بینابینی (۰٫۱۵) با غضب و أسف هم‌رتبه است؛ مرکزیتِ ویژه‌بُردارِ آن (۰٫۵۰) نیز پس از کره (۰٫۵۰) دومین مقدار شبکه است. این الگوی شبکه‌ای، فرضیه‌ی ف.۴ پژوهشِ حاضر را تأیید می‌کند: \textbf{بَغی به‌مثابهِ گرهِ پلی، میان میدانِ هیجانی (غضب، غیظ) و میدانِ اخلاقی-اجتماعیِ ظلم/جرم/فسق/عداوت ایفای نقش می‌کند.} جدول ۴ نشان می‌دهد که بَغی با ۱۹ هم‌رخدادگیریِ آیه‌ای با واژگانِ چترِ مفهومی، بالاترین تعدادِ پل‌های بینامیدانی را در میانِ همه‌ی چهارده ریشه‌ی کانونی دارد.

\textbf{۴-۸-۲. طُغیان (۳۹ بسامد، ۲۹ مکی / ۱۰ مدنی).} ریشه «ط-غ-ي» در لغت به «سرریز شدنِ آب از حدّ خود» اشاره دارد؛ معنایی استعاری که به‌صورتِ گسترده برای توصیفِ سرکشیِ انسان در برابرِ حدودِ الهی به‌کار می‌رود. آیات ۶ و ۷ سوره علق، که از کوتاه‌ترین و کلیدی‌ترین آیاتِ گفتمانِ مکی محسوب می‌شوند، الگوی شناختیِ طغیان را به‌نحوی فشرده ترسیم می‌کنند:

\begin{quote}
«كَلَّا إِنَّ الْإِنسَانَ لَيَطْغَىٰ ٭ أَن رَّآهُ اسْتَغْنَىٰ»

(حقا که انسان طغیان می‌کند، آن‌گاه که خود را بی‌نیاز می‌بیند).
\end{quote}

این آیه، یک شناخت‌مدلِ روان‌شناختی بسیار پیشرفته را در فشرده‌ترین صورتِ ممکن بیان می‌کند: \textbf{خودبسندگیِ ادراکی} \lr{(perceived self-sufficiency)}، شرطِ شناختیِ طغیان است. به دیگر بیان، طغیان تنها یک پدیده‌ی هیجانی نیست، بلکه پاسخی شناختی به یک ادراکِ خاص از خود (ادراکِ بی‌نیازی) است. این فرمول‌بندی با نظریه ارزیابی شناختیِ لازاروس \lr{(Lazarus 1991)} و مدل فرآیندیِ شِرر \lr{(Scherer 2001)} به‌نحوِ شایانِ توجهی هم‌سو است: هر دو نظریه، هیجان را به‌مثابه پاسخی به ارزیابیِ شناختی از موقعیت تعریف می‌کنند، نه پاسخی خام و غریزی.

قاری‌زاده و همکاران (۱۴۰۲) در پژوهش خود به نقشِ این چارچوبِ شناختی در طغیان اشاره کرده‌اند، هرچند آن را در پیوستارِ کلیِ هیجان قرار نداده‌اند. یافته‌ی پیکره‌ایِ پژوهش حاضر، گامی فراتر می‌رود: ۲۹ بسامدِ مکیِ طغیان (در برابر ۱۰ مدنی) نشان می‌دهد که این مفهوم به‌صورتِ تخصصی در فضای گفتمانیِ مکی، که به مواجهه با شناختِ نادرستِ مشرکان (به‌مثابه‌ی ریشه شناختیِ طغیان) معطوف است، شکل گرفته است.

\textbf{۴-۸-۳. عُتُوّ (۱۰ بسامد، ۹ مکی / ۱ مدنی).} شدیدترین واژه‌ی میدانِ مورد مطالعه، عُتُوّ از ریشه «ع-ت-و» است که با تنها ۱۰ بسامد، نقاطِ پایانیِ پیوستار را علامت می‌گذارد. در لغت، عتوّ به «سرکشیِ سختِ همراه با تکبر و خودبزرگ‌بینی» (لسان العرب) تعبیر می‌شود؛ یعنی مرحله‌ای که در آن عصیان نه‌تنها به رفتار، بلکه به یک موضع‌گیریِ وجودیِ تثبیت‌شده بدل می‌شود.

آیه ۲۱ سوره فرقان، بافتِ کانونیِ این واژه است: «وَقَالَ الَّذِينَ لَا يَرْجُونَ لِقَاءَنَا لَوْلَا أُنزِلَ عَلَيْنَا الْمَلَائِكَةُ أَوْ نَرَىٰ رَبَّنَا ۗ لَقَدِ اسْتَكْبَرُوا فِي أَنفُسِهِمْ وَعَتَوْا عُتُوًّا كَبِيرًا» (و کسانی که به دیدارِ ما امیدی ندارند، گفتند: چرا فرشتگان بر ما فرود نیامدند یا پروردگارمان را نمی‌بینیم؟ هر آینه در {[}اندیشه{]} خود، خود را بزرگ شمردند و سرکشیِ بزرگی کردند). نکته معناشناختیِ شایانِ توجه آن است که عتوّ، در هم‌نشینیِ بافتیِ خود، پیوسته با مفاهیمِ «استکبار»، «شناختِ نادرستِ خود»، و «انکارِ آگاهانه‌ی حق» همراه می‌شود؛ یعنی این واژه، پایانیِ سرکشی نه به‌مثابه یک کنشِ لحظه‌ای، بلکه به‌مثابه یک موضعِ شناختی-وجودیِ تثبیت‌شده را توصیف می‌کند.

توزیعِ به‌شدت مکیِ این واژه (۹ مکی / ۱ مدنی، آزمون دوجمله‌ای: \lr{p = 0.059}) بافتارِ تاریخی-گفتمانیِ آن را روشن می‌سازد: عتوّ در گفتمانِ مکی، پاسخی به موضعِ سرسختانه‌ی مشرکانِ مکه است که در بسیاری از داستان‌های اقوامِ گذشته (نظیر قوم نوح، عاد، ثمود، فرعون) نیز به‌مثابه نقطه پایانیِ مسیرِ طغیان و شرط نهاییِ نزولِ عذاب الهی بازنمایی می‌شود.

\textbf{۴-۸-۵. ملاحظات بر ریشه‌های مرحله ۶ (بَغی، طُغیان، عُتُوّ): علیتِ دوسویه.} سه ملاحظه‌ی مفسرانه باید برای مرحله ۶ گزارش شود؛ ملاحظاتی که در قرائتِ بازنگری‌شده‌ی پیوستار، کلیدی هستند. \emph{بَغی} در کاربردِ قرآنی، بسیار اوقات از طمع، برتری‌جویی (عُلوّ فی الأرض) یا منفعت‌طلبیِ ستمگرانه برمی‌خیزد، نه از خشمِ آنی؛ کنشِ فراروی، گاه \textbf{محتوای} میلِ تباه است (مقایسه با آیه ۲۱ سوره‌ی انعام: «إِنَّهُ لَا يُفْلِحُ الظَّالِمُونَ») و نه \textbf{پیامدِ} یک هیجان. \emph{طُغیان} --- به‌شکلِ پارادایمیِ فرعونی (آیه ۱۷ سوره‌ی نازعات؛ آیه ۴ سوره‌ی قصص) --- در استکبارِ ساختاری و اعمالِ قدرتِ مهار‌نشده ریشه دارد، نه به‌مثابهِ امتدادِ پدیدارشناختیِ خشم؛ آیه‌ی ۵۵ سوره‌ی شعراء («إِنَّهُمْ لَنَا لَغَائِظُونَ») دقیقاً به این دلیل قابلِ توجه است که قطبیت را معکوس می‌کند: خشمِ ستمگر از سرکشیِ ستمدیده برمی‌خیزد، نه از خشمِ خود او. \emph{عُتُوّ} عمدتاً به استکبار وابسته است و به‌مثابهِ سرپیچیِ ثابت و تثبیت‌شده --- ویژگیِ پایدارِ شخصیتِ اخلاقی و نه حالتی هیجانیِ گذرا --- صورت‌بندی می‌گردد. ما این سه ریشه را در مرحله ۶ نگه می‌داریم، چرا که پیکره به‌درستی آن‌ها را در قطبِ رفتاریِ پیوستار می‌نشاند؛ ولی فلشِ علّیِ آن‌ها \textbf{دوسویه} است: خشم می‌تواند این رفتارها را تولید کند، ولی این رفتارها نیز می‌توانند خشم تولید کنند (و در روایتِ قرآنیِ ستمگران و پیامبران، اغلب چنین می‌کنند). نزدیکی میانِ مراحلِ ۱ تا ۵ (۱۶۷) و مرحله ۶ (۱۴۵)، با \lr{χ²(1) = 1.55}، دقیقاً همان چیزی است که در صورتِ صحتِ این قرائت انتظار می‌رود.

\begin{figure}
\centering
\includegraphics[width=0.95\linewidth,height=\textheight,keepaspectratio]{fig4_cooccurrence.pdf}
\caption{شبکه‌ی هم‌رخدادگیریِ آیه‌ایِ چهارده ریشه‌ی کانونی. هر یال، حضورِ مشترکِ دو ریشه در یک یا چند آیه را نمایان می‌سازد؛ عددِ روی یال، تعدادِ آیاتِ مشترک است. قوی‌ترین یال --- میان حزن و ضيق با وزنِ ۳ --- تأییدِ پیکره-درونیِ مستقیمی برای دسته‌بندیِ مرحله ۲ (فشارِ درونی و انقباض) فراهم می‌کند: این دو واژه‌ی پیشاخشمی، بیشتر از هر جفتِ دیگری در طیف، آیه‌ی مشترک دارند، که نشان می‌دهد قرآن آن‌ها را پدیده‌های گفتمانی-هم‌بسته تلقی می‌کند («ضیقُ الصدر در حزن»، هود: ۱۲؛ انعام: ۳۳). \emph{بغی} (مرحله ۶) خوشه‌ی سمتِ راست را لنگر می‌اندازد، با یال‌هایی که به سویِ واژگانِ ارزیابیِ اخلاقی (ظلم، جرم، فسق) امتداد می‌یابند. ساختارِ گذرای شبکه (نبودِ یالِ مستقیمِ مرحله ۱ → مرحله ۶؛ میانجی‌گریِ مراحلِ ۲ تا ۵)، هم‌سو با فرضیه‌ی شیبِ شدتِ کنش است، در عین این‌که با قرائتِ علیتِ دوسویه‌ی مرحله ۶ که در §۴-۸-۵ بسط یافت، نیز سازگار می‌ماند.}
\end{figure}

\begin{figure}
\centering
\includegraphics[width=0.95\linewidth,height=\textheight,keepaspectratio]{fig7_centrality.pdf}
\caption{سه سنجه‌ی مرکزیت محاسبه‌شده روی شبکه‌ی هم‌رخدادگیریِ آیه‌ای: (الف) درجه‌ی وزن‌دار --- مجموعِ پیوندهای هم‌رخدادگیری؛ (ب) بینابینی --- تواترِ قرارگرفتنِ گره روی کوتاه‌ترین مسیرها، شاخصِ توانِ گرهِ پلی؛ (ج) نزدیکی --- معکوسِ میانگینِ فاصله تا سایرِ گره‌ها. بَغی \textbf{بالاترین مرکزیتِ نزدیکی} را به دست می‌آورد (۰٫۵۵) و در مرکزیتِ بینابینی (۰٫۱۵) با غضب و أسف هم‌رتبه است --- تأییدی بر نقشِ آن به‌مثابهِ یکی از \emph{لولاهای اصلی} میانِ طیفِ هیجانیِ خشم (مراحلِ ۱-۵) و میدانِ گسترده‌ترِ ارزیابیِ اخلاقیِ قرآن (۳۱۵ بسامدِ ظلم، ۶۶ بسامدِ جرم، ۵۴ بسامدِ فسق، ۱۰۶ بسامدِ عداوت). پنلِ نزدیکی، الگوی مکمل را آشکار می‌کند: ضيق و حزن (مرحله ۲) به‌رغمِ بسامدِ خامِ پایین، نزدیکیِ بالا کسب می‌کنند، که نشان می‌دهد خوشه‌ی فشارِ درونی، از حیثِ ساختاری \emph{مرکزیِ} طیف است هرچند از حیثِ بسامدی پیرامونی است. این سنجه‌ها هم‌نوا، خوانشی لایه‌مند از شبکه را پشتیبانی می‌کنند: مرحله ۲ هسته‌ی عاطفی را تأمین می‌کند، مرحله ۶ پلِ ساختاری به میدان‌های همسایه را، و مرحله‌های ۳ تا ۵ انتقالِ انرژی میانِ آن‌ها را میانجیگری می‌کنند.}
\end{figure}

\subsubsection{۴-۹. آزمون پایداری: واژگانِ چترِ مفهومی}\label{ux622ux632ux645ux648ux646-ux67eux627ux6ccux62fux627ux631ux6cc-ux648ux627ux698ux6afux627ux646-ux686ux62aux631-ux645ux641ux647ux648ux645ux6cc}

برای راستی‌آزماییِ پایداریِ پیوستارِ پیشنهادی، پنج ریشه‌ی گسترده‌تر (شنأ، جرم، فسق، ظلم، عدو) نیز از پیکره استخراج شدند (یادآور: ریشه‌ی «کره» که پیش‌تر چترِ مفهومی بود، اکنون به‌مثابهِ ریشه‌ی مرحله ۱ در دستگاهِ خود پیوستار گزارش شده، نک. جدول ۲). توزیعِ این پنج ریشه (جدول ۵) نشان می‌دهد که:

\textbf{جدول ۵.} ریشه‌های چترِ مفهومی و توزیعِ آن‌ها

\begin{longtable}[]{@{}
  >{\centering\arraybackslash}p{(\linewidth - 10\tabcolsep) * \real{0.1667}}
  >{\centering\arraybackslash}p{(\linewidth - 10\tabcolsep) * \real{0.1667}}
  >{\centering\arraybackslash}p{(\linewidth - 10\tabcolsep) * \real{0.1667}}
  >{\centering\arraybackslash}p{(\linewidth - 10\tabcolsep) * \real{0.1667}}
  >{\centering\arraybackslash}p{(\linewidth - 10\tabcolsep) * \real{0.1667}}
  >{\raggedright\arraybackslash}p{(\linewidth - 10\tabcolsep) * \real{0.1667}}@{}}
\toprule\noalign{}
\begin{minipage}[b]{\linewidth}\centering
ریشه
\end{minipage} & \begin{minipage}[b]{\linewidth}\centering
کل
\end{minipage} & \begin{minipage}[b]{\linewidth}\centering
مکی
\end{minipage} & \begin{minipage}[b]{\linewidth}\centering
مدنی
\end{minipage} & \begin{minipage}[b]{\linewidth}\centering
p-value (دوجمله‌ای)
\end{minipage} & \begin{minipage}[b]{\linewidth}\raggedright
تفسیر
\end{minipage} \\
\midrule\noalign{}
\endhead
\bottomrule\noalign{}
\endlastfoot
شنأ & ۳ & ۱ & ۲ & -- & n خیلی کوچک؛ بدون نتیجه‌گیریِ آماری \\
جرم & ۶۶ & ۶۰ & ۶ & \lr{< 0.0001} & غلبه قاطع مکی؛ گفتمانِ نقدِ شرک و انکار \\
فسق & ۵۴ & ۲۰ & ۳۴ & ۰٫۰۰۰۷ & تمایلِ معنادار به مدنی؛ گفتمانِ انضباطِ امّتی \\
ظلم & ۳۱۵ & ۲۱۱ & ۱۰۴ & -- & فراگیرترین واژه میدانِ اخلاق هنجاری \\
عدو & ۱۰۶ & ۴۷ & ۵۹ & ۰٫۰۰۸ & تمایلِ مدنی؛ بافتارِ منافقان و اهلِ کتابِ پیمان‌شکن \\
\end{longtable}

این یافته‌ها، چند نکته نظری را روشن می‌سازند:

اولاً، \textbf{ظلم} با ۳۱۵ بسامد، فراگیرترین واژه میدانِ اخلاق هنجاریِ قرآن است و نمی‌تواند به‌مثابه یک گرهِ منفرد در پیوستارِ هیجانی عمل کند؛ این واژه، خود یک «واژه چترِ مفهومی» است که میدان‌های متعددی را پوشش می‌دهد.

دوماً، \textbf{جرم} و \textbf{فسق} الگوهای مکی-غالب و مدنی-غالبِ متضاد دارند: جرم در گفتمانِ مواجهه با شرک‌گرایانِ مکی، و فسق در گفتمانِ انضباطِ درون‌امّتیِ مدنی شکل گرفته است. این الگوی متضاد، خود نمونه‌ای از «تخصص‌یافتگیِ گفتمانیِ-بافتیِ یک واژه» است که در فرضیه ۳ پژوهش حاضر برای سَخَط نیز مشاهده شد.

سوماً، \textbf{عدو} با ۱۰۶ بسامد و توزیعِ ۴۷:۵۹، در میدانِ مجاورِ هیجانی-اخلاقیِ «دشمنیِ دوسویه» قرار دارد و از مفهومِ یک‌سویه‌ترِ «بَغی» (تعدیِ یک‌جانبه) متمایز است.

این یافته‌های جانبی، الگوی کلیِ پیوستارِ پیشنهادی را تأیید می‌کنند، ولی وضوحِ تحلیل، اقتضا می‌کند که در ساختارِ اصلیِ پیوستار، تنها چهارده ریشه‌ی کانونی گنجانده شوند.

\subsubsection{۴-۱۰. تحلیل تطبیقیِ ترجمه‌ها}\label{ux62aux62dux644ux6ccux644-ux62aux637ux628ux6ccux642ux6cc-ux62aux631ux62cux645ux647ux647ux627}

برای ارزیابیِ این که پنج ترجمه فارسیِ معتبر چگونه با تمایزِ شدتیِ پیوستار در آیات کلیدی برخورد کرده‌اند، جدول ۶ را تنظیم کرده‌ایم. آیات منتخب: حجر ۹۷ (ضیق)، آل عمران ۱۳۴ (غیظ)، ملک ۸ (تمیّز من الغیظ)، علق ۶ (طغی).

\textbf{جدول ۶.} تحلیل تطبیقیِ ترجمه‌ها در چهار آیه کانونی

\begin{longtable}[]{@{}
  >{\centering\arraybackslash}p{(\linewidth - 10\tabcolsep) * \real{0.1667}}
  >{\raggedright\arraybackslash}p{(\linewidth - 10\tabcolsep) * \real{0.1667}}
  >{\raggedright\arraybackslash}p{(\linewidth - 10\tabcolsep) * \real{0.1667}}
  >{\raggedright\arraybackslash}p{(\linewidth - 10\tabcolsep) * \real{0.1667}}
  >{\raggedright\arraybackslash}p{(\linewidth - 10\tabcolsep) * \real{0.1667}}
  >{\raggedright\arraybackslash}p{(\linewidth - 10\tabcolsep) * \real{0.1667}}@{}}
\toprule\noalign{}
\begin{minipage}[b]{\linewidth}\centering
واژه
\end{minipage} & \begin{minipage}[b]{\linewidth}\raggedright
فولادوند
\end{minipage} & \begin{minipage}[b]{\linewidth}\raggedright
مکارم شیرازی
\end{minipage} & \begin{minipage}[b]{\linewidth}\raggedright
الهی قمشه‌ای
\end{minipage} & \begin{minipage}[b]{\linewidth}\raggedright
معزی
\end{minipage} & \begin{minipage}[b]{\linewidth}\raggedright
انصاریان
\end{minipage} \\
\midrule\noalign{}
\endhead
\bottomrule\noalign{}
\endlastfoot
ضیق (حجر:۹۷) & تنگ می‌شود & تنگ می‌شود & غمگین می‌شوی & تنگ گردد & تنگ می‌شود \\
غیظ (آل عمران:۱۳۴) & خشم خود را فرو می‌برند & خشم خود را فرو می‌برند & خشم خود را فرو نشانند & خشم را فرو می‌خورند & خشم خود را فرو می‌برند \\
تمیّز من الغیظ (ملک:۸) & چنان است که می‌خواهد از خشم بترکد & از شدت خشم پاره‌پاره می‌شود & از خشم متلاطم است & چنانکه از خشم بشکافد & از شدت خشم نزدیک است متلاشی شود \\
طغیان (علق:۶) & سرکشی می‌کند & طغیان می‌کند & به طغیان برمی‌خیزد & بسرکشی پردازد & طغیان می‌کند \\
\end{longtable}

نکات این تطبیق:

اولاً، در آیه حجر ۹۷، الهی قمشه‌ای ضیق را به «غمگین می‌شوی» ترجمه کرده که هم‌میدانیِ ضیق و حزن را نشان می‌دهد ولی تمایز شدتیِ آن دو را از میان می‌برد---خطایی که در پنج ترجمه دیگر مشاهده نمی‌شود.

دوماً، در آیه ملک ۸، چهار ترجمه از پنج، تصویرِ «فروپاشیِ ظرف» (تمیّز) را با عبارات قابل‌قبولی منتقل می‌کنند («می‌خواهد بترکد»، «پاره‌پاره می‌شود»، «بشکافد»، «متلاشی شود»)؛ تنها ترجمه الهی قمشه‌ای با «متلاطم است»، عمقِ تصویرِ فروپاشی را تضعیف می‌کند.

سوماً، در همه‌ی ترجمه‌های مورد بررسی، \textbf{تمایز ساختاریِ غضب از غیظ} تنها در آیه آل عمران ۱۳۴ تا حدی منعکس می‌شود (با عبارت «فرو می‌برند» که مهارِ غیظ را برجسته می‌کند). ولی هیچ‌یک از این ترجمه‌ها، در ترجمه آیه ملک ۸، از نشانه‌ی صرفیِ متفاوتی برای انتقال این که «غیظ» یک حالتِ متراکم‌تر از غضب است، استفاده نمی‌کند.

این تحلیل، یافته‌ای روشن دارد: ترجمه‌های فارسیِ موجود، در سطحِ تک‌آیه به‌خوبی عمل می‌کنند ولی در حفظِ ساختارِ شبکه‌ایِ پیوستار، نقصی سامانمند دارند. این نقص، در بخش بحث به‌مثابه پایه‌ی پیشنهاد عملیِ ترجمه‌ای پژوهش حاضر مطرح می‌شود.

\begin{center}\rule{0.5\linewidth}{0.5pt}\end{center}

\subsection{۵. بحث}\label{ux628ux62dux62b}

\subsubsection{۵-۱. منطقِ تشدیدیِ خشم: یک نظریه‌ی واحد}\label{ux645ux646ux637ux642-ux62aux634ux62fux6ccux62fux6cc-ux62eux634ux645-ux6ccux6a9-ux646ux638ux631ux6ccux647ux6cc-ux648ux627ux62dux62f}

یافته‌های بخش پیشین، در یک قاب تحلیلیِ واحد، یک کشفِ نظریِ کانونی را آشکار می‌سازد: قرآن کریم، در میدانِ معناییِ طیفِ هیجانیِ خشم، یک «منطقِ تشدیدیِ خشم» \lr{(logic of anger escalation)} را به‌نحوی نظام‌مند رمزگذاری کرده است. در این منطق، نشانه‌های خشم نه به‌مثابه پدیده‌هایی گسسته، بلکه به‌مثابه نقاطی متصل بر روی یک پیوستار درجه‌بندی‌شده درک می‌شوند.

این منطق در سه تحول هم‌زمان نمایان می‌شود (شکل ۱):

\textbf{تحول مکانِ هیجان:} هیجان از ساحتِ کلامیِ کم‌شدت (أُف) و درونِ فرد (کراهت، ضِيق الصدر، حُزن القلب) آغاز می‌شود، به ساحتِ ارزیابانه (نَقم، سَخَط، مَقت) و سپس بیانیِ فعال (غَضَب، حَرد) می‌رسد، و در نهایت به ساحتِ کنشیِ بیرونی-ساختاری (بَغی، طُغیان، عُتُوّ) می‌رسد.

\textbf{تحول شدت:} شدتِ هیجان، از کم‌ترین تَرازِ کلامی (أف) به فشارِ درونی، نفرتِ ارزیابانه، خشمِ فعال، و سرانجام به حالتی متراکم در آستانه‌ی فروپاشی (غَیظ، تَمَیُّز) ارتقا می‌یابد.

\textbf{تحول حوزه‌ی تأثیر:} حوزه‌ی پیامدهای هیجان، از کلامِ فردی (مرحله ۱) به روانِ فردی (مرحله ۲) و رفتارِ فردی (مرحله‌های ۳ تا ۵) و سرانجام به اختلالِ ساختارِ اجتماعی (مرحله ۶) گسترش می‌یابد.

این منطق را می‌توان به‌صورتِ یک فرمولِ فشرده بیان کرد:

\begin{quote}
\textbf{پیش‌از‌خشم} → \textbf{فشارِ درونی} → \textbf{نفرتِ ارزیابانه} → \textbf{خشمِ فعال} → \textbf{خشمِ متراکم/انفجاری} → \textbf{پیامدهای رفتاری (با علیتِ دوسویه)}
\end{quote}

با تأکید بر این‌که گذار به مرحله ۶، تنها از طریقِ شکستِ تنظیمِ مراحلِ پیشین رخ نمی‌دهد، بلکه می‌تواند از مسیرهای مستقلِ ساختاری (طمع، استکبار) نیز ظاهر گردد. به همین جهت، نزدیکیِ آماری مراحلِ ۱ تا ۵ (۱۶۷) و مرحله ۶ (۱۴۵) که در §۴-۲ گزارش شد، نه پارادوکس بلکه پشتیبانیِ مهم برای علیتِ دوسویه‌ی مرحله ۶ به‌حساب می‌آید.

\subsubsection{۵-۲. بازخوانیِ ساختارِ مرحله‌ای از منظرِ مدلِ فرآیندیِ گراس}\label{ux628ux627ux632ux62eux648ux627ux646ux6cc-ux633ux627ux62eux62aux627ux631-ux645ux631ux62dux644ux647ux627ux6cc-ux627ux632-ux645ux646ux638ux631-ux645ux62fux644-ux641ux631ux622ux6ccux646ux62fux6cc-ux6afux631ux627ux633}

مدلِ فرآیندیِ تنظیم هیجانِ گراس \lr{(Gross 1998, 2014, 2015)} که در روان‌شناسیِ معاصر به‌مثابهِ چارچوبِ مرجع شناخته می‌شود، تنظیم هیجان را در گام‌های متوالی فرض می‌کند. شش مرحله‌ی پیوستارِ قرآنی، با گام‌های این مدل به‌نحو ساختاری منطبق می‌شوند، با این ملاحظه که قرآن سویِ عاطفی را ریزتر تقسیم می‌کند و یک ثبتِ صریحِ پیامدهای رفتاری نیز اضافه می‌نماید:

\begin{longtable}[]{@{}
  >{\centering\arraybackslash}p{(\linewidth - 4\tabcolsep) * \real{0.3333}}
  >{\raggedright\arraybackslash}p{(\linewidth - 4\tabcolsep) * \real{0.3333}}
  >{\raggedright\arraybackslash}p{(\linewidth - 4\tabcolsep) * \real{0.3333}}@{}}
\toprule\noalign{}
\begin{minipage}[b]{\linewidth}\centering
مرحله‌ی قرآنی
\end{minipage} & \begin{minipage}[b]{\linewidth}\raggedright
فازِ گراس
\end{minipage} & \begin{minipage}[b]{\linewidth}\raggedright
تطبیقِ نظری
\end{minipage} \\
\midrule\noalign{}
\endhead
\bottomrule\noalign{}
\endlastfoot
۱ --- أف، کره (پیش‌از‌خشم) & پیشایی-محور: انتخابِ موقعیت & کم‌ترین نشانه‌های ناخشنودی؛ قابلِ تنظیم با اجتناب یا بازچارچوب‌بندی \\
۲ --- ضیق، حزن، أسف (فشار درونی) & تغییرِ موقعیت & انقباضِ درونی؛ کنشگر می‌تواند موقعیت را تعدیل کند \\
۳ --- نقم، سخط، مقت (ارزیابانه) & گسترشِ توجه / تغییرِ شناختی & پیامدِ ارزیابیِ شناختی-اخلاقی از موقعیت \\
۴ --- غضب، حرد (فعال) & تغییرِ شناختی / گزینشِ پاسخ & پاسخِ حرکتی-بسیج‌شده به موقعیتِ ارزیابی‌شده \\
۵ --- غیظ + کظم؛ تمیّز & تعدیل پاسخ (بازداری/شکست) & تنظیمِ درونیِ مستقیم؛ تمیّز نشانه‌ی شکستِ مهار \\
۶ --- بغی، طغیان، عتو (با علیتِ دوسویه) & فراتر از تنظیم: اختلالِ بیرونی‌شده & جایی که گراس ساکت می‌شود --- و قرآن، روان-انسان‌شناسیِ اخلاقی \textbf{و} ساختارگراییِ مستقل را اضافه می‌کند \\
\end{longtable}

این تطبیق برای پنج مرحله‌ی نخست دقیق است. مرحله‌ی ششم، اما، جایی است که پیوستارِ قرآنی از مدلِ گراس \textbf{فراتر} می‌رود: گراس بر موفقیتِ تنظیم متمرکز است و پیامدهای شکست را به‌مثابهِ یک طبقه‌ی باقیمانده در نظر می‌گیرد، در حالی که پیوستارِ قرآنی پاتولوژی‌های ساختاری-اخلاقیِ قطبِ رفتاری را به‌نحوِ نظام‌مند توصیف می‌کند --- پاتولوژی‌هایی که، چنان‌که در §۴-۸-۵ استدلال کردیم، \textbf{هم} می‌توانند پیامدِ خشم باشند \textbf{و هم} می‌توانند علتِ آن. قرآن بنابراین نه‌تنها یک مدل تنظیمِ خشم، بلکه \textbf{یک مدل ساختاری-اخلاقی برای شکستِ تنظیم و نیز برای فرماسیون‌های مستقلِ شخصیتی} ارائه می‌دهد.

\subsubsection{۵-۳. مقایسه با چرخه پلاتچیک، STAXI، و نظریه ارزیابی شناختی}\label{ux645ux642ux627ux6ccux633ux647-ux628ux627-ux686ux631ux62eux647-ux67eux644ux627ux62aux686ux6ccux6a9-staxi-ux648-ux646ux638ux631ux6ccux647-ux627ux631ux632ux6ccux627ux628ux6cc-ux634ux646ux627ux62eux62aux6cc}

پیوستارِ قرآنی را می‌توان با سه چارچوب دیگر روان‌شناسیِ هیجانِ معاصر نیز در گفت‌وگو قرار داد:

\textbf{چرخه‌ی پلاتچیک \lr{(Plutchik 1980, 2001)}:} خشم در مدلِ پلاتچیک، در یک طیفِ شدتیِ سه‌گانه (آزردگی → خشم → خشمِ شدید) سامان‌بندی می‌شود. این ساختار، در پیوستار قرآنی به‌نحوِ شایانِ توجهی بازتاب دارد: مراحلِ ۱ تا ۵ ≈ خشمِ صفتیِ پلاتچیکی در شدت‌های متوالی. ولی \textbf{مرحله ۶ قرآنی هیچ معادلِ مستقیمی در مدل پلاتچیک ندارد}؛ یعنی پیوستار قرآنی، بُعدِ کنشی-ساختاری را --- هرچند با احتیاطِ علیتِ دوسویه --- به مدل سه‌مرحله‌ای پلاتچیک افزوده است.

\textbf{STAXI اسپیلبرگر \lr{(Spielberger 1999)}:} سامانه‌ی سنجشِ خشمِ صفت-حالِ اسپیلبرگر، تفکیکِ سه‌گانه‌ای میان «خشم صفتی» (پایه‌ی عاطفیِ پایدار)، «خشم حالتی» (پاسخِ هیجانیِ موقت)، و «کنترل خشم» (مهار خودخواسته) قائل می‌شود. این تفکیک، با تمایز سه‌گانه‌ی قرآنیِ مراحلِ پایدارِ ضیق/حزن (به‌مثابه خشم صفتی)، مراحلِ ۳ و ۴ (واکنش‌های ارزیابانه و حرکتی به‌مثابه خشم حالتی)، و کظم غیظ (به‌مثابه کنترل خشم) قابل تطبیق است.

\textbf{نظریه ارزیابیِ شناختیِ لازاروس و مدل فرآیندیِ شِرر \lr{(Lazarus 1991; Scherer 2001, 2005)}:} هر دو نظریه، هیجان را به‌مثابه پاسخی به یک ارزیابیِ شناختیِ چندلایه از موقعیت تعریف می‌کنند. آیه ۶ سوره علق («إِنَّ الْإِنسَانَ لَيَطْغَى أَن رَّآهُ اسْتَغْنَى») با این چارچوب به‌نحوِ خیره‌کننده‌ای هم‌سو است: طغیان، در این آیه، نه یک هیجانِ خام، بلکه پاسخی به یک ارزیابیِ شناختیِ خاص (ادراکِ بی‌نیازی) معرفی می‌شود.

این هم‌سویی‌های متعدد، نشان می‌دهد که گفتمانِ قرآنی، در دستگاه واژگانیِ خود، یک نظریه‌ی پیشرفته از روان‌شناسیِ هیجان را---نه به‌صورتِ صورت‌بندیِ صریح نظری، بلکه به‌صورتِ معماریِ پنهان معناشناسیِ متنی---در خود نهفته دارد.

\subsubsection{۵-۴. تعادلِ پدیدارشناختی/پیامدها: چرا پیکره دیگر یک کجیِ ۲٫۴ برابره را نشان نمی‌دهد}\label{ux62aux639ux627ux62fux644-ux67eux62fux6ccux62fux627ux631ux634ux646ux627ux62eux62aux6ccux67eux6ccux627ux645ux62fux647ux627-ux686ux631ux627-ux67eux6ccux6a9ux631ux647-ux62fux6ccux6afux631-ux6ccux6a9-ux6a9ux62cux6cc-ux6f2ux6f4-ux628ux631ux627ux628ux631ux647-ux631ux627-ux646ux634ux627ux646-ux646ux645ux6ccux62fux647ux62f}

یک بازنگریِ مرکزی نسبت به نسخه‌ی پیشینِ این تحلیل، به وزنِ نسبیِ هسته‌ی پدیدارشناختی (مراحلِ ۱ تا ۵) و قطبِ پیامدهای رفتاری (مرحله ۶) مربوط می‌شود. در تحلیلِ ده‌ریشه‌ایِ پیشین، پیکره چنین می‌نمود که ۲٫۴ برابر بیشتر از سه مرحله‌ی نخست را به مرحله ۴ (که آن‌گاه «کنشِ عصیانگرانه» نام داشت) اختصاص می‌دهد؛ در یک قرائتِ ساده‌انگارانه، این کجی نشانه‌ی تأکیدِ گفتمانی بر پیامدها بود.

تحلیلِ بازنگری‌شده‌ی چهارده‌ریشه‌ایِ شش‌مرحله‌ای، روایتی متفاوت ارائه می‌دهد. فهرستِ گسترش‌یافته‌ی واژگانِ پدیدارشناختی (أف، کره، نقم، مقت، حرد) مجموعِ مراحلِ ۱ تا ۵ را به ۱۶۷ می‌رساند، در برابرِ ۱۴۵ مرحله ۶؛ آزمونِ خی-دو در برابرِ تقسیمِ یکنواخت \textbf{رد نمی‌شود} (\lr{χ²(1) = 1.55}). این نزدیکی، در قرائتِ ما، دقیقاً همان چیزی است که در صورتِ صحتِ این دو فرض انتظار می‌رود: (۱) سویِ پدیدارشناختیِ پیوستار غنی‌تر از آن است که تحلیلِ ده‌ریشه‌ایِ پیشین اجازه می‌داد، و (۲) ریشه‌های مرحله ۶ صرفاً از خشم برنمی‌خیزند و در رابطه‌ی \textbf{دوسویه‌ی} علّی با خشم قرار دارند. واقعیتِ پیکره‌ای بنابراین با ملاحظه‌ی مفسرانه‌ی ما درباره‌ی بَغی، طُغیان و عُتُوّ سازگار است --- و در واقع آن را پشتیبانی می‌کند.

این، یک پیامدِ روش‌شناختیِ مهم دارد: تحلیلِ توزیعیِ پیکره‌ای، نباید به‌مثابه‌ی بازنماییِ خنثای واقعیتِ بیرونی تلقی شود. این تحلیل‌ها، نقشه‌هایی از \textbf{تأکیدهای گفتمانیِ} متن‌اند؛ تفسیرشان به‌نحوِ بنیادی با تصمیماتِ نظری درباره‌ی عضویتِ ریشه‌ها (کدام ریشه‌ها به طیف تعلق دارند) و علیتِ آنها (آیا رابطه یک‌سویه است یا دوسویه) درآمیخته است. تغییرِ از تحلیلِ ۱۰/۴-مرحله‌ای به ۱۴/۶-مرحله‌ای، صرفاً آرایشی نبوده: هم الگوی تجربی و هم تفسیرِ نظریِ آن را تغییر داده.

\subsubsection{۵-۵. الگوی تخصص‌یافتگیِ بافتیِ مدنی/مکی}\label{ux627ux644ux6afux648ux6cc-ux62aux62eux635ux635ux6ccux627ux641ux62aux6afux6cc-ux628ux627ux641ux62aux6cc-ux645ux62fux646ux6ccux645ux6a9ux6cc}

چندین یافته‌ی توزیعی نظریه‌نمای مهم در پژوهش حاضر، الگوی انحصاریِ توزیعِ سَخَط (۰ مکی / ۴ مدنی) و أَسَف (۵ مکی / ۰ مدنی)، گرایشِ مکیِ نَقم (۱۲/۵) و عُتُوّ (۹/۱)، هاپَکس‌مکیِ حَرد، و گرایشِ مدنیِ کره (۱۴/۲۷) را در بر می‌گیرد. تحلیلِ بافتاری نشان می‌دهد که این دو الگو، نمونه‌هایی از پدیده‌ی گسترده‌تری به‌نام \textbf{«تخصص‌یافتگیِ گفتمانی-بافتیِ یک واژه»} \lr{(discursive-contextual specialization of a lexeme)} هستند: هر دو واژه، کارکردِ معناییِ خود را در یک بافتارِ گفتمانیِ خاص تثبیت کرده‌اند.

سَخَط در گفتمانِ مدنی، به‌طور خاص، در مواجهه با منافقان و شبه‌مؤمنانِ نیمه‌پیمان شکل می‌گیرد. این بافت گفتمانی در سور مکی غایب است (چرا که مفهومِ منافقین در ساختارِ گفتمانیِ مکی، که عمدتاً به مواجهه با مشرکانِ آشکار می‌پردازد، شکل نگرفته است). از سوی دیگر، أَسَف در روایت‌های پیامبرانه‌ی پیشین (به‌ویژه ماجرای موسی (ع) با فرعون و قومش، و آن‌چه در سور مکی پربسامد است) به‌مثابه یک حالتِ گذاریِ ادبی-روایی به‌کار می‌رود. در فضای مدنی، که عمدتاً به سامان‌دهیِ نخستینِ امّت و رویاروییِ عملی با چالش‌های نوظهور می‌پردازد، نیازی به این حالتِ گذاریِ روایی نیست.

این الگوی تخصص‌یافتگی، یک یافته‌ی نظریِ بدیع را پشتیبانی می‌کند:

\begin{quote}
\textbf{پدیده تخصص‌یافتگیِ گفتمانیِ-بافتی:} در دستگاه واژگانیِ قرآن، برخی از واژگانِ هم‌میدان نه به‌صورتِ تصادفی، بلکه به‌صورتِ نظام‌مند، در یک بافتارِ گفتمانیِ خاص تثبیت شده‌اند. این تخصص‌یافتگی، یک نوع تنظیمِ معناییِ ظریف است که در سطحِ تحلیل تک‌آیه‌ای از دیدِ پژوهشگر مخفی می‌ماند ولی در سطحِ تحلیلِ پیکره‌ای آشکار می‌شود.
\end{quote}

این یافته را می‌توان با پدیده‌های مشابه در زبان‌شناسیِ عمومی (نظیر «تخصصِ ثبتی» در زبان‌های با چندین رجیستر) و در مطالعاتِ Sitz im Leben در نقدِ متن‌شناختیِ قرآنی هم‌مرز دانست.

\subsubsection{۵-۶. آیه ۶۷:۸ به‌مثابه نمونه‌ی شاخصِ «فروپاشیِ ظرف»}\label{ux622ux6ccux647-ux6f6ux6f7ux6f8-ux628ux647ux645ux62bux627ux628ux647-ux646ux645ux648ux646ux647ux6cc-ux634ux627ux62eux635-ux641ux631ux648ux67eux627ux634ux6cc-ux638ux631ux641}

یکی از سهم‌های نظریِ شایانِ توجه پژوهش حاضر آن است که آیه‌ی ۸ سوره‌ی ملک، \textbf{یک نمونه‌ی بسیار شفاف از پدیده‌ی «فروپاشیِ ظرف» در یک متنِ پیشامدرن} است (با این تذکر که ادعاهایِ تطبیقیِ قوی‌تر، نیازمندِ پیمایشِ نظام‌مندِ ادبیاتِ مذکور است). لیکاف و کوچش (۱۹۸۷) و کوچش (۲۰۰۰) این پدیده را به‌مثابه «گذار از فشار به انفجار» در زبان انگلیسی شناسایی کردند، و آن را با عبارت‌هایی نظیر «\lr{to explode with rage}», «\lr{she blew her top}», «\lr{I lost it}» نشان دادند. ولی هیچ‌یک از این مثال‌های انگلیسی به وضوحِ تصویر قرآنیِ آیه ۶۷:۸ نزدیک نیست. در حالی که عبارات انگلیسی، \textbf{انفجار را گزارش می‌کنند}، آیه قرآنی \textbf{فروپاشیِ ساختاری خود ظرف را به‌نحو صریح بازنمایی می‌کند}.

این نکته، چندین پیامدِ نظری دارد. اولاً، آیه ۶۷:۸ شواهدی برای \textbf{عمومیت تاریخی} استعاره مفهومیِ «خشم به‌مثابه‌ی ظرف فشرده» فراهم می‌آورد---نشان می‌دهد که این استعاره صرفاً یک ویژگیِ زبان مدرن غربی نیست، بلکه در ساختار ادراکیِ بنیادینِ بشر، احتمالاً جهان‌شمول است. دوماً، آیه ۶۷:۸ نمونه‌ای از آن چیزی است که سوک‌چِلسکی \lr{(Sokolewicz 2018)} «استعاره استعاره» \lr{(meta-metaphor)} می‌نامد: استعاره‌ای که خود نسبت به استعاره‌ی پایه \lr{(anger as pressurized container)} یک گام جلوتر می‌رود و فروپاشیِ خودِ ساختار را تصویر می‌کند.

\begin{figure}
\centering
\includegraphics[width=1\linewidth,height=\textheight,keepaspectratio]{fig6_metaphor_diagram.pdf}
\caption{طرحواره‌ی «خشم به‌مثابهِ ظرفِ تحتِ فشار» (لیکاف و کوچش، ۱۹۸۷) که با مثلثِ قرآنیِ غضب → غیظ + کظم → تمیّز عملیاتی شده است. \textbf{پنلِ نخست} (مرحله ۴): غضب --- ظرفِ \emph{باز}؛ هیجان حاضر است ولی متراکم نشده، و آزادانه به سمتِ زبان (بیشترین مولّدبودنِ نومینال در طیف) تخلیه می‌شود. \textbf{پنلِ دوم} (مرحله ۵، آل عمران: ۱۳۴): غیظ --- مایعِ مهر شده با کَظم؛ فعلِ کَظَمَ در اصل بستن دهانه‌ی مَشک برای جلوگیری از سرریز را افاده می‌کند، یعنی مهارِ درون‌فردیِ کنشگر که در برابرِ افزایشِ فشار، وضعیتِ بسته بودن را حفظ می‌کند. \textbf{پنلِ سوم} (تجلیِ مرحله ۵، ملک: ۸): تمیّز --- شکستِ ساختاری؛ این فعل از ریشه‌ی م-ی-ز («جدا شدن، شکافته شدن») می‌آید و نه خروجِ محتوای فشرده (همچون \lr{explode} انگلیسی) بلکه \emph{جدا شدنِ خودِ ظرف} را تصویر می‌کند. عینیت‌یابیِ قرآنیِ این طرحواره، بنا بر استدلالِ پژوهشِ حاضر، شفاف‌تر از هر نمونه‌ی انگلیسی در پیکره‌ی لیکاف و کوچش (۱۹۸۷) است و شواهدِ بین‌زبانیِ نیرومندی برای جهان‌شمولیِ شناختیِ این طرحواره فراهم می‌آورد (§۵-۶).}
\end{figure}

\subsubsection{۵-۷. پیامدهای کاربردی برای ترجمه قرآن}\label{ux67eux6ccux627ux645ux62fux647ux627ux6cc-ux6a9ux627ux631ux628ux631ux62fux6cc-ux628ux631ux627ux6cc-ux62aux631ux62cux645ux647-ux642ux631ux622ux646}

تحلیل تطبیقیِ ترجمه‌ها (بخش ۴-۱۰، جدول ۶) نشان می‌دهد که ترجمه‌های فارسیِ موجود، در حفظ ساختار شبکه‌ایِ پیوستار، نقصی سامانمند دارند. یافته‌های پژوهش حاضر سه توصیه عملی برای ترجمه‌های آینده دارد:

\textbf{توصیه ۱.} برای حفظ تمایز شدتیِ پیوستار، مترجمان باید معادل‌های متمایزی برای هر یک از چهارده واژه‌ی کانونی به‌کار برند، حتی اگر این معادل‌ها در نگاه نخست، دست‌و‌پاگیر به نظر برسند. مثال پیشنهادی: أف → آهِ ناخوشایندی؛ کره (کراهت) → ناخواستنیِ درونی؛ کره (إکراه) → تحمیلِ ساختاری (که خود هیجان نیست)؛ ضیق → تنگیِ درونی؛ ضاق ذرعاً → تنگیِ توانِ کنش / پرخاشگریِ بازداشته؛ حزن → اندوه (با ارزش‌های متعدد)؛ أسف → افسوسِ آمیخته با خشم؛ نقم → نکوهشِ همراه با اراده‌ی انتقام؛ سخط → خشمِ اخلاقیِ مستمر؛ مقت → نفرتِ اخلاقیِ شدید؛ غضب → خشم؛ حرد → خشمِ همراه با قصدِ منع؛ غیظ → خشمِ متراکم؛ تمیّز → از هم گسستن از خشم؛ بغی → تعدی؛ طغیان → طغیان؛ عتوّ → سرکشیِ لجوجانه.

\textbf{توصیه ۲.} در ترجمه‌های پژوهشی، تعبیه‌ی پاورقی یا فهرستِ ضمیمه‌ای که جایگاه هر واژه در پیوستار شدتی را روشن می‌سازد، ضرورت دارد.

\textbf{توصیه ۳.} در ترجمه‌های کاربردی-عمومی، استفاده از یک دستگاه نمادهایِ تصویری (نظیر نمادهای رنگی برای هر مرحله) می‌تواند به خواننده کمک کند تا ساختار پیوستار را در طول قرائت متن دنبال کند.

\subsubsection{۵-۸. پیامدها برای روان‌شناسی اسلامیِ کاربردی}\label{ux67eux6ccux627ux645ux62fux647ux627-ux628ux631ux627ux6cc-ux631ux648ux627ux646ux634ux646ux627ux633ux6cc-ux627ux633ux644ux627ux645ux6cc-ux6a9ux627ux631ux628ux631ux62fux6cc}

یافته‌های پژوهش حاضر، چارچوبی نظری برای طراحی مداخلات روان‌شناختیِ مبتنی بر فرهنگ اسلامی فراهم می‌آورد:

\textbf{در سطحِ پیشگیری}، این چارچوب می‌تواند به فرد کمک کند که نشانه‌های خشم را در سطحِ مراحلِ ۱ تا ۲ (أف، کره، ضیق، حزن) شناسایی کند، یعنی پیش از آن که به مراحلِ بعدی ارتقا یابد. این رویکرد با اصلِ «شناختِ زودهنگام» \lr{(early recognition)} در روان‌درمانیِ شناختی-رفتاری همخوان است.

\textbf{در سطحِ مداخله‌ی بحرانی}، چارچوبِ کظم غیظ به‌مثابهِ استراتژیِ تعدیلِ پاسخِ مرحله ۵، می‌تواند با تکنیک‌های مدرنِ تنظیم هیجان (نظیر بازنگری شناختی، تنفس عمیق، توجه‌آگاهی) ترکیب شود.

\textbf{در سطحِ توان‌بخشی}، چارچوبِ پیوستار به فرد در حالتِ مرحله ۶ (که در ادبیاتِ بالینی به اختلال‌های خشمِ بیمارگونه نزدیک است) ارائه می‌دهد که مسیرِ معکوس به سمت مراحلِ پیشین، نه‌تنها ممکن، بلکه از منظرِ متنیِ قرآنی، یک «بازگشت» (توبه) قابلِ دستیابی است.

\begin{center}\rule{0.5\linewidth}{0.5pt}\end{center}

\subsection{۶. نتیجه‌گیری}\label{ux646ux62aux6ccux62cux647ux6afux6ccux631ux6cc}

\subsubsection{۶-۱. سهم‌های پژوهش}\label{ux633ux647ux645ux647ux627ux6cc-ux67eux698ux648ux647ux634}

پژوهش حاضر سه سهمِ متمایز را در حوزه‌های پژوهشی مرتبط ارائه می‌کند:

\textbf{سهم روش‌شناختی.} برای نخستین بار در پژوهش‌های قرآنی، چارچوبِ ترکیبیِ سه‌گانه «معناشناسی میدانیِ ایزوتسو + معناشناسی طیفیِ کندی-ساسون + نظریه استعاره مفهومیِ لیکاف-کوچش» را در کنار تحلیلِ پیکره‌ای محاسباتی و بازتولیدپذیر به‌کار می‌گیرد. این ادغام، نشان می‌دهد که هیچ‌یک از سه سنت به‌تنهایی برای واکاویِ پیچیدگیِ معماریِ معناییِ متونِ مقدس کفایت نمی‌کند، ولی ترکیب آن‌ها در یک ساختار سه‌لایه‌ای، ابزاری توانمند فراهم می‌آورد که می‌تواند به سایر میدان‌های مفهومیِ قرآن (نظیر شادی-فرح، ایمان-اطمینان، بیم-امید) نیز گسترش یابد.

\textbf{سهمِ توصیفی.} پژوهش حاضر یک کنکوردانسِ کاملاً بازتولیدپذیر و در دسترسِ عموم برای ۳۱۲ بسامد در چهارده ریشه‌ی کانونیِ طیفِ هیجانیِ خشم در قرآن فراهم آورده است.

\textbf{سهمِ نظری.} پژوهش حاضر «فرضیه‌ی پیوستارِ شدتیِ قرآنی» را به‌مثابه‌ی یک چارچوبِ نظریِ کانونی صورت‌بندی می‌کند: گزاره‌ی مرکزیِ این فرضیه آن است که میدانِ معناییِ طیفِ هیجانیِ خشم در قرآن، نه شبکه‌ای از واژگانِ مترادف، بلکه یک پیوستارِ درجه‌بندی‌شده بر روی محورِ «شدتِ کنش» است که در شش مرحله‌ی متمایز سامان‌بندی می‌شود. این فرضیه، با شواهدِ کیفی و کمّی پشتیبانی می‌شود. ملاحظه‌ی صریح بر علیتِ دوسویه‌ی مرحله ۶ بخشِ مهمی از این صورت‌بندیِ بازنگری‌شده است.

\subsubsection{۶-۲. پاسخ به سؤال‌های پژوهشی}\label{ux67eux627ux633ux62e-ux628ux647-ux633ux624ux627ux644ux647ux627ux6cc-ux67eux698ux648ux647ux634ux6cc}

\textbf{پاسخ به س۱:} میدانِ معناییِ طیفِ هیجانیِ خشم در قرآن، یک شبکه‌ی چهارده‌گرهی است که با مرکزیتِ بَغی (به‌مثابهِ گرهِ پلی)، مرکزیتِ حُزن و غَضَب (به‌مثابهِ گره‌های مرحله‌ای)، و حضورِ چندین گرهِ منزویِ تخصصی (سَخَط، أَسَف، حَرد، عُتُوّ) ساختار یافته است.

\textbf{پاسخ به س۲:} این چهارده گره را می‌توان بر روی یک پیوستارِ شش‌مرحله‌ای رده‌بندی کرد: مرحله ۱ (أف، کره)، مرحله ۲ (ضیق، حزن، أسف)، مرحله ۳ (نقم، سخط، مقت)، مرحله ۴ (غضب، حرد)، مرحله ۵ (غیظ، با تمیّز به‌مثابهِ تجلی)، و مرحله ۶ (بغی، طغیان، عتو، با ملاحظه‌ی علیتِ دوسویه).

\textbf{پاسخ به س۳:} الگوی توزیعیِ پیکره‌ای، رده‌بندیِ نظریِ پیشنهادی را تأیید می‌کند: (الف) آزمونِ خی-دوی فرضِ یکنواختی شش‌مرحله‌ای رد می‌شود (\lr{χ²(5) = 227.15, p < 0.001})؛ (ب) آزمونِ مقایسه‌ی هسته‌ی پدیدارشناختی با قطبِ پیامدهای رفتاری \textbf{رد نمی‌شود} (\lr{χ²(1) = 1.55})، که در قرائتِ ما، پشتیبانی برای علیتِ دوسویه‌ی مرحله ۶ است؛ (ج) شبکه‌ی هم‌رخدادگیریِ آیه‌ای ساختارِ پیش‌بینی‌شده‌ی پل را در بَغی نشان می‌دهد؛ (د) پروفایلِ ریخت‌شناختی، مراحلِ فعلیِ ۱ و ۲ را از مرحله ۶ اسم‌محور متمایز می‌کند.

\textbf{پاسخ به س۴:} پیوستارِ قرآنی با چهار مدلِ مدرنِ روان‌شناسیِ هیجان \lr{(Gross, Plutchik, Spielberger, Lazarus-Scherer)} به‌نحوِ ساختاری هم‌سو است، ولی از حیثِ گستره‌ی مفهومی، فراتر از این مدل‌ها می‌رود: قرآن مرحله ۶ را به‌نحوِ نظام‌مند بازنمایی می‌کند --- همراه با ملاحظه‌ی §۴-۸-۵ بر علیتِ دوسویه.

\textbf{پاسخ به س۵:} پیامدهای پژوهش برای ترجمه و روان‌شناسیِ اسلامیِ کاربردی در توصیه‌های عملی صورت‌بندی شده است: (الف) استفاده از معادل‌های متمایز برای چهارده واژه‌ی کانونی در ترجمه‌های آینده؛ (ب) طراحی مداخلاتِ روان‌شناختیِ پیشگیرانه در سطحِ مراحلِ ۱ تا ۲؛ (ج) ادغامِ چارچوبِ کظم غیظ با تکنیک‌های مدرنِ تنظیمِ هیجان.

\subsubsection{۶-۳. محدودیت‌ها و افق پژوهش‌های آینده}\label{ux645ux62dux62fux648ux62fux6ccux62aux647ux627-ux648-ux627ux641ux642-ux67eux698ux648ux647ux634ux647ux627ux6cc-ux622ux6ccux646ux62fux647}

پژوهش حاضر، با وجودِ سهم‌های مذکور، محدودیت‌هایی روش‌شناختی و مفهومی دارد که هم نقاطِ آغازِ پژوهش‌های آینده را مشخص می‌سازند و هم چارچوبِ معتبرِ تفسیرِ یافته‌ها را می‌سازند:

\paragraph{محدودیت‌های روش‌شناختیِ بنیادین.}

\textbf{(۱) قضاوتِ تفسیری در انتخابِ ریشه‌ها.} انتخابِ چهارده ریشه‌ی کانونی، با وجودِ معیارهای روشنِ مذکور در §۳-۲ (وقوعِ قرآنی، پشتیبانیِ لغوی در سه فرهنگ‌نامه، عضویتِ معناییِ مستقیم در میدانِ خشم، و حذفِ ریشه‌های اخلاقیِ بسیار گسترده مگر در نقشِ چترِ مفهومی)، در نهایت، بر یک قضاوتِ تفسیری استوار است. ریشه‌هایی نظیر «شنأ» (با ۳ بسامد) به دلیلِ پراکندگیِ بافتاری از پیوستارِ اصلی کنار گذاشته شده‌اند، در حالی که خوانندگانِ دیگر ممکن است گزینشِ متفاوتی داشته باشند. تکرارِ تحلیل با مجموعه‌های جایگزینِ ریشه‌ها (مثلاً افزودن «شنأ»، حذفِ «حرد» به‌مثابه هاپَکس) می‌تواند استواریِ یافته‌ها را در برابرِ این انتخاب بسنجد.

\textbf{(۲) بسامد ≠ شدت.} پژوهش حاضر، \emph{بسامدِ پیکره‌ای} را به‌مثابهِ شاخصی برای \emph{برجستگیِ گفتمانیِ} یک ریشه به‌کار می‌گیرد، نه به‌مثابهِ سنجشی مستقیم از \emph{شدتِ هیجانی}. ادعای مرکزیِ ما درباره‌ی \emph{ترتیبِ شدتی} میانِ ریشه‌ها، بر مدارکِ لغوی، تفسیری، استعاره-مفهومی، و ریخت‌شناختی استوار است؛ بسامد، نقشِ آزمونِ همگرایی را ایفا می‌کند، نه نقشِ مدرکِ اولیه. خواننده باید توجه داشته باشد که یک ریشه‌ی کم‌بسامد (نظیر سَخَط با ۴ یا حَرد با ۱) می‌تواند شدتِ تجربیِ بالایی داشته باشد در حالی که یک ریشه‌ی پربسامد (نظیر بَغی با ۹۶) عمدتاً نه شدتِ تجربی بلکه گستره‌ی گفتمانی-اخلاقی دارد.

\textbf{(۳) پیچیدگی‌های رده‌بندیِ مکی/مدنی.} رده‌بندیِ سور به مکی و مدنی، به‌ویژه برای سوره‌های میانه (مثلاً نحل، رعد) که ترکیبی از آیاتِ هر دو دوره را در بر دارند، در سنتِ علومِ قرآنی همواره موضوعِ اختلاف بوده است. خطِ پایه‌ی ۰٫۷۴ (نسبتِ آیاتِ مکی به کلِ آیات) که در آزمون‌های دوجمله‌ای به‌کار رفته، بر مبنای رده‌بندیِ سوره‌ای مستخرج از پیکره‌ی تنزیل محاسبه شده است؛ تحلیلِ حساسیت در دامنه‌ی ۰٫۷۰ تا ۰٫۷۸ نشان داد که نتیجه‌ی سَخَط (\lr{p < 0.01}) در سرتاسرِ این دامنه استوار باقی می‌ماند. با این وجود، تحلیلِ ظریف‌ترِ آیه‌محور (به‌جای سوره‌محور) که در پژوهش‌های آینده ممکن است، می‌تواند تصویرِ دقیق‌تری ارائه دهد.

\textbf{(۴) مدلِ مرحله‌ای، یک \emph{بازسازی} است نه یک \emph{فرآیندِ زمانی}.} مدلِ شش‌مرحله‌ایِ پیشنهادی، یک \emph{ترتیبِ منطقی-معنایی} است نه یک \emph{مسیرِ زمانیِ روان‌شناختی}. هیچ ادعایی مبنی بر اینکه «أف» همواره در زمان قبل از «غضب» تجربه می‌شود مطرح نشده. مرحله‌ها، نقاطِ متمایزی در یک طیفِ شدتیِ کنش‌ای هستند، نه ایستگاه‌های الزامیِ یک سفرِ هیجانیِ هر فرد.

\textbf{(۵) خطاهای احتمالیِ برچسب‌گذاریِ پیکره.} پیکره‌ی ریخت‌شناختیِ قرآنی (دوکس، ۲۰۱۱) در نسخه‌ی ۰٫۴ خود، اگرچه دقتِ بالایی دارد، احتمالاً موارد محدودی از خطاهای ریشه‌گذاری یا برچسبِ جزء سخن را در بر می‌گیرد. راستی‌آزماییِ دستیِ ۲۰ آیه‌ی نمونه‌گیری‌شده‌ی تصادفی در §۳-۵ نشان داد که نرخِ خطا در محدوده‌ی ۲٪ یا کم‌تر است؛ ولی این خطاها می‌توانند، به‌ویژه برای ریشه‌های کم‌بسامد، بر شمارش‌ها اثر بگذارند.

\textbf{(۶) پیامدهای ترجمه‌ای، نیازمندِ آزمونِ مستقل‌اند.} توصیه‌های ترجمه‌ای در §۵-۷ بر مبنایِ یک نمونه‌ی کوچک از پنج ترجمه‌ی فارسی-انگلیسی روی چهار آیه‌ی کانونی استوار است. ارزیابیِ نظام‌مند که نمونه‌ی تصادفیِ طبقه‌بندی‌شده‌ای از همه‌ی ۳۱۲ بسامد را در ده ترجمه‌ی شناخته‌شده پوشش دهد، یک پروژه‌ی مستقلِ آینده است.

\paragraph{محدودیت‌های دامنه و چارچوب.}

اولاً، \textbf{محدودیتِ سطحِ تحلیل}: تحلیل بر روی سطح واژگانیِ سطحی صورت گرفته است؛ گفتمان‌کاویِ بافتاریِ گسترده‌تر \lr{(analysis of wider discourse)} که آیات هر مرحله را در روایتِ کلیِ سوره‌ها بازخوانی می‌کند، در دستور کارِ پژوهش‌های آینده باید قرار گیرد.

دوماً، \textbf{محدودیتِ منابعِ تفسیری}: تحلیلِ تفسیری به‌صورتِ گزینشی به چهار تفسیرِ معتبر (المیزان، الکشاف، مجمع البیان، التحریر و التنویر) محدود شد. گسترشِ این تحلیل به سنتِ گسترده‌ترِ تفاسیر (به‌ویژه تفاسیر روایی-حدیثی، عرفانی، و معاصرِ فارسی) می‌تواند ابعاد جدیدی از پیوستار را آشکار سازد.

سوماً، \textbf{محدودیتِ چارچوبِ نظری}: چارچوبِ زبان‌شناسیِ شناختیِ به‌کار رفته در این پژوهش، عمدتاً ریشه در سنتِ آنگلوساکسون دارد. مقایسه با چارچوب‌های شناختیِ زبان‌شناسیِ عربی-اسلامیِ کلاسیک (به‌ویژه نظریه‌های جرجانی در دلائل الإعجاز) می‌تواند مدلی غنی‌تر و بومی‌تر فراهم آورد.

چهارماً، \textbf{محدودیتِ روایت قرآنی-روایی}: حدیث و سنتِ روایی، به‌مثابه‌ی مکمل قرآن، در این پژوهش به‌صورتِ نظام‌مند گنجانده نشده‌اند. آیا روایاتِ معصومان (ع) همین پیوستار را بازتاب می‌دهند یا گستره‌ای متفاوت ارائه می‌کنند؟ این پرسش، می‌تواند موضوع پژوهشِ بعدی باشد.

پنجماً، \textbf{محدودیتِ تحلیلِ بالینی}: پیامدهای پیشنهادیِ بالینی در بخش ۵-۸ به‌صورت مفهومی مطرح شده‌اند ولی هنوز در یک کارآزمایی بالینیِ تجربی آزموده نشده‌اند. طراحیِ یک مطالعه پایلوت برای ارزیابی اثربخشیِ مدلِ بومیِ تنظیم هیجانِ مبتنی بر این پیوستار، یک گامِ کاربردیِ مهم در پژوهش‌های آینده خواهد بود.

\subsubsection{۶-۴. سخن پایانی}\label{ux633ux62eux646-ux67eux627ux6ccux627ux646ux6cc}

در پایان، می‌توان گفت که این پژوهش، به‌گونه‌ای فروتنانه، نشان داد که قرآن کریم، در دستگاه واژگانیِ خود، یک نظریه‌ی پیشرفته از پدیدارشناسیِ طیفِ هیجانیِ خشم را---نه به‌صورتِ صورت‌بندیِ صریحِ نظری، بلکه به‌صورتِ معماریِ پنهانِ معناشناسیِ متنی---در خود نهفته دارد. بازخوانی این نظریه با ابزارهای ترکیبیِ زبان‌شناسیِ شناختی، معناشناسیِ طیفی و تحلیلِ پیکره‌ای، نه تنها افقی تازه برای مطالعاتِ قرآنی-زبانی می‌گشاید، بلکه می‌تواند زمینه‌ای استوار برای پژوهش‌های آتی در حوزه‌های روان‌شناسیِ اسلامی، ترجمه قرآن، علومِ تربیتی و فلسفه ذهن فراهم آورد.

ادغامِ سنتِ کلاسیکِ زبان‌شناسیِ عربی، چارچوب‌های نوینِ زبان‌شناسیِ شناختی، و ابزارهای محاسباتیِ معاصر، در یک پروژه واحد، نشان می‌دهد که سنت و نوآوری در پژوهش‌های قرآنی نه ناسازگار، بلکه مکمل‌اند. امیدواریم این پژوهش، در کنار آثار ارزشمندِ پیشینی نظیر آثار ایزوتسو، قائمی‌نیا، پاکتچی، نَریمانی و قاری‌زاده، گامی به سوی این ادغام مولّد باشد.

\begin{center}\rule{0.5\linewidth}{0.5pt}\end{center}

\subsection{فهرست منابع}\label{ux641ux647ux631ux633ux62a-ux645ux646ux627ux628ux639}

\subsubsection{الف) منابع عربی و قرآنی}\label{ux627ux644ux641-ux645ux646ux627ux628ux639-ux639ux631ux628ux6cc-ux648-ux642ux631ux622ux646ux6cc}

\begin{itemize}
\tightlist
\item
  ابن فارس، احمد بن فارس. (۱۴۰۴ ق / ۱۹۸۴ م). \emph{معجم مقاييس اللغة}. تحقیق عبد السلام محمد هارون. بیروت: دار الفکر.
\item
  ابن منظور، محمد بن مکرم. (۱۴۱۴ ق / ۱۹۹۳ م). \emph{لسان العرب} (الطبعة الثالثة، ۱۵ مجلد). بیروت: دار صادر.
\item
  ابن عاشور، محمد الطاهر. (۱۹۸۴ م). \emph{التحرير والتنوير} (۳۰ مجلد). تونس: الدار التونسية للنشر.
\item
  الراغب الأصفهاني، حسین بن محمد. (۱۴۱۲ ق / ۱۹۹۲ م). \emph{المفردات في غريب القرآن}. تحقیق صفوان عدنان داوودی. دمشق و بیروت: دار القلم و الدار الشامية.
\item
  الزبیدی، محمد مرتضی. \emph{تاج العروس من جواهر القاموس} (۴۰ مجلد). الکویت: التراث العربی.
\item
  الزمخشري، محمود بن عمر. (۱۴۰۷ ق). \emph{الكشاف عن حقائق غوامض التنزيل} (۴ مجلد). بیروت: دار الکتاب العربی.
\item
  الطباطبائي، محمد حسین. (۱۳۹۰ ق / ۱۹۷۰ م). \emph{الميزان في تفسير القرآن} (۲۰ مجلد). بیروت: مؤسسة الأعلمي للمطبوعات.
\item
  الطبرسي، الفضل بن الحسن. \emph{مجمع البيان في تفسير القرآن} (۱۰ مجلد). بیروت: دار المعرفة.
\item
  الفخر الرازي، محمد بن عمر. (۱۴۲۰ ق / ۱۹۹۹ م). \emph{مفاتيح الغيب (التفسير الكبير)} (۳۲ مجلد). بیروت: دار إحياء التراث العربي.
\item
  الفیروزآبادي، محمد بن یعقوب. \emph{القاموس المحيط}. بیروت: مؤسسة الرسالة.
\item
  المصطفوي، حسن. (۱۴۱۶ ق / ۱۹۹۵ م). \emph{التحقيق في كلمات القرآن الكريم} (۱۴ مجلد). تهران: وزارة الثقافة و الإرشاد الإسلامي.
\end{itemize}

\subsubsection{ب) منابع فارسی}\label{ux628-ux645ux646ux627ux628ux639-ux641ux627ux631ux633ux6cc}

\begin{itemize}
\tightlist
\item
  ایزوتسو، توشیهیکو. (۱۳۶۱). \emph{خدا و انسان در قرآن: معناشناسی جهان‌بینی قرآنی} (ترجمه احمد آرام). تهران: شرکت سهامی انتشار. (نسخه اصلی منتشرشده در ۱۹۶۴).
\item
  ایزوتسو، توشیهیکو. (۱۳۶۰). \emph{مفاهیم اخلاقی-دینی در قرآن مجید} (ترجمه فریدون بدره‌ای). تهران: نشر و پژوهش فرزان روز. (نسخه اصلی منتشرشده در ۱۹۶۶).
\item
  پاکتچی، احمد. (مجموعه آثار). \emph{روش‌شناسی معناشناسی واژگان قرآن کریم}. تهران: دانشگاه امام صادق (ع).
\item
  قائمی‌نیا، علیرضا. (۱۳۹۰). \emph{معناشناسی شناختی قرآن}. تهران: پژوهشگاه فرهنگ و اندیشه اسلامی.
\item
  قائمی‌نیا، علیرضا. (۱۳۹۳). \emph{استعاره‌های مفهومی و فضاهای قرآنی}. تهران: پژوهشگاه فرهنگ و اندیشه اسلامی.
\item
  قاری‌زاده، محمدعمر؛ سلمانی مارَوست، محمدعلی؛ میمندی، وصال؛ فارسی، بهنام. (۱۴۰۲). «چندمعنایی واژۀ «طغیان» در قرآن کریم با رویکردی بر زبان‌شناختی». \emph{پژوهش‌های زبان‌شناختی قرآن}، دانشگاه اصفهان.
\item
  نَریمانی، زهرا؛ اقبالی، مسعود؛ چاری، مجید. (۱۴۰۰). «بازشناسی معنایی واژگان «غیظ»، «غضب» و «سَخَط» در قرآن کریم با رویکرد تفسیری». \emph{پژوهش‌های قرآن و حدیث}، دانشگاه تهران.
\end{itemize}

\subsubsection{ج) منابع انگلیسی}\label{ux62c-ux645ux646ux627ux628ux639-ux627ux646ux6afux644ux6ccux633ux6cc}

\begin{itemize}
\tightlist
\item
  Averill, J. R. (1982). \emph{Anger and Aggression: An Essay on Emotion}. New York: Springer-Verlag.
\item
  Berkowitz, L. (1993). \emph{Aggression: Its Causes, Consequences, and Control}. New York: McGraw-Hill.
\item
  Dukes, K. (2011). \emph{Quranic Arabic Corpus \lr{(Morphology, Version 0.4)}}. University of Leeds. https://corpus.quran.com.
\item
  Gross, J. J. (1998). The emerging field of emotion regulation: An integrative review. \emph{Review of General Psychology}, 2(3), 271--299.
\item
  Gross, J. J. \lr{(Ed.)}. (2014). \emph{Handbook of Emotion Regulation} \lr{(2nd ed.)}. New York: Guilford Press.
\item
  Gross, J. J. (2015). Emotion regulation: Current status and future prospects. \emph{Psychological Inquiry}, 26(1), 1--26.
\item
  Izutsu, T. (1959). \emph{The Structure of the Ethical Terms in the Koran: A Study in Semantics}. Tokyo: Keio Institute of Philological Studies.
\item
  Izutsu, T. (1964). \emph{God and Man in the Koran: Semantics of the Koranic Weltanschauung}. Tokyo: Keio Institute of Cultural and Linguistic Studies.
\item
  Izutsu, T. (1966/2002). \emph{Ethico-Religious Concepts in the Qur'ān}. Montreal: McGill-Queen's University Press.
\item
  Kassinove, H., \& Tafrate, R. C. (2002). \emph{Anger Management: The Complete Treatment Guidebook for Practitioners}. Atascadero, CA: Impact Publishers.
\item
  Kennedy, C. (2007). Vagueness and grammar: The semantics of relative and absolute gradable adjectives. \emph{Linguistics and Philosophy}, 30(1), 1--45.
\item
  Kövecses, Z. (1986). \emph{Metaphors of Anger, Pride, and Love: A Lexical Approach to the Structure of Concepts}. Amsterdam: John Benjamins.
\item
  Kövecses, Z. (2000). \emph{Metaphor and Emotion: Language, Culture, and Body in Human Feeling}. Cambridge: Cambridge University Press.
\item
  Kövecses, Z. (2010). \emph{Metaphor: A Practical Introduction} \lr{(2nd ed.)}. Oxford: Oxford University Press.
\item
  Lakoff, G., \& Johnson, M. (1980). \emph{Metaphors We Live By}. Chicago: University of Chicago Press.
\item
  Lakoff, G., \& Kövecses, Z. (1987). The cognitive model of anger inherent in American English. In D. Holland \& N. Quinn \lr{(Eds.)}, \emph{Cultural Models in Language and Thought} \lr{(pp.~195--221)}. Cambridge: Cambridge University Press.
\item
  Lazarus, R. S. (1991). \emph{Emotion and Adaptation}. New York: Oxford University Press.
\item
  Maalej, Z. (2004). Figurative language in anger expressions in Tunisian Arabic: An extended view of embodiment. \emph{Metaphor and Symbol}, 19(1), 51--75.
\item
  Plutchik, R. (1980). \emph{Emotion: A Psychoevolutionary Synthesis}. New York: Harper \& Row.
\item
  Plutchik, R. (2001). The nature of emotions. \emph{American Scientist}, 89(4), 344--350.
\item
  Sassoon, G. W. (2010). The degree functions of negative adjectives. \emph{Natural Language Semantics}, 18(2), 141--181.
\item
  Scherer, K. R. (2001). Appraisal considered as a process of multilevel sequential checking. In K. R. Scherer, A. Schorr, \& T. Johnstone \lr{(Eds.)}, \emph{Appraisal Processes in Emotion} \lr{(pp.~92--120)}. New York: Oxford University Press.
\item
  Scherer, K. R. (2005). What are emotions? And how can they be measured? \emph{Social Science Information}, 44(4), 695--729.
\item
  Sharifian, F. (2011). \emph{Cultural Conceptualisations and Language: Theoretical Framework and Applications}. Amsterdam: John Benjamins.
\item
  Spielberger, C. D. (1999). \emph{Professional Manual for the State-Trait Anger Expression Inventory--2 \lr{(STAXI-2)}}. Odessa, FL: Psychological Assessment Resources.
\item
  \emph{Tanzil Project}. \lr{(n.d.)}. \emph{The Tanzil Quran Text} \lr{(Uthmani version)}. https://tanzil.net.
\item
  Wierzbicka, A. (1999). \emph{Emotions Across Languages and Cultures: Diversity and Universals}. Cambridge: Cambridge University Press.
\end{itemize}

\begin{center}\rule{0.5\linewidth}{0.5pt}\end{center}

\subsection{پیوست: داده‌ها، کد، و قابلیت بازتولید}\label{ux67eux6ccux648ux633ux62a-ux62fux627ux62fux647ux647ux627-ux6a9ux62f-ux648-ux642ux627ux628ux644ux6ccux62a-ux628ux627ux632ux62aux648ux644ux6ccux62f}

تمامِ کد، داده‌های تفصیلی و فایل‌های نمودار در انباره عمومیِ پژوهش به آدرسِ زیر در دسترس‌اند:

\begin{quote}
\textbf{\lr{\url{https://github.com/ali-kin4/quranic-emotion-spectrum}}}
\end{quote}

این انباره شامل: (الف) خط لوله‌ی استخراج به زبان پایتون؛ (ب) پیکره‌ی ریخت‌شناختیِ قرآنی \lr{(Dukes 2011)}؛ (ج) فایل‌های CSV کاملِ کنکوردانس، توزیعِ سوره‌ای، تحلیل آماری، و تحلیل شبکه؛ (د) فایل‌های PDF و PNG هفت نمودارِ منتشرشده در این مقاله؛ (ه) متنِ کاملِ نسخه‌ی فارسی و نسخه‌ی انگلیسی.

\paragraph{قابلیتِ بازتولید گام‌به‌گام.} پایپ‌لاین به‌صورت قطعی طراحی شده و ورودیِ یکسان، خروجی بایت-به-بایت یکسان تولید می‌کند. برای بازتولیدِ کامل، خواننده باید موارد زیر را اطمینان حاصل کند:

\begin{itemize}
\tightlist
\item \textbf{محیطِ پایتون}: \lr{Python 3.10} یا بالاتر.
\item \textbf{بسته‌های مورد نیاز}: \lr{pandas >= 2.0}, \lr{numpy >= 1.24}, \lr{networkx >= 3.0}, \lr{matplotlib >= 3.7}, \lr{arabic-reshaper}, \lr{python-bidi} (آخرین فهرستِ قطعی در فایلِ \lr{analysis/requirements.txt}).
\item \textbf{داده‌ی ورودی}: فایلِ \lr{quranic-corpus-morphology-0.4.txt} از پروژه‌ی \emph{Quranic Arabic Corpus} \lr{(corpus.quran.com)}، که در پوشه‌ی \lr{data/quran/} انباره قرار داده شود.
\end{itemize}

سپس، اجرای ترتیبیِ فرمان‌های زیر، تمامِ یافته‌های گزارش‌شده را بازتولید می‌کند:

\begin{verbatim}
pip install -r analysis/requirements.txt
python analysis/extract_concordance.py
python analysis/advanced_metrics.py
python analysis/visualize_spectrum.py
python analysis/visualize_advanced.py
python analysis/visualize_spectrum_fa.py
python analysis/visualize_advanced_fa.py
\end{verbatim}

خروجی‌های مرحله‌ی نخست در پوشه‌ی \lr{data/concordance/}، خروجی‌های آماری و شبکه‌ای در همان پوشه، و خروجی‌های گرافیکی در پوشه‌های \lr{figures/} (انگلیسی-برچسب) و \lr{figures\_fa/} (فارسی-برچسب) ذخیره می‌شوند.

\paragraph{مجوزها.} وضعیتِ مجوزها به تفکیکِ هر مؤلفه:

\begin{itemize}
\tightlist
\item \textbf{پیکره‌ی منبع}: \emph{Quranic Arabic Corpus, version 0.4} \lr{(Dukes, 2011)} --- منتشر شده تحت \lr{GNU General Public License (GPL)}.
\item \textbf{متنِ قرآن (نسخه‌ی عثمانی)}: \emph{Tanzil Project} --- تحت \lr{Creative Commons Attribution-NoDerivs 3.0 (CC-BY-ND 3.0)}.
\item \textbf{کدِ تحلیلیِ نگارندگان}: کلیه‌ی اسکریپت‌ها در \lr{analysis/}، تحت \lr{MIT License} منتشر می‌شود.
\item \textbf{فایل‌های CSV خروجی (کنکوردانس‌ها، آمارها، شبکه)}: محصولِ مشتق از پیکره‌ی \lr{GPL}-محور هستند و، در حالتِ احتیاط، تحت همان \lr{GPL} توزیع می‌شوند.
\item \textbf{فایل‌های نموداری (PDF/PNG)}: تحت \lr{Creative Commons Attribution 4.0 (CC-BY 4.0)} منتشر می‌شوند تا استفاده در کلاس‌های درسی و بازنشرِ علمی، با درجِ منبع، آزاد باشد.
\item \textbf{متنِ مقاله}: حقوقِ مادی و معنوی نزدِ نگارندگان است؛ ناشر، در صورتِ پذیرش، حقوقِ نشر را بر اساس قراردادِ خود تنظیم خواهد کرد.
\end{itemize}

\paragraph{تضادِ منافع و تأمین مالی.} نگارندگان اعلام می‌کنند که هیچ‌گونه تضادِ منافعی در رابطه با پژوهشِ حاضر وجود ندارد. این پژوهش بدون دریافت کمکِ مالیِ مشخص از نهادهای دولتی، بازرگانی، یا غیرانتفاعی صورت گرفته است.


\end{document}
